% =====================================================================
%  Harald Høffding, SPINOZA'S ETHICA.  Analyse og Karakteristik.  D. Kgl.
%  Danske Vidensk. Selsk. Skrifter, 7. Række, historisk og filosofisk Afd.
%  III, 3.  København: Andr. Fred. Høst & Søn; Bianco Lunos Bogtrykkeri,
%  1918.  97 pp.
%
%  TRANSCRIPTION -- GENERATED; DO NOT EDIT BY HAND.  Made in three steps, all rerunnable:
%    python3 ebook-method/spinozas-ethica/draft.py --emit --force
%      this template + the scan's text layer (Google's OCR, in Google's scan of the
%      Princeton copy of the volume; ~/bibliotek/Høffding, Harald/
%      1918-vidensk-selsk-skrifter-7r-hist-filos-III.pdf, see SCANS.tsv), then
%      patches.PATCHES (corrections read at the image that a word decision cannot express);
%    python3 ebook-method/collate/check.py spinozas-ethica --italics --apply
%      the decisions in ebook-method/spinozas-ethica/decisions.tsv, taken against an
%      independent tesseract reading of the images (rules.py writes the ones settled by
%      rule, each marked with its rule; the rest were read at the image), and the
%      italics decisions (--italics must be in the same call as --apply);
%    python3 ebook-method/spinozas-ethica/draft.py --post
%      patches.POST: corrections written against the collated text.
%  Run them chained (&&): each step exits without writing if a patch fails to match.
%  STATE (2026-10-04): collated, OPEN 0 (words and italics; ~1470 decisions, ~400 of
%  them at the image); all notes read at the image; headings, labels, small capitals
%  and italic titles set by rule in draft.py (layout) from the images; pp. 38 and 80
%  read in full at the image (2 small-capital faults on p. 38, since fixed by rule;
%  p. 80 clean); 35 misprints marked.  See ACCURACY-LEDGER.tsv.
%  No e-book of this book is known.
%
%  CONVENTIONS.  "% --- p. N ---" + \opage{N} at the first word of printed page N,
%  in the book's own pagination (the print also carries the volume's, N + 336, not
%  transcribed); a word broken over the turn is split with %.  Footnotes as printed,
%  1), 2), ... numbered per page (\footnote[k]).  Names in small capitals \textsc;
%  italics \textit; letterspacing \emph.  A note running over a page is set whole at
%  its mark.  Running heads, folios, the series line and sheet numbers at a sheet's
%  foot, and Google's watermark are not transcribed.  Print governs; misprints are
%  kept and marked PRINTED AS IS.
% =====================================================================
\documentclass[12pt]{article}
\usepackage[utf8]{inputenc}
\usepackage[T1]{fontenc}
\usepackage{libertinus}
\usepackage{libertinust1math}
\usepackage[danish]{babel}
\usepackage{textalpha}            % Greek: the subdivisions α, β, ... and quotations
\usepackage[protrusion=true,expansion=true]{microtype}
\usepackage{setspace}
\usepackage[margin=1.3in]{geometry}
\usepackage{fancyhdr}
\usepackage{hyperref}
\hypersetup{pdftitle={Spinoza's Ethica. Analyse og Karakteristik}, pdfauthor={Harald Høffding}, hidelinks}
\pagestyle{fancy}
\fancyhf{}
\fancyhead[LE,RO]{\thepage}
\fancyhead[C]{\textit{Harald Høffding: Spinoza's Ethica (1918)}}
\setstretch{1.3}
\usepackage{marginnote}
\usepackage{graphicx}
\DeclareUnicodeCharacter{0254}{\reflectbox{c}}  % ɔ: the old Danish id-est mark (a reversed c)
\renewcommand{\thefootnote}{\arabic{footnote})}
\newcommand{\tocline}[2]{\noindent #1\dotfill\ #2\par}
\newcommand{\opage}[1]{\marginnote{\footnotesize\textit{[#1]}}}

\begin{document}

% --- p. [1], title (PDF 389), read at the image 2026-10-04.  p. [2] is blank.
\begin{center}\vspace*{4em}
{\huge SPINOZA'S\ \ ETHICA}\\[2em]
{\large ANALYSE OG KARAKTERISTIK}\\[1em]
{\scriptsize AF}\\[1em]
{\large HARALD HØFFDING}\\[2em]
% (a short rule)
{\scriptsize D.\ \textsc{Kgl. Danske Vidensk. Selsk. Skrifter}, 7.\ \textsc{Række, historisk og filosofisk Afd.} III.\ 3}\\[8em]
% (an ornament)
{\large KØBENHAVN}\\[0.5em]
{\scriptsize HOVEDKOMMISSIONÆR: ANDR.\ FRED.\ HØST \&\ SØN, KGL.\ HOF-BOGHANDEL}\\
{\tiny BIANCO LUNOS BOGTRYKKERI}\\
{\tiny 1918}
\end{center}
\clearpage

% --- p. 3 ---
\setcounter{footnote}{0}%
\opage{3}\section*{Indledning.}


1. Et filosofisk System er et Forsøg paa at føre det i Erfaringen Givne, de
foreliggende Emner, tilbage til et Grundemne, et Urfænomen og dermed samle al
Erkendelse om én Grundtanke. Gyldigheden af denne Grundtanke kan søges paavist
dels gennem en Analyse af de Emner, der staar til Raadighed for Filosofen,
dels gennem et Forsøg paa at udlede disse Emner af Grundemnet. Ved den første
Fremgangsmaade bliver Grundemnet Konsekvens af de enkelte Emner; ved den
anden bliver disse Konsekvenser af hint. I intet System mangler der Forsøg paa at
anvende begge Fremgangsmaader. Men selv om saadanne Forsøg kunde lykkes,
vilde alligevel det Spørgsmaal staa tilbage, om den Kreds af Emner, Tænkeren
raader over, ikke er begrænset, og da nye Emner stedse kan opdukke, maa dette
Spørgsmaal besvares med ja. Intet System kan derfor blive andet end en paa Analyse
af en mere eller mindre omfattende Kreds af Emner bygget Hypotese. Hvad
den anden Fremgangsmaade angaar, vil der ogsaa rejse sig principielle Vanskeligheder,
idet det allerede rent logisk viser sig at være umuligt at udlede en Mangfoldighed
af Emner af et eneste Grundemne; til enhver Slutning behøves der nu
engang mere end én Præmis.

Denne Betragtning kan vel siges at være i den Grad trængt igennem i den
filosofiske Verden, at Systemernes Tid nu er forbi. I en vis Forstand vil der naturligvis
altid være „System“ i alvorlig Tænkning, idet hver enkelt Tanke maa
kunne „staa sammen“ med de andre Tanker, som fastholdes. Men „System“ i Betydning
af ikke blot afrundet, men afsluttet og definitiv Verdensforklaring eller
Verdensanskuelse hører Fortiden til.

Studiet af Fortidens filosofiske Systemer har dog mere end rent historisk Interesse.
Saadanne Tankebygninger var i og for sig naturlige, saalænge ingen
gennemført kritisk Prøvelse af den menneskelige Erkendelse var foretagen. Og det
viser sig ofte ved nærmere Betragtning, at Problemer, vi er tilbøjelige til at betragte
som helt moderne, ligger skjulte som Elementer i de uvilkaarlige Krystallisationer,
der foregik hos Fortidens Tænkere. Dersom nu de Emner, Tænkningen sysler med,
bevarer deres Betydning, --- dersom de Grundtanker, med hvilke der arbejdes, hænger
nøje sammen med den menneskelige Tænknings Natur, --- og dersom der arbejdes
med Alvor og med Sans for Problemernes Dybde, --- saa kan et saadant
% --- p. 4 ---
\setcounter{footnote}{0}%
\opage{4}System ikke blot staa som Udtryk for en svunden Tids Tænkning, men ogsaa tjene
til Forbillede for, hvorledes der maa tænkes under andre historiske Forudsætninger.


Om ikke mange Systemer gælder dette i den Grad som om \textsc{Spinoza}'s.


Et fuldstændigt, for alle Tider afsluttet System mente \textsc{Spinoza} nu ikke heller
at have opstillet, saaledes som en hyppig, men urigtig Opfattelse mener om ham.
I et Brev til \textsc{Oldenburg} (Ep. 30 i \textsc{van} \textsc{Vloten}'s Udgave) siger han, at han ikke véd,
hvorledes Tingene i Virkeligheden hænger sammen, og hvorledes enhver Del passer
ind i den Helhed, den hører til; --- dertil vilde jo udfordres, at man kendte hele
Naturen og alle dens Dele! I et Brev til \textsc{Boxel} (Ep. 56) siger han, at den største
Del af Naturens Grundegenskaber (Attributer) kender vi ikke; vi kender kun de to
Grundegenskaber, som kundgør sig i vort eget Væsen, Tænkning og Udstrækning
(Ep. 64 smlgn. Ethica II. Ax. 5). Og da en ung Ven af ham var gaaet over til
Katolicismen og nu i et Brev, i hvilket han forsøgte at omvende ham, haanende
spurgte, om han da mente at have fundet den bedste Filosofi af alle, svarede \textsc{Spinoza}
(Ep. 76), at han ingenlunde mente at være naaet til den bedste Filosofi, men
at han var vis paa at kende sand Filosofi. Ikke blot overfor Naturens Grundegenskaber
(Attributer), ogsaa med Hensyn til de enkelte, bestemte Tings Existens (existentia
modorum) staar vi paa den blotte Erfarings Grund (Ep. 10), og hvad Forstaaelsen
af disse enkelte i Erfaringen givne Ting angaar, staar vi overfor en uendelig
Opgave, idet den enkelte Ting (det enkelte Emne) forstaas ved sin Sammenhæng
med andre enkelte Emner, saa at Aarsagernes Række gaar i det Uendelige (Ethica
I, 28). Erfaringserkendelsen kan altsaa aldrig blive afsluttet.

Hvorvidt de her fremdragne Udtalelser kan bringes i strengt nødvendig
Sammenhæng med \textsc{Spinoza}'s hele øvrige Opfattelse, vil vise sig ved Gennemgang
af Hovedværkets enkelte Partier. Men før vi gaar over til en saadan Undersøgelse,
maa vi forsøge at finde Spor af det Tankearbejde, af hvilket Systemet er fremgaaet.
Naar det færdige System kan sammenlignes med en Krystal, bliver det kun forstaaeligt,
naar man kan efterspore Krystaldannelsen, det vil sige, finde de Forudsætninger,
der har været afgørende for Systemets Bygning.

Spinozaforskerne har, især i den nyeste Tid, lagt Hovedvægten paa at finde
historiske Forgængere for \textsc{Spinoza}'s ejendommelige Tankegang. Trods al den Lærdom,
der herved er udfoldet, kan man ikke sige, at man er kommen videre end
til blotte Muligheder. \textsc{Spinoza}'s første Biografer (Lucas, \textsc{Colerus}) henviste til \textsc{Descartes}.
Moderne Forskere har undersøgt hans Forhold til jødisk Middelalderteologi
(\textsc{Joel}), til Nyskolastiken (\textsc{Freudenthal}), til \textsc{Giordano} \textsc{Bruno} (\textsc{Sigwart}, \textsc{Avenarius}),
til de friere Bevægelser i Tidens hollandske Aandsliv (\textsc{Meinsma}), til
Fortidens Teologer i det hele (\textsc{Dunin}-\textsc{Borkowski})\footnote[1]{Det har været \textsc{Spinoza}'s Skæbne, at hans Personlighed og hans Lære gentagne Gange er blevne forsvarede af Mænd, som stod fjernt eller endog fjendtlig overfor hans hele Retning. Saaledes rensede Præsten \textsc{Colerus}, der selv var meget forarget over \textsc{Spinoza}'s Lære, hans Minde fra forskellige Bagvaskelser og er derfor bleven kaldt en moderne Bileam. En saadan Bileam var ogsaa \textsc{Jacobi}, hvis Følelsesteologi dog lod ham se med Rædsel paa \textsc{Spinoza}'s Filosofi. Og i vore Dage har Jesuiten \textsc{Dunin}-\textsc{Borkowski} ikke blot med stor Lærdom, men ogsaa med stor Sympati gransket \textsc{Spinoza}'s Udviklingsgang. --- Sympatien har sin Grund i, at han ser \textsc{Spinoza} som Modvægt mod Datidens og Nutidens Skepsis og dilettantiske Fritænkeri. Men i Sammenligning med, hvad \textsc{Dunin}-\textsc{Borkowski} kalder Verdensfilosofien, nemlig den middelalderlige Skolastik med \textsc{Thomas af Aqvino} som første Mand, staar \textsc{Spinoza}'s Filosofi dog som en stor Illusion. (\textsc{Stanislaus v. Dunin-Borkowski}: \textit{Der junge Spinoza}. Münster 1910, p. 351---485).}. Det er forlængst godtgjort, at
% --- p. 5 ---
\setcounter{footnote}{0}%
\opage{5}\textsc{Spinoza} ikke har spundet sin Visdom ud af sig selv, og at han i det hele ikke var
den Eneboer, hans første Disciple fandt det raadeligst at give ham Udseende af at
være. Han var en belæst Mand paa mange Omraader og har fra mange Sider kunnet
tage, hvad han kunde bruge. Men det Afgørende har været en Tankestrømning,
der uvilkaarlig og fra først af uden klar Bevidsthed udsprang hos kam, og som betingede % PRINTED AS IS: kam (ham)
Valget af, hvad han optog af Tidens Tanker, og tillige afgjorde, hvilken
Plads de fik i hans Tankebygning. Inderst inde har en aandelig Fornødenhed, der
tidlig opstod hos ham, været Tankestrømningens Kilde og tillige bestemmende for
dens Retning.

2. Det Tankearbejde, der tidlig kom i Gang hos \textsc{Spinoza}, begyndte indenfor
den jødiske Teologi og genoptoges paa Grundlag af den vesterlandske Naturvidenskab
og Filosofi, som han, især efter sin Udstødelse af Synagogen, havde levet sig
ind i. Paa begge Omraader er han gaaet Analysens og Konsekvensens Vej og
har søgt at fastholde og gennemføre Forudsætninger, som Forgængerne ikke havde
været sig klart bevidst. Og under hele dette Arbejde har han været ledet af en
Trang til Forstaaelse af sig selv og af sin Sammenhæng med den store, lovmæssige
og uendelige Tilværelse, en Trang, der hos ham paa én Gang havde en intellektuel
og en religiøs Karakter. Disse tre Elementer eller Tendenser vil jeg i Korthed karakterisere,
før jeg gaar over til gennem Analyse af Hovedværket (Ethica) at paavise,
hvorledes de brydes i hans Tænkning. I de Definitioner og Axiomer, hvormed
Ethica's første Bog begynder, har der krystalliseret sig Tanker, der vil kunne føres
tilbage til de tre antydede Kilder. Systemets Ejendommelighed og dets Forfatters
Personlighed lægger sig for Dagen i den Maade, paa hvilken de tre Motiver virker
sammen og stræber at komme til deres Ret.

\textit{a.} \textsc{Spinoza}'s Ungdomsskrift, den „korte Traktat“, har et mere teologisk og
mystisk Præg end Hovedværket Ethica. \textsc{Spinoza} gaar der ud fra Gudsbegrebet,
saaledes som dette havde udviklet sig i de højere Folkeligioner, og han drager alle
Konsekvenserne af, at Gud er et ubetinget og uendeligt Væsen, --- et Væsen uden
Forudsætninger og uden Grænser. Det er ad denne Vej, han kommer til det strenge
Begreb om Substans, der fremsættes i de første Definitioner i Ethica, --- Begrebet
om det, der bestaar ved sig selv og forstaas ved sig selv, det vil sige, hvis Begreb
ikke forudsætter noget andet Begreb, af hvilket det skulde dannes. Der viser sig
her en Trang til at naa et Sidste, et Forudsætningsløst, der ikke igen viser Tanken
ud over det selv. \textsc{Spinoza} imødegaar den Indvending, at enhver Definition jo dog
netop bestaar i, at et Begreb henføres under et højere Begreb (K. Tr. I, 7, 9---10).
Det gælder, hævder han, ikke de Grundbegreber, der betinger al Erkendelse;
% --- p. 6 ---
\setcounter{footnote}{0}%
\opage{6}hvis ogsaa de forudsatte andre Begreber, vilde intet nogensinde kunne blive fuldt
erkendt! Der fremtræder her en ejendommelig Forening af Mystik og Erkendelsesteori.
Trangen til at finde afsluttende Begreber er ikke rent intellektuel; den er ét
med Trangen til at komme til Ro. Skal vi ikke kastes ud i endeløse Tankerækker,
maa vi holde os til Tanken om en eneste Substans, et eneste Grundvæsen (K. Tr. I,
2, 10). Vor Tanke finder hverken Ro ved at blive staaende ved Erkendelsen af det
Aandelige eller ved Erkendelsen af det Legemlige, men maa gaa tilbage til det, der
betinger disse to Omraaders Bestaaen og gør Forstaaelsen af dem mulig: Gud eller
Substansen (K. Tr. II, 22, 5).

Baade i den jødiske, i den muhammedanske og i den kristelige Teologi har
Gudsbegrebet teoretisk Betydning, afgiver Principet for Forstaaelse, i hvert Tilfælde
for den afsluttende Forstaaelse af alt. Og det er denne Side af Gudsbegrebet, hvorom
\textsc{Spinoza}'s Tænken mere og mere koncentrerer sig. Hans Forhold til Teologien
kan betegnes ved, at han gør Alvor af Spørgsmaalet om, hvilken Karakter den afsluttende
Tanke maatte have. Den kan ikke være uforstaaelig, og naar den ikke
skal føre ud over sig selv, maa den forstaas ved sig selv. For at klare os, hvad
\textsc{Spinoza}'s Substans maa betyde, maa vi derfor spørge, hvad det er, der forstaas ved
sig selv, og som kan gøre alt andet forstaaeligt. Det kan kun være den nødvendige
Sammenhæng mellem Subjekt og Prædikat eller mellem Forudsætninger og
Slutning, den Sammenhæng, der ligger til Grund ved enhver Tankeovergang. Længere
tilbage end til Principet for al Tænken og derved tillige for al den Rationalitet,
vi kan finde i Tilværelsen, kan vi ikke komme. Vi kan ikke give nogen Begrundelse
af, hvorfor Konklusionen følger af sine Præmisser. Det er ved en analytisk
Intuition\footnote[1]{Se om dette Begreb min Afhandling om Intuition i Vid. Selsk. Oversigter 1914, p. 77---80.}, at vi fremdrager Principet for al Rationalitet. Principet for al Forstaaelse
staar i \textsc{Spinoza}'s Metafysik tillige som Substansen, det Væsen, der ligger til Grund
for alt. Hans Metafysik er en Projektion af hans Erkendelsesteori. Han begynder
teologisk med at sige, at Gud er det, der gør alt forstaaeligt for os, og han ender
med at sige, at det, der gør alt forstaaeligt for os, er Gud. Medens han i den korte
Traktat gaar fra Gudsbegrebet til Substansbegrebet, gaar han i de første Definitioner
i Ethica's første Bog omvendt fra Substansbegrebet til Gudsbegrebet. Denne Omvending
forberedes i det første Tillæg til den korte Traktat, og den fremtræder afgjort
i de første Breve (se især Ep. 2).

Det Gudsbegreb, \textsc{Spinoza} tog med sig fra sin teologiske Periode og stedse fastholdt
i sit Hovedværk, var naturligvis meget forskelligt fra Folkereligionernes Gudsbegreb,
og han gør selv opmærksom paa Forskellen (se f. Ex. Ep. 19. 21. 73). Han
konstruerer sit Gudsbegreb ud fra Forudsætningerne for al rationel Forstaaelse af
Tilværelsen, men maa saa (fortsættende hvad jødiske Religionsfilosofer i Middelalderen,
især \textsc{Maimonides} og \textsc{Ibn} \textsc{Gebirol}, havde indledet) søge at udrense alle Antropomorfismer,
alle konkrete psykologiske Egenskaber fra dette Begrebs Indhold.

I de samme Aar, i hvilke \textsc{Spinoza}'s Brud med den Folkereligion, i hvilken han
var opvokset, forberedtes ved hans Tænknings Selvstændighed, vendte en anden stor
% --- p. 7 ---
\setcounter{footnote}{0}%
\opage{7}Aand, \textsc{Blaise} \textsc{Pascal}, gennem en religiøs Krise tilbage til sin Folkereligions Gudsbegreb
med udtrykkelig Hævden af alle Antropomorfismer. „Abrahams, Isaks og
Jakobs Gud, Jesu Christi Gud, ikke Filosofernes og Videnskabsmændenes Gud!“
udraaber han i den lidenskabelige Udtalelse, som han nedskrev den 25. Novbr. 1654
og senere stedse bar hos sig indsyet i sin Klædning.

\textsc{Spinoza}'s Gud var nu netop „Filosofernes og Videnskabsmændenes Gud“, forsaavidt
som hans Gudsbegreb er et metafysisk Udtryk for al Tænknings og Videnskabs
sidste Forudsætning. Da nu denne Forudsætning ikke staar i udvortes Forhold
til de enkelte Emner, men kundgør sig ved en Forstaaelse af dem og ved al
Fastsætten af Betingelserne for deres Bestaaen, saa maa \textsc{Spinoza}'s Gud være en
iboende (immanent) Gud, ikke en ydre (transscendent) Gud. Sin Modsætning til det
sædvanlige Gudsbegreb betegner \textsc{Spinoza} saaledes (Ep. 73): „Jeg har om Gud og
Naturen en helt anden Opfattelse end den, moderne Kristne hylder. Jeg hævder
nemlig, at Gud er iboende (immanent), ikke ydre (transeunt) (= transscendent) Aarsag
til alle Ting“. I den korte Traktat (første Dialog) havde han belyst dette Punkt
ved Sammenligning med Forstanden, der rører sig i alle enkelte Tanker, men ikke
er ydre Aarsag til dem\footnote[1]{\textsc{Hume}, der ikke kan have kendt den korte Traktat, og som vel i det hele kun kender \textsc{Spinoza} fra \textsc{Bayle}'s Dictionnaire, benytter \textsc{Spinoza}'s Lære til at forskrække sine spiritualistiske Modstandere. Naar disse hævdede, at vore enkelte Fornemmelser og Forestillinger alle var Led i en og samme Bevidsthed, ikke, som \textsc{Hume} mente, stod i et rent udvortes, mekanisk Associationsforhold til hverandre, saa benyttede de jo, siger han, den samme Tankegang (Begrebet Iboen), som førte \textsc{Spinoza} til hans ugudelige Meninger! (\textit{Treatise on Human Nature}. I, 4, 5).}. --- \textsc{Spinoza}'s indre eller iboende Aarsag svarer til, hvad
\textsc{Bruno} kalder „Princip“, medens \textsc{Bruno} ved „Aarsag“ mener, hvad \textsc{Spinoza} kalder
ydre Aarsag\footnote[2]{\textsc{Bruno}: \textit{De la causa, principio et uno}. Opere ed. Lagarde, p. 230, 272.}. ---

Deraf, at Gud er Princip for al Forstaaelse, følger ikke blot, at han maa være
„iboende“, ikke „ydre Aarsag“, idet han forholder sig til alle Emner som Grunden
forholder sig til Følgen. Der følger deraf tillige, at der ingen Grænse kan sættes for
Guds Væsen. Alle Tilværelsens Grundegenskaber maa være Guds Egenskaber. Det
er ved denne Tankegang, at Gudsbegrebet for \textsc{Spinoza} ikke blot er ét med Substansbegrebet,
som vi allerede har set, men ogsaa ét med Naturbegrebet, saa at
han kan bruge de tre Udtryk Gud, Substans og Natur som ensbetydende. I den
første Dialog, der synes at være den ældste Del af den korte Traktat, spørges der,
om der kan tænkes „et Væsen, som er højst fuldkomment, idet det ikke kan begrænses
ved noget andet Væsen“, og Svaret lyder paa, at Naturen er et saadant
Væsen, idet alt er indbefattet i den. Og det paavises desuden, at denne Uendelighed
hænger nøje sammen med Begrebet iboende Aarsag. Hvis Emnerne (Tingene i
Verden) skyldtes en ydre, fra dem forskellig Aarsag, vilde denne Aarsag være begrænset.
--- Der er saaledes en indre Sammenhæng mellem Forstaaelsesprincip,
Iboen og Uendelighed, og det er \textsc{Spinoza}'s konsekvente Gennemtænken af det overleverede
Gudsbegreb, som fører ham til alle tre Begreber.

% --- p. 8 ---
\setcounter{footnote}{0}%
\opage{8}\textsc{Avenarius}\footnote[1]{\textsc{Rich}. \textsc{Avenarius}: \textit{Die beiden ersten Phasen des Spinozischen Pantheismus}. (1868).} og \textsc{Sigwart}\footnote[2]{\textsc{Chr. Sigwart}: \textit{Spinoza's kurzer Tractat\textsuperscript{2}}. Prolegomena. (1869).} antager, at Naturbegrebet er blevet Led i \textsc{Spinoza}'s
Tænkning gennem Indfiydelse fra Renaissancens Tænkere, maaske fra \textsc{Bruno}. Dette % PRINTED AS IS: Indfiydelse
er højst sandsynligt, og det vil ved nærmere Betragtning føre os over til det andet
store Tankemotiv hos \textsc{Spinoza}. Men som Naturbegrebet dukker op i den første
Dialog, i Svaret paa det Spørgsmaal, om der gives et Væsen, der ikke er begrænset
af noget andet, forudsættes det, at der allerede forud har rørt sig en Trang til at
afslutte Tænkningen i et Begreb, der igen ikke fører ud over sig selv. Det er \textsc{Spinoza}'s
fundamentale Trang, og den har udviklet sig under hans Gennemtænken af
det overleverede Gudsbegreb. Det er fra Tanken om Gud som Principet for al Forstaaelse,
han er kommen til Tanken om Guds Iboen og om Guds eller Naturens
Uendelighed.

Den Tankegang, der fører \textsc{Spinoza} til at hævde Uendelighedsbegrebets Gyldighed,
baade hvad Substansen (som absolut uendelig) og hvad dens Attributer (som
hver i sin Art uendelige) angaar, kan minde om den Tankegang, der førte \textsc{Bruno}
til at forkaste Gyldigheden af det gamle, begrænsede Verdensbillede. De to første
Sætninger i Ethicas anden Bog antyder den Vej, \textsc{Spinoza} er gaaet, og som har ført
ham til anden og sjette Definition i første Bog. Vi kan, hedder det (Eth. II, 1. Schol.),
stedse føje Tanke til Tanke, uden at der her er nogen Grænse, og ligesaa kan vi
(Eth. II. 2) føje Udstrækning til Udstrækning uden Grænse. Tankens og Udstrækningens
Verden er hver for sig uendelige indenfor deres Art (in suo genere), idet
der indenfor dem stedse kan skrides videre frem efter samme Lov. Og paa samme
Maade forholder det sig med det Begreb, der ligger til Grund for alle Begreber,
nemlig Begrebet om noget, der kan fremtræde i uendelig mange Former eller Arter,
idet disse forskellige Arters (genera) eller Attributers Række stedse kan fortsættes,
ligesom Rækken af de forskellige Emner indenfor samme Art. Derved fremkommer
Definitionen af Gud i Begyndelsen af Ethica (Eth. I, Def. 6): „Ved Gud forstaar jeg
et absolut uendeligt Væsen, d. v. s. en Substans, bestaaende af uendelig mange
Attributer, af hvilke ethvert udtrykker et Evigt og Uendeligt“. --- I Ep. 36 antydes
hele denne Tankegang klart.

Ogsaa hos \textsc{Bbuno} anvendes, hvad vi kan kalde den kontinuerlige Rækkedannelses % PRINTED AS IS: BBUNO
Lov. Baade Sansning, Fantasi og Tanke viser os, siger han, at vi stedse kan
føje Led til Led. Den sanselige Horizont kan stedse udvides, og vi kan føje Billede
til Billede, Størrelse til Størrelse, Tal til Tal, Tanke til Tanke i det uendelige. Derfor
spørger han, om vi kan tænke os Verden begrænset: paa selve den Grænse, vi
forsøger at sætte for Verden, vilde der jo, naar vi tænker os hensatte der, danne sig
en ny Horizont omkring os\footnote[3]{Se herom \textit{Den nyere Filosofis Historie\textsuperscript{2}}, I, p. 118---123. --- Især i Bogen \textit{De Immenso} (Francofurti 1591) kommer denne Tanke frem. Dette Skrift fremstiller tillige Forholdet mellem Frihed og Nødvendighed paa samme Maade som \textsc{Spinoza}. (Eth. I, Def. 7).}. Uendelighedsbegrebets logiske Genesis viser sig at
være den samme hos de to Tænkere, og \textsc{Bruno} er her Forgængeren. Forskellen
% --- p. 9 ---
\setcounter{footnote}{0}%
\opage{9}mellem deres Anvendelse af dette Begreb hænger sammen med Forskellen mellem
deres Problemer. Hos \textsc{Bruno} er det Udvidelsen af Verdensbilledet, der er Hovedsagen,
--- hos \textsc{Spinoza} er det Gennemførelsen af Lovmæssigheden indenfor Verdensbilledet.
Men der foreligger her en Mulighed for Indflydelse af \textsc{Bruno} paa \textsc{Spinoza},
som angaar et væsentligere Punkt end dem, man i Regelen har holdt sig til ved Drøftelsen
af Forholdet mellem de to Filosofer.

Hos dem begge staar naturligvis det Spørgsmaal tilbage, med hvilken Ret vi
tillægger den psykologiske eller logiske Mulighed\footnote[1]{Det er ikke tydeligt, hvilket Princip for Rækkedannelse \textsc{Spinoza} tænker paa, naar han siger, at vi blot ved at holde os til Tænkningens Attribut kommer til Tanken om et uendeligt Væsen (Eth. II, 1. Schol.). Det kunde være en rent psykologisk Række af Sansninger (Eth. II, 18. Schol.), eller (sandsynligere) en Række af Grunde og Følger (Eth. I, 28), eller (usandsynligt) en Række af Led, i hvilket hvert forholder sig til det foregaaende som Subjekt til Objekt (Eth. II, 21. Schol.). --- Smlg. om forskellige Rækkedannelser \textit{Den menneskelige Tanke}, p. 164---172, og specielt om Rækkedannelser i Kraft af Gentagelsens Princip \textit{Totalitet som Kategori} (Vid. Selsk. Skr. 1917), p. 34---37.} af fortsat Rækkedannelse real
Gyldighed. Dette Spørgsmaal kommer først frem i den kritiske Filosofi. Sammenligner
vi i denne Henseende \textsc{Bruno} og \textsc{Spinoza}, staar \textsc{Spinoza} mere dogmatisk end
hans Forgænger. Medens for \textsc{Spinoza} det, at vi kan konstruere et Uendelighedsbegreb,
er Bevis nok for, at der gives et uendeligt Væsen, saa skelner \textsc{Bruno} mellem
Tankekonstruktion og Virkelighed, og specielt mellem Matematik og Fysik\footnote[2]{F. Ex. \textit{La cena de le ceneri} (Opere ital. ed. Lagarde, p. 184): Altro è giocare con la geometria, altro è verificare con la natura.}.

\textit{b.} Efter sin Udtrædelse af Synagogen gjorde \textsc{Spinoza} Bekendtskab med Vestens
Videnskab, Filosofi og Literatur. Og naar han nu, efter at have opstillet Gudsbegrebet
som Princip for al Forstaaelse, igen spørger, hvad Forstaaelse da egentlig
vil sige, hævder han, at vi faar Besked derom ved at holde os til de Grundtanker
(notiones communes), som viser sig at finde Anvendelse overalt i Naturen, i det Enkelte
som i det Hele, altsaa til de universelle Love for alt Existerende (Tract. theol.
polit. Kap. 6---7). Kun ad denne Vej kan vi lære Guddommen at kende. „Hvad
mig angaar,“ siger \textsc{Spinoza} i et Brev (Ep. 21), „saa har jeg ikke lært og ikke kunnet
lære nogen af Guds evige Egenskaber at kende af den hellige Skrift.“ --- Teologien
havde lært ham, at Gud er Forstaaelighedens Princip, og han henvender sig
nu til Videnskabens Forudsætninger for at faa at vide, hvad Forstaaelse er.

Man har, tror jeg, ikke været tilstrækkelig opmærksom paa, hvad en Sammenligniug % PRINTED AS IS: Sammenligniug
mellem den korte Traktat og Ethica tydelig viser, at naturvidenskabelige
Grundsætninger danner Baggrund for \textsc{Spinoza}'s Filosofi i dens definitive Form.
Naar han (Eth. I, 28) kræver, at ethvert Fænomen (Emne) skal forklares ved et
andet, dette ved et tredie, o. s. v., saa indeholdes heri den af \textsc{Kepler} opstillede
Fordring (som \textsc{Kepler}\footnote[3]{\textit{Den nyere Filosofis Historie\textsuperscript{2}}, I, p. 165. --- Der fandtes et Skrift af \textsc{Kepler} i \textsc{Spinoza}'s Bibliotek, men det var et kronologisk Værk, ikke et af hans rent naturvidenskabelige. Af matematiske og naturvidenskabelige Værker har \textsc{Spinoza} besiddet Bøger af \textsc{Vieta}, \textsc{Huyghens}, \textsc{Riolan}, \textsc{Steno} og \textsc{Boyle}, foruden af \textsc{Descartes}.} selv anvendte til Kritik af sin Ungdoms Spekulationer), at
de Aarsager, som kan anerkendes i Videnskaben, selv maa kunne paavises i Er%
% --- p. 10 ---
\setcounter{footnote}{0}%
\opage{10}faringen. Dette Princip falder for ham sammen med den af \textsc{Galilei} og \textsc{Descartes}
opstillede Inertisætning, som han Eth. II, Lemma 3 opstiller som gyldig for materielle
Emner og ogsaa (Eth. III, 6. kfr. II, 1. Schol.) anvender paa aandelige Emner. Dertil
slutter sig Sætningen om det konstante Forhold mellem Hvile og Bevægelse, et
andet Udtryk for \textsc{Descartes}'s Sætning om Bevægelsens Bestaaen og ligesom denne
en ufuldkommen Forløber for den moderne Sætning om Energiens Bestaaen (Eth.
II, Lemma 7, Schol. kfr. Ep. 32. 64). --- I den „geometriske“ Fremstilling, \textsc{Spinoza}
uheldigvis har givet af sin Lære, udledes de nævnte Sætninger af Begrebet
Attribut, der (Eth. I, def. 4) defineres som en Egenskab ved Substansen, der forstaas
ved sig selv og derfor ikke forstaas ved Hjælp af nogen anden Egenskab. Deraf
følger bl. A., at et udstrakt Fænomen maa forklares ved et andet udstrakt Fænomen.
Men i Virkeligheden er \textsc{Spinoza} kommen til sit strenge Begreb om Attribut
ved Analyse af naturvidenskabelige Grundbegreber. Hvad der ved syntetisk Fremstilling
stilles først, er i Analysen kommet tilsidst. At Udstrækning ikke kan begrænses
af Tanke, eller Tanke af Udstrækning, at man altsaa i sin Naturforklaring
ikke maa springe over fra et Attribut til et andet (Eth. I, def. 2; I, 10 og II, 6); det
er tilsidst et metodisk Princip, der er fundet ved Analyse af de nævnte naturvidenskabelige
Sætninger, skønt \textsc{Spinoza} fremsætter det som en evig Sandhed.

En af \textsc{Spinoza}'s Korrespondenter indvendte mod \textsc{Spinoza}'s Definition af Attribut,
at det dog ikke var utvivlsomt, at Tanke ikke kunde begrænses ved Legeme
eller Legeme ved Tanke, da Tanken maaske kunde være en legemlig Bevægelse
(motus corporeus) (Ep. 3). \textsc{Spinoza} svarer (Ep. 4), at Udstrækning, som Udstrækning
betragtet, ikke er Tanke, og at det er Alt, hvad han behøver for at fastholde
sin Sætning, at Tanke og Udstrækning ikke kan begrænse hinanden. --- Rækkedannelsens
Princip, som ovenfor er omtalt, ligger tydelig til Grund herved. \textsc{Spinoza}
fordrer Rækker af ensartede Led. Og derved kommer vi over til den mest fundamentale
Sætning i \textsc{Spinoza}'s System, Aarsagssætningen, og til den Maade, paa hvilken
\textsc{Spinoza} opfatter den.

Aarsagsforholdet er os ifølge \textsc{Spinoza} kun forstaaeligt, naar dets Led, det, vi
kalder „Aarsag“, og det, vi kalder „Virkning“, er af samme Natur, saa at der bestaar
et Identitetsforhold mellem dem. I Afhandlingen om „Forstandens Forbedring“
siger han, at Erkendelsen af Virkningen ikke er Andet end en fuldkommen Forstaaelse
af Aarsagen og forudsætter denne. „Hvis to Ting ikke har nogen fælles
Egenskab, kan de ikke forstaas ved hinanden, eller den enes Begreb indeholder
ikke den andens“ --- hedder det i fjerde og femte Axiom i Ethica's første Bog. Og
i Ep. 4 udtales, at hvis to saadanne Ting skulde staa i Aarsagsforhold til hinanden,
vilde der være Noget i Virkningen, som ikke er i Aarsagen, og dette Noget vilde da
have sin Grund i Intet.

\textsc{Spinoza} forlanger altsaa Ækvivalens mellem Aarsag og Virkning, ligesom der
i Slutninger, som kan vendes om, er Ækvivalens mellem Præmisser og Konklusion.
Han vil indenfor Fænomenernes (Emnernes) Verden finde en ligesaa streng og logisk
Sammenhæng, som mellem Led i en Række af logiske Slutninger, der kan
% --- p. 11 ---
\setcounter{footnote}{0}%
\opage{11}vendes om. Man skal ikke begynde med Abstraktioner, men med Tanker, der
gør Emnernes, baade de aandeliges og de materielles, Opstaaen og Rækkefølge forstaaelige.
Jo mere vi forstaar de enkelte Emner, des mere forstaar vi Gud (Eth. V,
24), ti det, der gør os dem forstaaelige, er netop Gud. Hvad der bestaar ved sig
selv og forstaas ved sig selv, er ikke selv en enkelt Ting, et enkelt Emne. Det er
det Begrundelsens eller Lovmæssighedens Princip, som er al Videnskabs sidste Forudsætning,
og som godtgør sig i det Hele som i alle enkelte Dele af Tilværelsen.
Gudserkendelse og Naturvidenskab var for \textsc{Spinoza} saa langt fra at udelukke hinanden,
at de tværtimod faldt sammen. Det at tænke og forske var for ham en religiøs
Akt.

I et berømt Værk, der udkom en halv Snes Aar efter \textsc{Spinoza}'s Ethica, findes
tildels en lignende Tankegang. I sine „Principia“ (og senere i „Optics“) opfattede
\textsc{Newton} Rummet som Guds Sanseorgan (sensorium), Organet for Allestedsnærværelse.
Matematisk Forsken blev da ogsaa Gudserkendelse. Baade \textsc{Newton}'s og
\textsc{Spinoza}'s Tanke vilde, ved historisk at forfølges, føre tilbage til \textsc{Platon}. I andre
Henseender er \textsc{Spinoza}'s og \textsc{Newton}'s Religionsfilosofi højst forskellige. \textsc{Newton} vil
væsentlig begrunde Antagelsen af Guds Existens ad teleologisk Vej, ud fra Hensigtsmæssigheden
i Verdensbygningen. Men for \textsc{Spinoza} var det at tænke sig Gud
handle efter Formaal ét med at tænke sig ham som et endeligt og ufuldkomment
Væsen, der endnu ikke besidder alt, hvad der stemmer med dets Natur. De skarpeste
Udtalelser, der findes hos \textsc{Spinoza}, fremkommer netop i hans Kritik af Teleologien.


Naar der i det Foregaaende er lagt Vægt paa, at \textsc{Spinoza}'s Filosofi, særlig hans
afsluttende Opfattelse af Begreberne Substans og Attribut, skyldes en Analyse af
Erfaringsvidenskabens Aarsagsforklaring, har en skarpsindig Spinozaforsker derimod
indvendt, at \textsc{Spinoza}'s Substansbegreb ingen kausale Bestemmelser indeholdt, og at
man af hans Definition af Substans (Eth. I, def. 3) ikke heller kan udlede saadanne
Bestemmelser\footnote[1]{\textsc{Gabriel} \textsc{Huan}: \textit{Le dieu de Spinoza}. Paris 1914, p. 44 f.}. Denne i sig selv rigtige Bemærkning indeholder en Kritik af \textsc{Spinoza},
ikke en Kritik af min Opfattelse af ham. Naar \textsc{Spinoza} faktisk opfatter Kausalitet
som et rent logisk Forhold, behøvede han i sin Definition af Substansen
ikke at optage andre end rent logiske Bestemmelser. \textsc{Spinoza}'s Filosofi er et Forsøg
paa at tænke Verden ved de formale Kategorier, til hvilke han førte alle reale
Kategorier (som Kausalitet) og alle ideale Kategorier (Værdibegreberne) tilbage\footnote[2]{Om de forskellige Grupper af Kategorier se \textit{Den menneskelige Tanke}.}.
Man kan kritisere ham dels saaledes, at man viser, at han ikke i Virkeligheden
har gennemført Reduktionen af alle Kategorier til de formale, dels saaledes, at man
godtgør Umuligheden af en saadan Reduktion. Men som hans Værk faktisk foreligger,
vil man ikke kunne nægte, at det er karakteriseret ved et i sig selv storartet
Forsøg i den nævnte Retning. I den følgende Analyse af Værkets enkelte Dele vil
der blive Lejlighed til Belysning af dette Hovedpunkt.

% --- p. 12 ---
\setcounter{footnote}{0}%
\opage{12}\textit{c.} Efter \textsc{Spinoza}'s egen Erklæring i Begyndelsen af den ufuldendte Afhandling
„om Forstandens Forbedring“ har hans personlige Livserfaring været af afgørende
Betydning for ham under Udviklingen af hans Hovedtanker. Naar han skulde
gøre sig Rede for, hvilken Betydning filosofisk Tænkning havde for ham, eller rettere,
hvorledes filosofisk Tænkning for ham var bleven det, hvori han tilsidst fandt
Livets Værdi, saa laa for ham Forklaringen i, at det eneste, der staar i vor Magt
og gør os uafhængige af ydre Forhold, er selve Tanksarbejdet og den Erkendelse, % PRINTED AS IS: Tanksarbejdet
der gennem det kan vindes af vor Aands Sammenhæng med hele Naturen. Andre
Goder --- Nydelser, Ære, Rigdom --- beror paa noget, der ikke udspringer af vort
inderste Væsen. Men gennem Tankearbejdet naas der en Erkendelsesglæde og en
Kærlighed til Guddommen, som hæver over al Sorg og Uro. Her saa han et Maal
for sit Liv, et Maal, som tillige muliggjorde og medførte en Bestræbelse for, at andre
Mennesker kan naa det samme Maal. Ti, som han senere udtalte (Eth. IV,
App. cap. 9), den bedste Maade, paa hvilken et Menneske kan vise, hvilken Evne
og Aand han raader over, er den at paavirke Andre saaledes, at de kommer til at
leve efter sand Erkendelses eget Herredømme.

Det var ikke for at fly det virkelige Liv, at \textsc{Spinoza} valgte denne Vej. Han
var ikke Pessimist. Han vilde ikke have sympatiseret med den unge \textsc{Schopenhauer}'s
Ord til \textsc{Wieland}: „Livet er en mislig Sag; jeg vil tilbringe mit med at tænke over
det“. Tankens Liv var for \textsc{Spinoza} det højeste Liv, ikke fordi det stod i Modsætning
til Virkelighedens Liv, men netop fordi det paa sin Vis formaaede at give et
Udtryk for alle Livets Sider, ligesom det selv var en Side eller Form af Liv.

Hvorvidt \textsc{Spinoza}'s ædle og alvorlige Livsførelse var konsekvent ud fra hans
eget Standpunkt, var et Spørgsmaal, der blev meget drøftet paa og efter \textsc{Spinoza}'s
egen Tid, medens man vel nutildags --- med fornødent Kendskab til hans Tankegang
--- neppe vil se noget Problem deri. Der er et andet Spørgsmaal, som for os
har mere Interesse, nemlig, i hvilken Grad Eftervirkningen af det positivt religiøse
Standpunkt har havt vedvarende Indflydelse paa hans Sind.

Det har sikkert ikke været blot Akkommodation, naar han brugte Ordene
Gud, Natur og Substans som ensbetydende. I sin Filosofis Grundbegreb har han
paa én Gang set et religiøst og et rationelt Indhold udtrykt. Der har for ham ikke
gjort sig nogen Modsætning gældende her. Det har særlig ikke været blot Akkommodation,
naar han ofte, især i sine Breve, anvender kristelige Udtryk som Symboler
for Tanker, han selv gav et helt andet Indhold. Det havde sikkert været heldigst,
om han havde fulgt den Regel, han selv opstiller i et af sine Breve (Ep. 23),
at en Filosof helst skal undgaa at bruge teologiske Udtryk, men \textsc{Wundt} har afgjort
Ret i den Bemærkning, han gør i sin Karakteristik af \textsc{Spinoza}, at dennes Opfattelse
var helt forskellig fra Oplysningstidens: „Medens Oplysningstiden med velberaad Hu
tildækker den religiøse Tænknings mystiske Dybder, er det netop Gudsbegrebets
og den religiøse Følelses mystiske Indhold, som \textsc{Spinoza} søger at omsætte
i fornuftig Erkendelses Form“\footnote[1]{\textit{Ethik\textsuperscript{4}}, II, p. 113.}. Gennem selve den arbejdende Tanke, hvis Betyd%
% --- p. 13 ---
\setcounter{footnote}{0}%
\opage{13}ning for Systemets Tilblivelse jeg i denne Afhandling søger at betone, mente \textsc{Spinoza}
at høre Guddommens Stemme. I de fire første Bøger af Ethica træder Systemets
mystiske Karakter tilbage i Sammenligning med den korte Traktat, men i den
sidste Del af femte Bog gør den sig meget stærkt gældende. Det mystiske Element
kommer her til at skyde det rationelle Element tilside.

I nyeste Tid har \textsc{Dunin}-\textsc{Borkowski} hævdet, at \textsc{Spinoza} selv var hildet i en
Illusion med Hensyn til sin „moralske Energi“. „\textsc{Spinoza}“, siger han\footnote[1]{\textit{Der junge Spinoza}, p. 145.}, „troede, at
det Lys, der strømmede ind over hans Liv, kom fra hans selvskabte filosofiske
Firmament. Men dette Lys skyldtes faktisk andre Stjerner\dots{} [dem], der havde
bestraalet ham i hans Ungdom!“ Denne Bemærkning kan støttes ved en psykologisk
Kendsgerning af stor Betydning. Følelsesvirkningerne af en Oplevelse eller en
Tankegang kan holde sig meget længe efter, at det oprindelige Motiv er faldet bort;
de kan ogsaa gaa ind som Elementer i Følelser, som opstaar under Indflydelse af
nye Oplevelser eller nye Tanker. Jeg har i min Psykologi kaldt dette Fænomen
for Følelsens Expansion. Men den lærde og overfor \textsc{Spinoza} saa sympatiske Jesuit
lægger for liden Vægt paa Indflydelsen af \textsc{Spinoza}'s store og dybtgaaende Tankearbejde
i de sidste tyve Aar af hans Liv, et Tankearbejde, i hvilket hans hele Personlighed
virkede med, og som derfor førte til en Koncentration, en Totaldrift, i
Forhold til hvilken andre Tilskyndelser maatte blegne. Den Sætning, at en Følelse
kun kan kues eller fortrænges ved en anden Følelse af modsat Retning og større
Styrke (Eth. IV, 7), har \textsc{Spinoza} kunnet lære af egen Erfaring, alt som hans hele
aandelige Energi samledes om Erkendelsestrangen.

Naar en Søfarende kommer fra den sydlige Halvkugle op paa den nordlige,
er det nye Stjerner, han faar Øje paa og maa orientere sig efter. Erindringen om
Sydens Himmel falder derfor ikke bort, men faar ikke mere afgørende Betydning.

De tre Tankemotiver, der kan spores hos \textsc{Spinoza}, og hvis Brydning karakteriserer
hans System, træder i vore Dage mere eller mindre skarpt op overfor hverandre
og afføder forskellige Problemer. Denne Forskel mellem Problemer existerer
næsten ikke for \textsc{Spinoza}. Han søgte Tanker, ved hvis Hjælp de forskellige Problemer
paa én Gang kunde finde deres Løsning. Det er Problemernes Differentiation,
der især skiller os fra \textsc{Spinoza}. Men da der dog stedse maa forudsættes en indre
Sammenhæng mellem de forskellige Problemer, der sysselsætter menneskelig Tænken,
vil \textsc{Spinoza}'s store Tankeforsøg stedse bevare sin Betydning som orienterende, skønt
vi nu ikke kunde opføre Tankebygninger i samme Stil som hans.

% --- p. 14 ---
\setcounter{footnote}{0}%
\opage{14}\section*{Første Bog.}
\subsection*{a. Definitioner og Axiomer.}

De Definitioner og Axiomer, med hvilke Ethica begynder, er, som vi i det
Foregaaende har set, ikke opstillede vilkaarlig, men vundne ved Analyse dels af
overleverede religiøse Forestillinger, dels af den Fremgangsmaade, Tidens nye Videnskab
anvendte. Definitioner, siger \textsc{Spinoza} i et Brev (Ep. 9), er kun vilkaarlige,
naar de blot opstilles for at prøves, ikke naar de opstilles, fordi Erfaringen leder
os til dem. Vel kan Erfaring, siger han i et følgende Brev (Ep. 10), ikke lære os
Tingenes Væsen at kende, men den kan begrænse vor Tanke til visse væsentlige
Sider ved Tingene (mentem nostram determinare, ut circa certas tantum rerum essentias
cogitet).

Ved nærmere Betragtning finder vi i Ethica's første Definitioner og Axiomer en
Formulering af de Forudsætninger, paa hvilke al Videnskab ifølge \textsc{Spinoza} hviler.
Som vi i Indledningen har set, betragtes disse Forudsætninger af \textsc{Spinoza} ikke som
rent formale; de er Omskrivninger af den tilstrækkelige Grunds Princip, af det
strenge logiske Forhold mellem Grund og Følge, men den Sandhedsnorm, som de
udtrykker, tænkes som et existerende Væsen.

I den første Definition gaas ud fra, at der er Noget, hvis Existens man ikke
kan nægte uden at komme i Modsigelse med sig selv. Naar dette Noget kaldes
Aarsag til sig selv (causa sui), saa maa erindres, at han i Afhandlingen om Forstandens
Forbedring erklærer dette Udtryk for rent populært. Det hedder nemlig
der: „Hvis Noget er i sig selv, eller, som man populært siger (ut vulgo dicitur), er
Aarsag til sig selv (causa sui), saa maa det forstaas ved sit eget Væsen alene“.
\textsc{Descartes} havde allerede brugt dette Udtryk, og det var blevet indvendt imod
ham (i den første og den fjerde Gruppe af Indvendinger mod hans Meditationer),
at for at et Væsen skulde kunne være Aarsag til sig selv, maatte det være til, før
det existerede. Det er maaske derfor, at \textsc{Spinoza} kritiserer Udtrykket i sin erkendelsesteoretiske
Afhandling. Sagen er dog den, at saaledes som \textsc{Spinoza} bruger
Ordet Aarsag, spiller Tidsforholdet mellem Aarsag og Virkning tilsidst ingen Rolle.
Derfor bliver Grund (ratio) og Aarsag (causa) for ham et og det samme. Derfor
falder den første Definition sammen med den tredie, der definerer Substans som
„det, der er i sig selv og forstaas ved sig selv, det vil sige, det, hvis Begreb ikke
trænger til nogen anden Tings Begreb, af hvilket det skulde dannes“.

% --- p. 15 ---
\setcounter{footnote}{0}%
\opage{15}Hvad kan det nu være for et Begreb, der saaledes skulde kunne afslutte vor
Erkendelse uden at pege ud over sig selv? Som jeg har søgt at vise i Indledningen,
er det for \textsc{Spinoza} Principet for al Lovmæssighed i Verden --- det Princip, der
gør det muligt at forstaa Emnernes Rækkefølge som en Række af Grunde og Følger.
Tanken om Forholdet mellem Grund og Følge bliver den første (eller sidste) Tanke,
og dens Indhold er, for \textsc{Spinoza}, det Væsen, der ligger til Grund for alt, idet det
selv ikke har nogen Grund udenfor sig.

I den anden og sjette Definition skelnes der mellem Noget, som er endeligt i
sin Art, Noget, som er uendeligt i sin Art, og Noget, som er absolut uendeligt. Dermed
er Forholdet mellem Enkeltvæsen (Modus), Grundegenskab (Attribut) og Grundvæsen
(Substans) angivet.

Som vi har set i Indledningen, grunder Attributets Begreb sig paa, at der er
Rækkedannelser, som kan fortsættes i det uendelige efter samme Regel, som den,
hvormed de begyndte. Dette vises i første og anden Sætning i anden Bog for
Tænkningens og Udstrækningens Vedkommende. Der kan ingen Grænse sættes,
hverken for Tankens eller Udstrækningens Verden, og de kan heller ikke begrænse
hinanden; ti i saa Fald vilde jo Rækkedannelsens Lov forandres. Medens
der indenfor et Attribut ingen Grænse kan sættes, er Attributet selv dog forsaavidt
begrænset, som det forudsætter en bestemt Regel for Rækkedannelse. Derfor kan
det kun udtrykke en enkelt Side eller Egenskab ved Tilværelsen. Og der kan da
dannes en ny Række af Tilværelsens forskellige Attributer. Denne Række kan
heller ikke have nogen Grænse (skønt vi, som det siden vil vise sig, kun kender to
Led i den), og Gudsbegrebet faar netop sin Definition i Begrebet om den under
uendelig mange Attributer bestaaende Substans (Def. 6).

Gennem de Rækker, i hvilke Attributerne lægger sig for Dagen, udtrykker
Substansen sig. De enkelte Led (Modi) i disse Rækker holdes nemlig sammen ved
Begrundelsens Lov. Men der er (ifølge Def. 3 og 5) den Forskel mellem Substans
og Modus, at hin har sin Grund i sig selv, denne sin Grund i noget andet end
sig selv. Den fælles Forudsætning er, at en Grund maa der være, hvori den saa
ligger. Den tilstrækkelige Grunds Princip falder for \textsc{Spinoza} sammen med Aarsagsprincipet,
som vi allerede har set, og som det udtrykkelig indskærpes i fjerde og
femte Axiom: „Erkendelsen af Virkningen afhænger af Erkendelsen af Aarsagen
og forudsætter den. Ting, som ikke har noget fælles indbyrdes, kan heller ikke
forstaas ved hinanden, eller den enes Begreb forudsætter ikke den andens Begreb“.
Forholdet mellem Aarsag og Virkning er altsaa et rent Begrebsforhold, ved hvilket
Tidsforskellen ingen Rolle spiller.

\textsc{Spinoza} havde herefter egentlig ingen Grund til at opstille Aarsagssætningen
som en særlig Sætning. Vi finder dog, at han i speciellere Tilfælde, hvor Diskussionen
syntes at kræve det, udtrykkelig formulerer den. Saaledes Eth. I, 8. Schol. 2:
„Der gives nødvendigvis for enhver existerende Ting en Aarsag, paa Grund af hvilken
den existerer“. Her tillægges der altsaa Aarsagssætningen nødvendig Gyldighed,
og paa samme Maade udtrykker \textsc{Spinoza} sig i Ep. 34. Hvori denne Nødven
% --- p. 16 ---
\setcounter{footnote}{0}%
\opage{16}dighed ligger, siger \textsc{Spinoza} ikke, og at han ikke ser noget Problem heri, forstaas,
naar man fastholder, at Aarsag og Grund for ham er ét, saaledes som han udtrykkelig
siger i et Bevis for ellevte Sætning i første Bog: „For enhver Tings Existens
eller Ikkeexistens maa der angives en Aarsag eller Grund (causa seu ratio)“.

Men nu medfører alligevel den Forskel, der gøres mellem Substans og Modi
et dobbelt Aarsagsforhold, idet der er Forskel mellem at have sin Grund i sig selv
og at have den i noget andet. Eth. I, 28 (denne vigtige Sætning, som allerede
omtaltes i Indledningen, og som vi senere kommer tilbage til) kræver, at hver enkelt
Modus skal forklares ved en anden Modus, denne ved en tredie, og saaledes i
det uendelige. Men tillige skal hver enkelt Modus forstaas som en særegen Bestemmelse
(affectio) ved Substansen (Def. 5), og den maa derfor finde sin Forklaring i
denne --- foruden, at den skal finde sin Forklaring i andre Modi. Dette dobbelte
Aarsagsforhold --- man kunde kalde det det ideale og det elementære --- maatte for
\textsc{Spinoza} egentlig reduceres til et eneste, det ideale. Substansen er jo gennem sine
Attributer Princip for al Lovmæssighed i Verden, altsaa ogsaa for de Love, der
gælder mellem de forskellige Modi, og paa den anden Side forstaas disse Modi gennem
deres indbyrdes lovmæssige Indhold, altsaa ved at betragtes som Momenter i
Substansen (affectiones substantiæ). Desuden maa der ifølge Ax. 4 og 5 ikke være
Forskellighed mellem Aarsag og Virkning. Det elementære Aarsagsbegreb, der
skulde gælde for Forholdet mellem de indbyrdes forskellige Modi, kan altsaa ikke
være noget sandt Aarsagsforhold. Og tilsidst gælder det udvortes Forhold mellem
Modi, og dermed det elementære Aarsagsforhold, kun paa den sanselige Forestillings
(Imaginationens) Standpunkt. Og det sidste Problem bliver da, hvorledes de forskellige
Standpunkter, --- det, hvorfra der kun existerer det rent logiske (resp. matematiske)
Forhold mellem Grund og Følge i gensidig Ækvivalens, og det, for
hvilket der existerer forskellige Ting i udvortes Aarsagsordning, --- kan være mulige.
Reduktionen af disse to Standpunkter er endnu den Dag idag ikke gennemført\footnote[1]{Smgln. min Kritik af Marburgerskolen i \textit{Totalitet som Kategori} (Vid. Selsk. Skr. 1917, p 46---48.} --- % PRINTED AS IS: no ) and no stop after p
og kan sikkert ikke gennemføres.
--- En speciel Form for Aarsagssætningen, Inertisætningen, er indeholdt i den
anden Definition. Ti naar det her siges, at Udstrækning ikke kan begrænses af
Tanke, saa begrundes dette, som vi har set i Indledningen, ved, at vi kan tænke
os enhver Linie eller Flade eller ethvert Volumen forlænget eller forstørret, altsaa danne
Led i en kontinuerlig Række uden at forlade Udstrækningens Attribut. Men tillige
indeholder denne Sætning for \textsc{Spinoza} det metodiske Princip, at ethvert udstrakt
Emne skal forklares ved et andet udstrakt Emne, altsaa ethvert Legemes Forandringer
ved Forandringer i et andet Legemes Tilstand eller Stilling. Efter dette
Princip arbejder \textsc{Spinoza}, som vi vil se i det Følgende, overfor de specielle Spørgsmaal.
Og naar han forlanger (Eth. I, 10, til hvilken Sætning, foruden til Ax. 4,
Eth. II, 6 viser tilbage), at ethvert Attribut (og altsaa alt, hvad der hører ind under
dette Attribut) skal (debet) forklares ved det selv, saa følger det vel tilsidst af, at
ethvert Attribut, ifølge Def. 4, udtrykker Substansens Væsen og altsaa ligesom Sub
% --- p. 17 ---
\setcounter{footnote}{0}%
\opage{17}stansen maa have sin Grund i sig selv. Men selve denne Definition er i Grunden
et metodisk Princip, der begrundes ved, at der maa være Ensartethed mellem Aarsag
og Virkning, hvis Aarsagsforholdet skal være forstaaeligt. --- Der foreligger her
en interessant Forskel mellem den korte Traktat og Ethica. Hist opstilles allerede
Inertisætningen (II, 19, 8), men der drages endnu ikke Konsekvenser af den med
Hensyn til Metode (eller, efter \textsc{Spinoza}'s Terminologi, med Hensyn til Forholdet
mellem Attributerne), og der opstilles derfor ikke, som i Ethica en ny psykofysisk
Hypotese. \textsc{Descartes} havde i sin Formulering af Inertisætningen (Principia Philosophiæ
II, 37---41) udtrykkelig gjort et Forbehold med Hensyn til Engle og Sjæle.
Og i Henhold til et saadant Forbehold lader den korte Traktat trygt Sjælen paavirke
de fysiologiske Processers Retning, og lader paa den anden Side Legemet
paavirke Sjælen og derved erkendes af denne (K. Tr. II, 19, 9---14). Her er altsaa i
Mellemrummet mellem den korte Traktat og Ethica foregaaet en Ændring i \textsc{Spinoza}'s
Opfattelse, sandsynligvis under Indflydelse af det i Indledningen omtalte
metodiske Princip (vera causa).

Karakteristisk for \textsc{Spinoza} er det nu, at disse forskellige Principer (der iøvrigt
for ham ikke er forskellige) ikke opfattes som Former eller Betingelser for menneskelig
Forstaaelse, bestemte ved den menneskelige Tænknings Natur og ved Beskaffenheden
af de Opgaver, der stilles den, men betragtes som Udtryk for Tingenes
Væsen, tilsidst for den evige Tilværelsesgrunds Væsen. Naar en Tanke er sand,
hedder det i Ax. 6, maa den stemme med sin Genstand\footnote[1]{Dette er den rette Oversættelse af Ax. 6 (idea vera debet cum suo ideato convenire). I selve Tankens egen Lov er Sandhedsnormen given, der tillige afgør, hvad der er falsk (veritas norma sui et falsi, Eth. II, 43 Schol.). I Afhandlingen om Forstandens Forbedring udvikler \textsc{Spinoza}, at Sandhed ikke beror paa noget fra Erkendelsen forskelligt; enhver falsk Forestilling vil opløses gennem sine Konsekvenser. Vi kommer tilbage til dette Spørgsmaal senere. \textsc{--- Freudenthal} betragter (Festskrift til \textsc{Zeller} p. 128) med Urette Ax. 6 som Definition af Sandhed og finder deri en Rest af Skolastik.}. Og hvad der er Indhold
af streng Tænken ---
hedder det i Eth. I, 30 Dem.\footnote[2]{Ordet „objective“ betyder i \textsc{Spinoza}'s Sprogbrug „subjektivt“. Smlgn. nedenfor ved II, 7.} som nærmere Forklaring af
Ax. 6 --- maa nødvendigvis gives i Naturen.

At vi i vor Forsken efter Erkendelsens sidste Forudsætninger ikke kan komme
længere tilbage end til visse Principer, er for \textsc{Spinoza} Bevis for, at disse Principer
udtrykker et evigt og nødvendigt Væsens Natur, som ligger til Grund for alt (res
omnium prima, som det hedder i Afhandlingen om Forstandens Forbedring), og
som ikke selv har sin Grund i noget fra det forskelligt. Det er \textsc{Parmenides}' gamle
Lære, som stadig gør sig gældende: „Tanken og Tankens Genstand er ét og det
samme“ (ταὐτὸν δ᾽ἐστιν νοεῖν τε καὶ οὓνεκεν ἐστι νοήμα).

Vil vi vide mere om Substansen, henvises vi til Attributerne, der hvert for
sig udtrykker Substansens Væsen fra en enkelt Side (in suo genere). Attribut er
ifølge Def. 4 det, som Forstanden opfatter om Substansen som udgørende (konstituerende)
dens Væsen. Af denne Definition sluttede \textsc{Hegel}\footnote[3]{\textit{Vorlesungen über die Geschichte der Philosophie\textsuperscript{2}}, Berlin 1844, III, p. 344: „Wie kommt ausser Gott der Verstand herbeigelaufen, der die beiden Attribute des Denkens und der Ausdehnung auf die absolute Substanz anwendet?“}, og efter ham J. E.
% --- p. 18 ---
\setcounter{footnote}{0}%
\opage{18}\textsc{Erdmann}, at det kun er for vor Erkendelse, at der gives saadanne Attributer. Men
„som“ (tanquam) betyder her ikke „som om“. \textsc{Spinoza} bruger flere Steder „tanquam“
om gyldig Opfattelse (se f. Ex. Eth. II, 45 Schol.). Den kritiske Filosofis Sondring
mellem Tingene i sig selv og vor Opfattelse af dem tør ikke uden videre lægges ind
hos \textsc{Srinoza}. Han lærer jo udtrykkelig, at Attributet ligesom Substansen har sin % PRINTED AS IS: SRINOZA
Grund i sig selv. Hvis \textsc{Hegel} havde Ret, maatte det hedde, at Attributet havde
sin Grund i den menneskelige Forstand. \textsc{Spinoza}'s karakteristiske Lære gaar netop
ud paa, at Erkendelsens sidste Grundlag ogsaa er Tilværelsens sidste Grundlag.

I det Forhold, \textsc{Spinoza} har angivet mellem Begreberne Substans (Grundvæsen),
Attribut (Grundegenskab) og Modus (Enkeltvæsen, Emne), ligger egentlig hele hans
Filosofi indeholdt. Den specielle Række af Sætninger, han opstiller, er kun videre
Udviklinger af, hvad der indeholdes i Definitionerne. De fremkommer ikke, som en
virkelig „geometrisk“ Metode vilde kræve, ved kombinerende Konstruktion ud fra
de opstillede Definitioner, men er Explikationer, Omskrivninger, ikke egentlige
Demonstrationer. Den „geometriske“ Metode er blot en Skal, som maa sønderbrydes,
naar man vil have fat paa \textsc{Spinoza}'s egentlige, arbejdende Tanke. Og den har
den store Ulempe, som \textsc{Spinoza} selv undertiden føler, at de forskellige Læresætninger
stilles paa lige Linie trods deres indbyrdes Over- og Underordningsforhold.

\subsection*{Exkurs om Definitioner, Axiomer og Postulater hos \textsc{Spinoza}.}

{\small
Med Hensyn til \textsc{Spinoza}'s „Definitioner, Axiomer og Postulater“ bemærkes følgende.

1) Definitionerne opstilles som Navneforklaringer af Grundbegreber, men viser sig ved nærmere Undersøgelse at skyldes en Analyse af \textsc{Spinoza}'s Videnskabsbegreb. Denne Analyses Resultater slaas dogmatisk fast under den (ikke formulerede) Forudsætning, at Tænkningeus Grundlag ogsaa er Tilværelsens. % PRINTED AS IS: Tænkningeus
\par
2) Axiomerne har en forskellig Karakter i de forskellige Bøger. I første Bog udtales de Grundsætninger, der betragtes som selvfølgelige, men, ligesom Definitionerne, viser sig at bero paa Analyse af Videnskabsbegrebet. --- I anden, tredie og fjerde Bog er Axiomerne empiriske Sætninger, der for \textsc{Spinoza} staar som utvivlsomme. At „Mennesket tænker“ (II, Ax. 2), og at „vi ikke mærker eller opfatter andre Enkeltting end Legemer og Tanker“ (II, Ax. 5), staar fast som Erfaringskendsgerninger. Ligesaa, at enhver Enkeltting vil kunne finde sin Overmand (IV. Ax.). Efter II, 13 opstilles --- i Forbindelse med Laanesætninger (Lemmata) fra Naturvidenskaben --- Axiomer, der støtter sig til Erfaringen (f. Ex. at Legemer enten er i Hvile eller i Bevægelse o. s. v.). Disse Axiomer er de naturvidenskabelige Sidestykker til de i Bogens Begyndelse fra psykologisk Erfaring hentede Axiomer. --- Af de to Axiomer i femte Bog er det ene en empirsk-psykologisk Sætning, det andet en erkendelsesteoretisk Sætning, der udledes af en Sætning i tredie Bog. % PRINTED AS IS: empirsk-psykologisk
\par
3) Udtrykket Postulat bruger \textsc{Spinoza} kun i anden og tredie Bog, og han gør ingen principiel Forskel mellem Axiom og Postulat (smlgn. III, Ax. 1: postulatum seu axioma). Den Række Postulater, der opstilles i tredie Bog efter Sætning 13, er kun Resuméer af, hvad der indeholdtes i de forudgaaende Laanesætninger (Lemmata) og Axiomer. Til disse Postulater vises tilbage, naar det i Sætning 17 Schol. hedder: „Alle de Postulater, som jeg har opstillet, indeholder neppe Noget, der ikke staar fast ved Erfaring“.\par}

% --- p. 19 ---
\setcounter{footnote}{0}%
\opage{19}\subsection*{b. Sætning 1---14.}


\begin{center}\textit{Udrensning af Substansbegrebet og Paavisning af dets Identitet med Gudsbegrebet.}\end{center}

\textit{α.} At der maa være Substans i Verden, d. v. s. Noget, der bestaar i sig selv
uden at forudsætte noget andet, er for \textsc{Spinoza} en Selvfølge, fordi vore Tankers
(Grundenes og Aarsagernes) Række maa finde Afslutning i Noget, som ikke igen
kræver fortsat Tænkning. Trangen til Afslutning viser sig, som allerede omtalt i
Indledningen, paa en anden Maade i Ethica end i den korte Traktat. I Ungdomsskriftet
betonedes især Trangen til Ro. „Hvis vi“, hedder det her (I, 2), „vilde søge
Aarsagen til Substansen, der er Principet for alle Ting, saa maatte vi igen søge
Aarsagen til denne Aarsag, og saaledes i det uendelige, saa at det --- naar vi nødvendigvis
maa stanse etsteds og komme til Ro, hvilket vi maa --- er nødvendigt at
komme til Ro ved denne ene Substans“. Denne Tankegang har \textsc{Spinoza} senere
udtrykkelig forladt. Begrundelsen af Substansens eller Guds Existens, siger han i
et Brev (Ep. 12), ligger ikke i, at vi ellers vilde faa en uendelig Række af Aarsager,
men deri, at der ellers ikke vilde være nogen i sig nødvendig Aarsag til de ikke
nødvendige Ting\footnote[1]{En lignende Tankegang har allerede \textsc{Aristoteles}: „Hvis der ikke er et Første, er der overhovedet ingen Aarsag“. (\textit{Metaphys.} p. 994).}. At Rækken af Modi (Enkeltvæsener, Emner) fører os fra A til B
og fra B til C o. s. v. i det uendelige, --- deri finder \textsc{Spinoza} nu ingen Vanskelighed.
Tværtimod er han --- ligesom \textsc{Bruno}\footnote[2]{Se \textit{Den nyere Filosofis Historie\textsuperscript{2}}. I, p. 123 f.} overbevist om, at der af Substansen
(Naturen, Guddommen) maa fremgaa en Uendelighed af Følger eller Virkninger.
Saalænge Tankearbejdet sysler med Modi, kan der altsaa ikke ventes nogen Afslutning.
Skal en Afslutning naas, maa den findes i det Princip, der ligger til Grund
for den lovmæssige Sammenhæng mellem alle Modi, saavel som for selve disse Modi
efter deres inderste Natur, og det er netop dette Princip, \textsc{Spinoza} kalder Substans,
Natur eller Gud. Ikke Hviletrang, men absolut rationel Tilfredsstillelse motiverer
nu Afslutningen.

\textsc{Spinoza} tænker sig altsaa ikke Aarsagernes Række afsluttet med et første Led,
ud af hvilket alle andre Led udspringer, men han finder den Afslutning, han søger,
i den indre Lov eller Kraft, som kundgør sig i de Aarsagsrækker, Erfaringsvidenskaben
paaviser. Enten existerer der Intet, hedder det i Eth. I, 11 Dem., eller ogsaa
existerer der et absolut nødvendigt Væsen. Og paa den anden Side gælder det
ogsaa, at jo mere vi forstaar Enkeltvæsenerne, des mere forstaar vi Gud (Eth. V, 24).

\textit{β.} Naar Substansens Existens staar for \textsc{Spinoza} som en evig Sandhed, er det
altsaa, fordi den udtrykker den Grundsandhed, hvoraf al anden Sandhed afhænger,
--- den nemlig, at der i Tilværelsen raader en lovmæssig Orden, som finder sit
Udtryk i Aarsagsbegrebet, der for \textsc{Spinoza} falder sammen med Begrebet Grund.

Som allerede bemærket, opstilles Aarsagssætningen udtrykkelig i andet Skolium
til ottende Sætning. Men denne Sætning forudsættes tydelig i de første otte Sætninger.
Ti naar \textsc{Spinoza} vil godtgøre, at der, ifølge Definitionen af Substansen, kun
% --- p. 20 ---
\setcounter{footnote}{0}%
\opage{20}kan gives en eneste Substans, sker det paa følgende Maade: Hvis der var flere
Substanser, maatte de have forskellige Attributer, ti gennem Attributer udtrykkes
Substansers Væsen; men Substanser med forskellige Attributer vilde ikke kunne
være Aarsager til eller Virkninger af hverandre, eftersom de intet fælles vilde have;
deres Begreber vilde ikke kunne føres tilbage til et og samme Begreb.

En Forudsætning for denne Tankegangs Gyldighed er der, som \textsc{Spinoza} ikke
antyder, den nemlig, at forskellige Substanser ikke kan tænkes uden at staa i Vexelforhold
til hverandre. Dette er ikke nødvendigt, da jo enhver Substans skal have
sin Grund eller Aarsag til sig selv. Det kunde jo tænkes --- og \textsc{Leibniz} har i sin
Monadelære tænkt sig denne Mulighed --- at hvad der udspringer af den ene Substans,
slet ikke kommer til at krydses med eller staa i Vexelvirkningsforhold til,
hvad der udspringer af en anden Substans. De forskellige Substanser kunde staa i
samme Forhold til hverandre, som hos \textsc{Spinoza} de forskellige Attributer staar i
indbyrdes.

Hvis man definerer Aarsagssætningen saaledes, at den udsiger, at alt, hvad vi
skal kunne erkende, maa staa i Forhold til noget andet som Grund til Følge (resp.
Aarsag til Virkning) eller omvendt, vilde det være rigtigt, at hvis der gaves flere
„Substanser“, maatte de staa i Vexelvirkning med hverandre. Men saa vilde Begrebet
Aarsag eller Grund til sig selv være faldet bort. Dette Begreb betegner i
Virkeligheden en blind Gyde og er et rent negativt Begreb, som kun udsiger, at
man ikke har fundet Grunden eller Aarsagen til noget og derfor ikke kan bringe
det i Sammenhæng med, hvad man ellers véd eller mener at vide. De første
Forudsætninger for vor Verdensopfattelse kan kun retfærdiggøres gennem deres
Konsekvenser, ikke selv igen udledes. De Begreber, de indeholder, er derfor Grænsebegreber.
--- Denne Betragtning vilde \textsc{Spinoza} ikke kunne gaa ind paa. Den gennemføres
først af den kritiske Filosofi.

Det har sin Interesse at sammenligne \textsc{Spinoza}'s Substansbegreb med Skolastikens
og \textsc{Descartes} Substansbegreber paa den ene Side, og med \textsc{Kant}'s „Ding an
sich“ paa den anden Side.

Det ældre Substansbegreb (Skolastikens og \textsc{Descartes}') stiller Substansen som
bestaaende i sig selv i skarp Modsætning til Egenskaberne, og den faar derved en
mystisk Karakter, da vi jo kun kender Tingene gennem deres Egenskaber og disses
indbyrdes Sammenhæng. Hos \textsc{Spinoza} er der derimod et indre Forhold mellem
Substansen og dens Egenskaber (Attributer), idet disse hver for sig udtrykker Substansens
Væsen. Hvad der gælder for Substansen, at den har sin Grund i sig selv,
det gælder ogsaa for hvert enkelt Attribut. Og hvis den ovenfor givne Tydning af
\textsc{Spinoza}'s Substans som selve Begrundelsens eller Kundskabens Princip er rigtig,
indeholder Attributbegrebet det videnskabelige Princip, at Naturens Omraader (det
aandelige som det legemlige) skal forklares hvert efter sine egne Love. Naturen
skal forklares ved Naturen selv. --- Et saadant indre Forhold mellem Teologi og
Videnskab, mellem Gudserkendelse og Naturerkendelse kunde Skolastiken ikke anerkende.
Gud kan ifølge den end ikke staa i realt Forhold til den Verden, han har
% --- p. 21 ---
\setcounter{footnote}{0}%
\opage{21}skabt. Verden staar i Forhold til Gud, nemlig i et Afhængighedsforhold; men Gud
staar ikke i Forhold til Verden. Skabelsen har ingen Betydning for Guds eget Væsen\footnote[1]{Smgln. min \textit{Religionsfilosofi} § 22 og de i Noten anførte Steder af \textsc{Hugo de St. Victore} og \textsc{Thomas af Aqvino}. Smgln. om \textsc{Augustinus} Slutningen af § 78.}.
\textsc{Dunin}-\textsc{Borkowski} finder da ogsaa den afgørende Forskel mellem \textsc{Spinoza} og
Skolastiken deri, at efter hin er der et Gensidighedsforhold mellem det Uendelige
og det Endelige (Gud og Verden), efter denne staar vel Verden i et Forhold til Gud,
men kun ad Analogiens Vej (det vil i Grunden her sige: billedlig) kan man tale om,
at Gud staar i Forhold til Verden\footnote[2]{\textit{Der junge Spinoza}, p. 354: „Dem Verhältnis der Abhängigkeit des Geschöpfes vom Schöpfer entspricht in der unendlichen Ursache kein wirkliches Gegenverhältnis“. \textsc{Dunin}-\textsc{Borkowski} tillægger denne Forskel „eine ungeheure Tragweite“.}.

\textsc{Kant}'s „Ding an sich“ er et Restbegreb, udtrykker, at der stedse bliver Noget
tilbage, naar det videnskabelige Arbejde slutter. Det er tillige et Grænsebegreb og
et problematisk Begreb, og der vil --- hvad \textsc{Kant} ikke altid fastholder --- stedse
kunne rejses Tvivl om, hvorvidt Grænsen er bestemt paa rette Maade. Som vi
har set i Indledningen, indrømmer \textsc{Spinoza}, at der stedse bliver en Rest tilbage;
kun hævder han, og med Rette, at dette ikke berøver den Erkendelse, som virkelig
er vunden, dens Betydning. Sandhedskriteriet er for ham (som vi i en senere Sammenhæng
vil se) Tankens Overensstemmelse med sig selv indenfor Erfaringens
Verden.

\textit{γ.} Begrebet Attribut betegner en Egenskab, i hvilken Substansens Væsen som
bestaaende i sig selv og forstaaelig ved sig selv lægger sig for Dagen. \textsc{Spinoza} indskærper
(10 Schol.), hvad allerede \textsc{Arnauld} og \textsc{Regius} havde hævdet overfor \textsc{Descartes},
at man ikke deraf, at to Attributer er indbyrdes irreduktible, har Ret til
at slutte, at der bag dem ligger forskellige Væsener eller Substanser. Tværtimod, jo
mere Realitet eller Fuldkommenhed et Væsen besidder, des flere forskellige Attributer
maa det besidde, og Gudsbegrebet er ifølge \textsc{Spinoza} Begrebet om en Substans,
der har uendelig mange Attributer, af hvilke ethvert udtrykker Evighed og Uendelighed
(Def. 6 og Sætning 11). Beviset for Guds Existens falder sammen med Beviset
for Substansens Existens, kun at ved Gudsbegrebet Tanken om Realitet eller Fuldkommenhed
er af afgørende Betydning. Det staar for \textsc{Spinoza} som absurd, at det
Fuldkomne ikke skulde existere, --- at der kun skulde existere endelige, begrænsede
Væsener med faa Egenskaber, men ikke et uendeligt Væsen med uendelig mange
Egenskaber. „Fuldkommenhed“, siger han (11 Schol.), „ophæver ikke en Tings
Existens, men begrunder den netop, medens Ufuldkommenhed ophæver den. Derfor
kan vi ikke være saa sikre paa nogen Tings Existens som paa Guds, det absolut
uendelige og fuldkomne Væsens Existens“. I den sjette Definition i anden Bog erklæres
udtrykkelig Realitet og Fuldkommenhed for ét og det samme, og i Tillægget
til første Bog forkaster \textsc{Spinoza} den Maalestok for Fuldkommenhed, der hentes fra
Tingenes Forhold til menneskelig Trang og menneskelige Ønsker: „Tingenes Fuldkommenhed
maa ene vurderes efter deres Natur og deres Magt“.

% --- p. 22 ---
\setcounter{footnote}{0}%
\opage{22}Der er nu alligevel en Vurdering skjult i \textsc{Spinoza}'s Begreb om Fuldkommenhed,
selv om det gøres til ét med Begrebet Realitet. Der udtrykker sig netop i
denne Identificeren en dyb Beundring for Tilværelsens Fylde som den, der overskrider
alle Grænser baade hvad Kvalitetsrigdom, Tid og Rum angaar. Der rører
sig her den samme Begejstring for Naturens Rigdom og Ubegrænsethed, der hos
\textsc{Montaigne} laa bag hans tilsyneladende Skepsis saavel som bag hans Sans for
Individualiteternes Mangfoldighed, --- der fik \textsc{Bruno} til at sprænge det gamle Verdensbilledes
Skranker, --- og som ogsaa besjælede de store italienske Kunstnere.
\textsc{Leonardo da Vinci} var, som \textsc{Julius} \textsc{Lange} har sagt\footnote[1]{\textit{Udvalgte Skrifter}, II, p. 239. --- For \textsc{Montaigne}'s og \textsc{Bruno}'s Vedkommende henvises til \textit{Den nyere Filosofis Historie\textsuperscript{2}}. I, p. 26; 30 og p. 111; 117---123.}, „forelsket i Verdens Mangfoldighed“.
Der er her ogsaa et æstetisk Element i \textsc{Spinoza}'s Verdensanskuelse.
Han beundrer den Naturens Magt og Fylde, der har kunnet lade baade det store
og det smaa, baade det gode og det onde, baade det gavnlige og det skadelige
blive til (se Slutningsordene af Tillægget til første Bog). Han gør da ogsaa ikke
blot Realitet til ét med Fuldkommenhed, men erklærer, at Fuldkommenhed og
Skønhed er nær beslægtede Begreber (Ep. 54).

\textsc{Spinoza} afviser egentlig kun Værdibegrebet, naar man vil bruge det til Forklaring
af enkelte, specielle Naturfænomener. Derimod anvender han det paa Naturens
inderste Grund og højeste Lov, naar han identificerer Gud og Substansen og
erklærer Realitet og Fuldkommenhed for ét. Vurderingstrangen er ikke forsvunden
hos \textsc{Spinoza}, men den har en Gang for alle lagt sig til Ro i Grundtanken om Tilværelsens
store Lovmæssighed som bestemmende for alt, hvad der for endelige Væsener
kommer til at staa som godt eller ondt. For \textsc{Spinoza} personlig var, som de
selvbiografiske Antydninger i Begyndelsen af Afhandlingen om „Forstandens Forbedring“
viser, Sandhedserkendelse kommen til at staa som det højeste menneskelige
Gode, og i Slutningen af femte Bog begrundes „den intellektuelle Kærlighed til
Gud“ ved, at det er den guddommelige Naturorden, der gør saadan Erkendelse
mulig. Her var der for \textsc{Spinoza} en Værdi, der maatte have sin Grund dybt i Tilværelsens
Natur.

\textit{δ.} I Antagelsen af uendelig mange Attributer viser sig paa ejendommelig
Maade Forholdet mellem \textsc{Spinoza}'s religiøs-spekulative Tankegang og den Erfaringserkendelse,
han stedse lægger Vægt paa. Ti \textsc{Spinoza} indrømmer, som vi allerede
har set i Indledningen, at han ikke kender alle guddommelige Attributer, end ikke
største Delen. I femte Axiom i anden Bog slaas det fast, at vi ikke mærker eller
opfatter andre Emner (Modi) end saadanne, som hører ind under Tænkningens eller
Udstrækningens Attributer, altsaa Sjæle og Legemer, og i Ep. 64 forklares, at dette
kommer deraf, at den menneskelige Natur hverken forudsætter eller frembyder
andre Attributer end disse to. Naar \textsc{Spinoza} alligevel (Ep. 55) tillægger sig en
exakt Viden om Gud, maa det være, fordi han er overbevist om, at det logiske
Forhold mellem Grund og Følge er det samme indenfor alle Attributer, hvad der
jo, ifølge den i det foregaaende givne Tydning, ligger i, at de alle er Attributer
% --- p. 23 ---
\setcounter{footnote}{0}%
\opage{23}ved samme Substans. Kun en uendelig Forstand (intellectus infinitus) vilde kunne
omfatte Tanker, som udtrykker de uendelig mange Attributer. At \textsc{Spinoza} antager
en saadan intellectus infinitus, vil vi senere komme til at se.

Selve Ideen om en uendelig Mængde Attributer er af nyplatonisk Oprindelse.
For \textsc{Plotinos} var Tilværelsesprincipet en absolut Enhed, som indesluttede en Fylde
i sig, der ikke kunde udtømmes i nogen Existens eller fattes af nogen Tanke. I
Middelalderens Skolastik laa denne Tanke stedse i Baggrunden, men brydes der
med den skolastiske Teologis rationalistiske Karakter. Efterat \textsc{Descartes} skarpt
havde afgrænset Sjæleliv og Materie som de to Omraader for vor Forsken, kom
naturlig det Spørgsmaal frem, med hvilken Ret vi antager, at disse to Omraader
fuldstændig udtrykker Tilværelsens Indhold. I et Skrift, der omtrent er samtidigt
med \textsc{Spinoza}'s Ethica, kommer \textsc{Malebranche} ind paa dette Spørgsmaal. „Man tør
ikke paastaa“, siger han (Recherche de la vérité. III, 2, 9), „at der kun gives Aander
og Legemer, Væsener, som kan tænke, og Væsener, som er udstrakte. Det
kunde være en Fejltagelse. Ti skønt de er tilstrækkelige til at forklare Naturen, og
skønt man derfor, uden at frygte for at tage fejl, kan slutte, at de Naturting, vi kan
erkende, afhænger af Udstrækning og af Tanke, kan det dog være, at der gives
andre Væsener, om hvilke vi ingen Forestilling har, og hvis Virkninger vi ikke ser.
Det er da en forhastet Dom, Menneskene fælder, naar de som utvivlsomt Princip
opstiller, at ethvert Væsen enten maa være en Sjæl eller et Legeme“. --- Forskellen
mellem \textsc{Malebranche} og \textsc{Spinoza} paa dette Punkt er kun, at hvor \textsc{Malebranche}
som Cartesianer taler om Væsener eller Substanser, taler \textsc{Spinoza} om Attributer ved
en og samme Substans.

Der er ikke, som man undertiden har ment, nogen Selvmodsigelse i Antagelsen
af uendelig mange Attributer. Men et positivt Bevis for denne Antagelse vilde
\textsc{Spinoza} ikke kunne give. Hans skarpsindigste Korrespondent, \textsc{Tschirnhausen}, opfordrede
ham da ogsaa forgæves til at give et saadant. Ligesaalidt har han kunnet
give en Forklaring af, hvorledes Attributernes Mangfoldighed kan udledes af Substansens
Enhed. Denne Vanskelighed er dog den samme, hvad enten der kun antages
to, eller der antages uendelig mange Attributer. Endnu en Vanskelighed berører
\textsc{Tschirnhausen}, nemlig hvoraf det vides, at alle Attributer staar paa lige Trin.
Disse Spørgsmaal af \textsc{Tschirnhausen} (Ep. 65---66; 70) staar endnu ved Magt. I et
tidligere Brev (Ep. 9) til en yngre Ven havde \textsc{Spinoza} erklæret, at han havde to
Beviser for Attributernes Mangfoldighed --- det ene, at jo mere Realitet et Væsen har,
des flere Attributer maa der tillægges det; --- det andet, at jo flere Attributer et
Væsen har, des mere nødes vi til at antage det for virkeligt. --- Overfor begge disse
„Beviser“ staar dog den af \textsc{Spinoza} selv anerkendte Kendsgerntng, at vi nu engang % PRINTED AS IS: Kendsgerntng
kun kender to Attributer, og at den sikreste Erkendelse af Virkelighed, vi har, er
opbygget paa Grundlag af dem. --- I den korte Traktat appellerede \textsc{Spinoza} til en
Anelse, vi har om, at der existerer andre Sider ved Tilværelsen end de, vi kender;
og han havde ikke opgivet Haabet om at lære dem at kende (K. Tr. I, 1, 7. Note og
VII, 1, 1. Note). I Brevene og i Ethica optræder denne Tanke ikke.

% --- p. 24 ---
\setcounter{footnote}{0}%
\opage{24}\textsc{Spinoza} kunde ganske simpelt have holdt sig til, at det Aandelige og det Materielle
ikke staar i kontradiktorisk, men kun i kontrært Forhold til hinanden, idet
vi ingen Garanti har for, at alle Sider ved Tilværelsen skulde fremtræde indenfor
vor Erfaring. Ligesaalidt som vor Erfaring kan siges at være udtømmende med
Hensyn til Antallet af Emner (Modi), aandelige og legemlige, ligesaalidt tør den siges
at være udtømmende med Hensyn til Antallet af de Hovedarter af Existens, den
aabenbarer os. --- Denne Betragtning er, som jeg andetsteds\footnote[1]{\textit{Den menneskelige Tanke}, p. 309 f; 334---338.} har søgt at vise, ikke
uden Betydning for vor Stilling overfor Problemet om Forholdet mellem Aand og
Materie.

Hvad de uendelig mange Attributers Ligestillethed angaar, beror den ifølge
\textsc{Spinoza} paa, at Substansen i dem alle udtrykkes efter sit Væsen. Hvis nu Substansens
Væsen bestaar i selve det logiske Forhold mellem Grund og Følge, kan
Ligestilletheden kun gælde, forsaavidt som alle Attributer har en rationel Karakter,
indeholder Mulighed for Anvendelse af Logik. Men der er i og for sig intet i Vejen
for, at de iøvrigt, med Hensyn til Beskaffenhed og Værdi kan være højst forskellige
--- med mindre (som \textsc{Spinoza} jo tenderer til) al Beskaffenhed reduceres til Logik og
al Værdi til „Realitet“. Tanken om uendelig mange Attributer aabner en langt
større Ventil for Muligheder, end \textsc{Spinoza} har været sig bevidst. Den Kendsgerning,
at vi kun kender to Attributer, sætter en snevrere Grænse for en videnskabelig Verdensanskuelse,
end han saa.

\subsection*{c. Sætning 15---28.}
\begin{center}\textit{Nærmere Bestemmelser af Gudsbegrebet i dets Forhold til Verdensbegrebet.}\end{center}
\begin{center}\textit{(Enhed og Mangfoldighed).}\end{center}

\textit{α.} Naar Gud er den under uendelig mange Attributer bestaaende Substans,
og naar Udstrækning er et af disse, saa er Udstrækning, den Grundform, i hvilken
alt Materielt fremtræder, en Form for Guds Væsen. Naar dette forundrer eller forarger
mange, er det, mener \textsc{Spinoza}, fordi de ved Materie forstaar et Kaos af sondrede
Elementer. Men det er kun for Fantasien, den sanselige Forestillen (imaginatio),
at der existerer sondrede Elementer, ligesom det kun er den, der lader Udstrækningen
staa som begrænset. Udstrækningen er kontinuerlig, og derfor ubegrænset og
udelelig\footnote[2]{\textsc{Huan} (\textit{Le Dieu de Spinoza}, p. 14) vender Forholdet om, idet han mener, at \textsc{Spinoza} støtter sig til Udeleligheden for at bevise, at Udstrækning er et guddommeligt Attribut. Men \textsc{Spinoza}'s Polemik mod Materiens Delelighed skal kun bortrydde en Hindring, ikke være et Bevis. Udeleligheden følger, ligesom Uendeligheden, af Kontinuiteten. Smlgn. det ovenfor om Rækkedannelser bemærkede.}. At Materien, hvis Grundform er Udstrækning i Rummet, er et guddommeligt
Attribut, betyder derfor ikke, at Gud er ét med Summen af de materielle
Ting, Sanserne viser os (15 Schol.). Det er Udstrækning som Geometriens Genstand,
\textsc{Spinoza} giver en principiel Plads i sin Filosofi, idet han ledes af den afgørende
Betydning, Geometrien har for al Naturvidenskab. Desuden gør sig Umuligheden
% --- p. 25 ---
\setcounter{footnote}{0}%
\opage{25}gældende af at udlede det Materielle af det Aandelige, hvis man ikke antager en
ubegribelig Skaberakt. Med lige Nødvendighed aabenbarer ifølge \textsc{Spinoza} Gud sig
i den materielle og i den aandelige Natur. \textsc{Spinoza}'s Udstrækningsattribut svarer
til, hvad \textsc{Descartes} i Oversigten over sine Meditationer (Synopsis Meditationum)
kalder „Legemet i Almindelighed“ (corpus generaliter sumptum) i Modsætning til
de enkelte Legemer. Som \textsc{Dunin}-\textsc{Borkowski} viser\footnote[1]{\textit{Der junge Spinoza}, p. 357---360; 573.}, gøres der hos Skolastikerne
fra Middelalderens Slutning Forsøg paa at danne et Begreb om Existens i Rummet,
der kunde passe baade paa Gud og paa endelige Væsener; men hos dem var der
kun Tale om Analogi mellem guddommeligt og materielt Rum, ikke om Identitet
af dem. --- Om Forholdet mellem \textsc{Spinoza} og \textsc{Newton} med Hensyn til Udstrækningens
Begreb er der talt i Indledningen.

Som det maatte ventes, er \textsc{Spinoza}'s Opfattelse af Udstrækningens Attribut i
dets Forhold til de Emner, Sansningen viser os, analog med hans Opfattelse af
Tænkningens Attribut i dets Forhold til de Emner, psykologisk Iagttagelse viser os.
Ligesom han advarer imod at gøre Materie som guddommeligt Attribut til ét med
sanselig Materie, saaledes advarer han ogsaa imod at opfatte Tænkningens Attribut
som ét med de sjælelige Tilstande og Virksomheder, psykologisk Iagttagelse lærer
os at kende. Forstand og Villie kan ikke høre til Guds Væsen i den Betydning, i
hvilken de kan bruges om Mennesker. Menneskelig Forstand lægger sig for Dagen
i tænkende Bearbejdelse og Iagttagelse og forudsætter altsaa et givet Stof; men
for Gud kan der ikke gives noget fra ham og hans Tanke forskelligt, som nu først
skulde bearbejdes i Tankens Form: saa vilde Gud jo være afhængig og begrænset!
Og menneskelig Villen stræber stedse efter et Maal, som endnu ikke er naaet; en
saadan Stræben kan kun tænkes hos endelige Væsener. For Gud staar alt i evig
Fuldendelse, uden Tidsfølge (simul natura). Kun Navnet kunde guddommelig og
menneskelig Forstand og Villie have fælles --- ligesom Hundestjernen og Hunden,
det gøende Dyr paa Jorden (17 Schol.). --- Det er af Interesse at se, at \textsc{Kant} i alle
sine tre Hovedværker kommer til samme Resultat. Ogsaa for ham bliver der kun
et Ord tilbage, naar man vil anvende psykologiske Udtryk om Gud\footnote[2]{\textit{Den nyere Filosofis Historie\textsuperscript{2}}, II, p. 95.}. Den grundlæggende
kritiske Filosofi kommer altsaa her til samme Slutning som den mest
gennemførte dogmatiske Filosofi.

Naar \textsc{Spinoza} alligevel opfatter „Tænkning“ som et guddommeligt Attribut,
foretager han en lignende Reduktion og Idealisering af aandelige Emner, som han
i Attributet „Udstrækning“ foretager af materielle Emner. Begge Steder maa det
forstaas saaledes, at Attributet betyder den Lovmæssighed, der gælder indenfor en
af de væsentlige Sider, fra hvilke Tilværelsen kan ses, --- en Lovmæssighed, der er
Betingelsen for de enkelte Emners Existens. Tanken som guddommeligt Attribut
betyder, at Erkendelsens Ideal existerer, ikke først skal opnaas, men er en evig
Realitet, som det er det sandhedssøgende Menneskes Opgave at faa Del i. Naar
\textsc{Spinoza} i sin teologisk-politiske Afhandling (Kap. 14) siger, at de positive Religioners
% --- p. 26 ---
\setcounter{footnote}{0}%
\opage{26}Gud har Betydning som praktisk Forbillede (exemplar veræ vitæ), saa gælder egentlig,
skønt \textsc{Spinoza} ikke vilde anerkende det, det Samme for hans Filosofis Gud.
Under Tænkningens Attribut er \textsc{Spinoza}'s Gud det evige Sandhedsideal, og tillige,
ved at have Grunden i sig selv, det evige Frihedsideal.

\textit{β.} At alle Enkeltvæsener (Modi, res particulares), der fremtræder som Emner
i vor Erfaring, udspringer af Guds Væsen som nærmere Bestemmelser ved dette
(affectiones substantiæ), begrunder \textsc{Spinoza} ved, at de forudsætter Substansen og
Attributerne (15---16; 25). Det er, som vi tidligere har set, den videnskabelige Undersøgelse
af de aandelige og de materielle Fænomener, der fører til Erkendelsen af
Naturens Lovmæssighed og derigennem til Erkendelsen af Gud eller Substansen.
Men --- hævdes det nu i Sætning 28 --- hvad der er endeligt og har en begrænset
Existens, kan ikke udledes af noget guddommeligt Attributs absolute Natur, ti af
denne kan kun følge, hvad der selv er uendeligt. Et endeligt Emne finder kun sin
Forklaring ved et andet endeligt Emne, dette ved et tredie, og saaledes i det Uendelige.
--- Det Princip, som \textsc{Spinoza} her hævder, falder, som allerede udtalt i Indledningen,
sammen med hvad man kunde kalde den gyldige Aarsags Princip, \textsc{Kepler}'s
vera causa. Hvad der fremtræder i Erfaringen, skal forklares ved en Aarsag,
der selv kan paavises i Erfaringen\footnote[1]{Eth. II, 17 Schol. bruger \textsc{Spinoza} Udtrykket vera causa paa denne Maade. Ellers bruger han Udtrykket causa proxima.}. \textsc{Galilei} gjorde i sine Dialoger Anvendelse af
dette Princip, naar han hævdede, at man ikke midt i en speciel Undersøgelse (f. Ex.
om Jordens Forhold til Solen) kunde paaberaabe sig en Indgriben af Gud. Og hans
Begrundelse er den samme som \textsc{Spinoza}'s: af en uendelig Aarsag følger uendelige
Virkninger; Gud er Aarsag til alt, men det gælder i det enkelte Tilfælde at finde
det specielle Aarsagsforhold.

Den Modsætning, for ikke at sige Modsigelse, der unægtelig er mellem Sætning
15 og Sætning 28, kan indenfor \textsc{Spinoza}'s Forudsætninger hæves, naar man lægger
den Opfattelse af \textsc{Spinoza}'s Substansbegreb til Grund, som er fremsat i det foregaaende.
Naar et i Erfaringen givet Emne kan forklares ved et andet, og der saaledes
kan dannes en uendelig Række af endelige Aarsager og Virkninger, saa ytrer
der sig netop i den Lov, som forbinder Leddene i en saadan Række, et Princip, en
Tingenes Orden, der ikke selv er Led i Rækken og ikke selv fremtræder som
Emne for Iagttagelse. Og det er et saadant Princip, \textsc{Spinoza} har udtrykt i sit Substansbegreb
og i det dermed partielt identiske Attributbegreb. Ved denne Opfattelse
kan de to Aarsagsbegreber, med hvilke \textsc{Spinoza} opererer, forbindes. At et Emne
forklares ved sin Plads i Emnernes Række, vil sige det samme, som at det forklares
ved Substansen (resp. Attributet). Derved forstaas en Sætning som Eth. V, 24:
„Jo mere vi forstaar de enkelte Ting, des mere forstaar vi Gud“.

Det forekommer mig, at den Overensstemmelse (eller dog Modsigelsesfrihed),
som paa denne Maade bliver mulig mellem de to vigtige Sætninger 15 og 28, taler
til Fordel for den Opfattelse af \textsc{Spinoza}'s Substansbegreb, som jeg har gjort gældende.
Men der bliver stadig en Vanskelighed tilbage, som følger af Identificerin%
% --- p. 27 ---
\setcounter{footnote}{0}%
\opage{27}gen af Grund og Aarsag i \textsc{Spinoza}'s System. Forholdet mellem Grund og Følge er
tidløst („evigt“); Forholdet mellem Aarsag og Virkning forudsætter (forsaavidt det er
forskelligt fra Forholdet mellem Grund og Følge) et Tidsforhold. Der er et Tidsforhold,
ikke blot et Begrundelsesforhold mellem Leddene i Aarsagsrækken. Derved
bestaar der alligevel en principiel Forskel mellem Modi (Emner) som Led i de Aarsagsrækker,
der kan bygges paa Erfaringens Grund, og det partielle Identitetsforhold,
i hvilket den enkelte Modus staar til Substansen. Desuden har de enkelte
Modi hver sin kvalitative Ejendommelighed. Problemet er hos \textsc{Spinoza}, ligesom for
moderne erkendelsesteoretiske Standpunkter, Forholdet mellem Logikens og Matematikens
tids- og kvalitetsløse Verden og Erfaringens Verden, som er de kvalitative
Forandringers Verden. Denne Verden forudsætter hin, naar den skal forstaas, men
kan ikke udledes af den. Allerede \textsc{Tschirnhausen} gjorde opmærksom paa, at Bevægelse
og Figur ikke kan udledes af Udstrækningens almene Attribut, og at overhovedet
Tingenes Forskellighed og Mangfoldighed ikke kan udledes af Substansen
eller Attributet i deres absolute Enhed. (Ep. 59. 80. 82).

At \textsc{Spinoza} har mærket Problemet, ses af, at han indskyder et Mellembegreb
mellem Substansen (og Attributerne) paa den ene Side, de konkrete Emner (Modi)
paa den anden Side. Dette Mellembegreb er Begrebet Modus infinitus (23). Uendelige
Modi udspringer (se 28 Schol.) umiddelbart eller middelbart af Substansens
eller Attributernes Væsen og er evige ligesom de. Som Exempler nævner han under
Tænkningens Attribut den uendelige Forstand og under Udstrækningens Attribut
det konstante Forhold mellem Hvile og Bevægelse. (Eth. I, 21. 30. V, 23. 40 Schol.
cfr. Ep. 32 og 64). Allerede i den korte Traktat fremtræder disse Begreber, og \textsc{Spinoza}
kalder der (I, 2), ligesom i Ep. 73, deres Genstande for „Gudesønner“ for at
betegne, at de staar i nærere Forhold til Substansen end de endelige Modi.

Forsøger vi at oversætte disse halvt mytologiske Udtryk paa Erkendelsesteoriens
Sprog, kan \textsc{Spinoza}'s Opfattelse siges at være den, at Tankelivets og Bevægelsens
Bestaaen i Verden under alle Omskiftelser er en Konsekvens af Begrundelsens
eller Kausalitetens Princip. Ligesom \textsc{Spinoza} identificerer Begrundelse og Kausalitet,
saaledes identificerer han Aarsagssætningen og Sætningen om Bevægelsens Bestaaen
(den ufuldkomne Forløber for Energiens Bestaaen). Ogsaa nyere Filosofer har foretaget
denne sidste Identificering. Men ligesom der er et Spring fra Begrundelse til
Kausalitet, saaledes er der ogsaa et Spring fra den Sætning, at alt har sin Aarsag,
til den, at Aarsag og Virkning er ækvivalente, og der er igen et Spring fra den almindelige
Sætning om Ækvivalens til de specielle, faktiske Omsætninger, der i ethvert
enkelt Tilfælde foregaar i en bestemt Retning. Her maa den kritiske Filosofi
sondre, hvad \textsc{Spinoza} identificerer\footnote[1]{Smlgn. om disse Forhold \textit{Den menneskelige Tanke}, p. 211---215; 281---287.}.

At de Mellemled, de uendelige Modi betegner, ikke løser Vanskeligheden, ligger
deri, at de endelige Modi ligesaalidt kan udledes af dem, som de direkte kan udledes
af Substansen. Det er en Forudsætning hos \textsc{Spinoza}, at der nu engang gives
endelige Emner. Derfor er, som vi senere vil se, Erfaringen det første Trin paa
% --- p. 28 ---
\setcounter{footnote}{0}%
\opage{28}Erkendelsens Skala, og derfor er det kun Erfaring, der belærer os om, at Modi
overhovedet existerer (Ep. 10). \textsc{Spinoza}'s Filosofi er et Forsøg paa ved Analyse af
Erfaringsdata at stige op til Attributer og Substans. De uendelige Modi er Trin paa
denne Stige. \textsc{Spinoza} trækker, uden at være sig det klart bevidst, Stigen op efter
sig, saa at en deduktiv Nedstigen ikke bliver mulig.

Den Modsætning mellem et elementært og et idealt Aarsagsforhold, som stedse
gør sig gældende i det videnskabelige Arbejde, har hos \textsc{Spinoza} krystalliseret sig til
en Modsætning mellem en Verden af kvalitativt forskellige Emner (Modi), der opstaar
og afløser hverandre i Tiden, og et System af evige Tanker, indenfor hvilket raader
Kontinuitet (indenfor hvert enkelt Attribut) og tilsidst Identitet (i Substansen). Hin
første Verden er en for Sansning, Erindring og Fantasi fremtrædende Aabenbaring
af det andet System. Hvorfor og hvorledes der foruden Substansen existerer Attributer,
foruden Attributernes evige Modi, og foruden disse evige Modi kvalitetsforskellige
og tidsbestemte Modi, --- eller hvorfor der foruden den rene Tænken (intellectus)
ogsaa existerer sanselig Forestilling (imaginatio), --- det er Spørgsmaal,
\textsc{Spinoza} ikke opkaster, og som iøvrigt hverken han eller nogen anden vilde kunne
besvare. Al Videnskab, ogsaa Filosofien, arbejder med visse Forudsætninger, der
fremskridende kan staa deres Prøve i Erfaringens og Kritikens Skærsild, men
som ikke kan paavises at være de eneste mulige. \textsc{Spinoza}'s Efterfølger i Filosofiens
Historie, \textsc{Leibniz}, har udtrykt dette i den Sætning, at en Hypoteses Held ikke beviser
dens Sandhed.

\subsection*{d. Sætning 29---36.}
\begin{center}\textit{Nærmere Undersøgelse af Forholdet mellem Grundvæsen og Enkeltvæsener.}\end{center}

\textit{α.} Foruden Begreberne Substans og Gud, der (naar Substansen tænkes med
uendelig mange Attributer) er identiske, er der endnu et tredie Begreb, der bliver
identisk med dem begge, saa at \textsc{Spinoza} siger „Gud eller Naturen“ (deus seu natura,
f. Ex. Eth. IV, Fortale), ligesom han siger Gud eller Substansen. Dette Begreb egner
sig til at antyde en Forbindelse mellem de to Verdener, der, som vi nys har set,
truer med at træde i skarp Modsætning til hinanden.

Efterat det i 29 er sagt, at der i Tingenes Natur (in rerum natura) ikke gives
noget tilfældigt, men at Alt bestemmes ifølge den guddommelige Naturs Nødvendighed
(ex necessitate divinæ naturæ) til at existere og virke paa en bestemt Maade,
hedder det i Anmærkningen (29 Schol.): „Før jeg gaar videre, vil jeg forklare, eller
rettere minde om, hvad jeg forstaar ved den frembringende Natur (natura naturans),
og hvad jeg forstaar ved den frembragte Natur (natura naturata). Jeg mener, at det
af det foregaaende er tilstrækkelig klart, at man ved den frembringende Natur maa
forstaa det, som er i sig selv og forstaas ved sig selv, eller, hvad der er det samme,
saadanne Egenskaber (attributa) ved Substansen, som udtrykker Evighed og Uendelighed,
det vil sige, Gud, forsaavidt han betragtes som fri Aarsag. Ved den frembragte
Natur forstaar jeg derimod Alt, som følger af Guds Naturs eller (hvad der er
% --- p. 29 ---
\setcounter{footnote}{0}%
\opage{29}det samme) af et guddommeligt Attributs Nødvendighed, det vil sige, alle ved de
guddommelige Attributer bestemte Enkeltvæsener (dei attributorum modi), forsaavidt
som de betragtes som Ting, der er i Gud og uden Gud hverken kan være eller
tænkes“. --- Med andre Ord, til natura naturans hører Attributerne, til natura naturata
Modi.

Som vi tidligere har set, har \textsc{Spinoza} opstillet Begrebet uendelige Modi og udtalt,
at saadanne enten kan følge umiddelbart af Attributerne, eller middelbart, nemlig
gennem de umiddelbart opstaaede uendelige Modi (21---23; 28 Schol.). Begge
disse Arter af uendelige Modi hører, ligesom de endelige Modi, til natura naturata.
\textsc{Spinoza}'s skarpsindigste Korrespondent, \textsc{Tschirnhausen}, ønskede ogsaa her en nærmere
Forklaring og udbad sig Exempler paa saadanne uendelige Modi, der umiddelbart,
og saadanne, der middelbart udsprang af Guds Natur. I sit Svar (Ep. 64)
nævner \textsc{Spinoza} som Exempler paa den første Art: indenfor Tænkningens Attribut
den uendelige Forstand, og indenfor Udstrækningens Attribut Bevægelse og Hvile
(hvorved vistnok, som allerede ovenfor bemærket, menes Forholdet mellem Bevægelse
og Hvile --- ratio motus ad qvietem, som det hedder i Ep. 32, altsaa Bevægelsens
Bestaaen). Som Exempel paa de middelbart frembragte uendelige Modi nævner
han kun „hele Universets Beskaffenhed“ (facies totius universi, „Altets Fysiognomi“),
der trods alle Forandringer i det enkelte dog vedbliver at være en og den samme,
og han henviser til anden Bog, Anmærkningen til den syvende Laanesætning (Eth.
II. Lemma 7. Schol.), hvor det hedder, at „hele Naturen er ét Individ, hvis Dele,
det vil sige alle Legemer, varierer paa uendelig mange Maader uden nogen Forandring
i det hele Individ“. En lignende Udtalelse findes i Ep. 32, hvor det hedder,
„at ethvert Legeme omgives af andre Legemer, og at alle Legemer indbyrdes bestemmes
til at være og virke paa en vis bestemt Maade, dog saaledes, at der i hele
Universet bevares det samme Forhold mellem Hvile og Bevægelse“. Længst nede i
Rækken kommer saa de endelige Modi, de sjælelige og legemlige Fænomener, der
er Genstand for Iagttagelse.

\begin{quote}
Vi faar altsaa følgende Skema for de vigtigste Begreber hos \textsc{Spinoza}:\\
\textit{Substantia.}
\end{quote}

Natura naturans.

\begin{quote}
Cogitatio. Extensio.\\
Intellectus infinitus. Motus et qvies.\\
(immediate producta).
\end{quote}

Modi infiniti.

\begin{quote}
Facies totius universi. Natura naturata.\\
(mediate producta).\\
Modi finiti. \{ Res singulares (mentes et corpora).
\end{quote}

Vi ser altsaa, at \textsc{Spinoza} som svarende til „Bevægelsens Bestaaen“ i den materielle
Natur antager en aandelig Bestaaen, en uendelig Tanke, der forbliver den
samme under de specielle aandelige Fænomeners Skiften. Af Antydninger, der gives
% --- p. 30 ---
\setcounter{footnote}{0}%
\opage{30}i femte Bog, ser man, at han tænker sig intellectus infinitus som indeholdende det
evige, der i de enkelte aandelige Væsener fremtræder i tidsbegrænset Form og som
Korrelat til en begrænset materiel Existens. Med Rette har derimod flere Spinozaforskere
savnet et Begreb, der indenfor Tænkningens Attribut og i Underordning
under intellectus infinitus kunde svare til „Universets Fysiognomi“, som under Udstrækningens
Attribut underordnes „Bevægelsens Bestaaen“. Man har paa forskellig
Maade søgt at konstruere dette manglende Led i Systemet. Dette Punkt er dog neppe
af stor Betydning, skønt det viser, hvorledes et System først efterhaanden og under
Indflydelse af Spørgsmaal og Tvivl naar sin fulde Artikulation. \textsc{Spinoza}'s eget Svar
paa det ovennævnte Spørgsmaal af \textsc{Tschirnhausen} er meget kort og synes noget
flygtig affattet; han giver sig end ikke Tid til at kritisere en Korrespondents Opfattelse,
ifølge hvilken (Ep. 63) de uendelige Modi, der umiddelbart fremgaar af
Substansen, skulde være Attributerne! Denne Korrespondent (\textsc{Schuller}) har ikke
gjort sig fortrolig med Attributets Begreb som det, der udtrykker selve Substansens
Væsen. Og det siges udtrykkelig (23 Dem.), at en uendelig Modus er en umiddelbar
eller middelbar Følge af et Attributs guddommelige Natur.

Nærmest vilde det ligge at tænke paa en Art Verdenssjæl som Korrelat til „Verdensfysiognomiet“
eller „Verdensindividet“. Et saadant Begreb opstiller \textsc{Bruno}\footnote[1]{\textit{Den nyere Filosofis Historie\textsuperscript{2}}, I, p. 127---129.}
som Udtryk for den Kraft, der gør Vexelvirkning i Rummet mulig og tillige rører
sig inderst inde i de enkelte Elementer. Naar \textsc{Bruno} siger om „Verdenssjælen“:
amina tota in toto et qvalibet totius parte! (hele Sjælen i Helheden og i enhver Del % PRINTED AS IS: amina (anima)
af Helheden!) --- saa mindes man derved om det oftere hos \textsc{Spinoza} forekommende
Udtryk (Ethica II, 37---38; 44): „hvad der forekommer paa samme Maade i Delen
som i Helheden“. Ved dette Udtryk tænker \textsc{Spinoza} paa de almindelige Love, medens
\textsc{Bruno} tænker paa en Kraft, der gennemfører disse Love i det hele og i det
enkelte. Dette Begreb var for mystisk og poetisk for \textsc{Spinoza}, og han har ikke
fundet noget, der kunde indtage dets Plads.

For \textsc{Spinoza} var den lovmæssige Naturorden, der betinges ved „Bevægelsens
Bestaaen“ og de tilsvarende Forhold indenfor andre Attributer, ikke et Abstraktum
(et universale), men en karakteristisk Side ved Tilværelsen, en individuel Ejendommelighed
ved det Unicum af en Verden, vi kender. Modi infiniti er baade individuelle
og universale. De falder sandsynligvis sammen med, hvad \textsc{Spinoza} i Afhandlingen
om „Forstandens Forbedring“ kalder „de faste og evige Ting, i hvilke
Naturlovene er indskrevne som i Lovbøger“. Om dem siges der, at „uagtet de er
individuelle (singularia), kan de dog for os staa som universelle (tanqvam universalia)
paa Grund af deres Allestedsnærværelse og omfattende Virken; de er som almene
Definitioner af de foranderlige Ting og de nærmeste Grunde til alt“. Altsaa
svarer disse faste og evige Ting til, hvad der i Ethica kaldes Modi infiniti af
anden Klasse (mediate producta). I hvert Tilfælde vidner dette Begreb om, at \textsc{Spinoza}
i den Naturorden, Videnskaben afslører, har set et stort konkret Faktum, et
Unicum. De almene Love tjener kun til Karakteristik af dette Unicum. Og som
% --- p. 31 ---
\setcounter{footnote}{0}%
\opage{31}det senere vil vise sig, er Erkendelsen af almene Love for \textsc{Spinoza} ikke den højeste
Art af Erkendelse\footnote[1]{Denne Opfattelse af „de faste og evige Ting“ er, forekommer det mig, klart antydet i \textit{Den nyere Filosofis Historie\textsuperscript{2}}, I, p. 302, hvorfor \textsc{Huan} (p. 251) med Urette nævner mig blandt dem, der fortolker Udtrykket for som Betegnelse for almene Love. I Modsætning til \textsc{Huan} mener jeg afgjort, at res fixæ et æternæ er modi infiniti \textit{mediate} producti. --- Oplysende vil maaske her Paavisningen af typiske Individualforestillingers erkendelsesteoretiske Betydning være: \textit{Totalitet som Kategori} § 12---13.}.

Det Spørgsmaal, hvorledes vi fra de endelige Modi stiger op til „de faste og
evige Ting“, og i det hele til de uendelige Modi, har \textsc{Spinoza} ikke nærmere undersøgt.
Han bebuder en Undersøgelse af „de moderne Empirikeres Fremgangsmaade“, men
hans erkendelsesteoretiske Skrift blev ikke fuldført, og derfor blev dette Spørgsmaal
ikke behandlet af ham.

\textit{β.} Naar \textsc{Spinoza} i Definitionen af natura naturans kalder Gud en fri Aarsag
(29 Schol.), maa den fra først af (Def. 7) givne Definition af Frihed som Bestemthed
ved det egne Væsens Natur fastholdes. I Sætning 33, Schol. 2 gaar \textsc{Spinoza} nærmere
ind paa dette Spørgsmaal for at afvise den hyppige Opfattelse af Frihed som
Indifferens, Mulighed i lige Grad for to Modsætninger. Frihed i denne sidste Betydning
vilde, mener han, stride mod Guds Fuldkommenhed. Man tænker sig gerne
først Gud som existerende og derpaa som tagende sine Beslutninger. Men for et
evigt Væsen gives der intet „før“ eller „efter“. Guds Beslutning er ét med hans
evige Væsen og kan derfor hverken begynde eller ændres. Der gives ingen fra Virkelighed
forskellig Mulighed i Gud. Dog finder \textsc{Spinoza} Antagelsen af en absolut
Indifferens hos Gud mindre urimelig end den Antagelse, at Gud skulde handle for
det Godes Skyld eller med det Gode for Øje (sub ratione boni): dette vilde være at
lade Gud stræbe efter et Forbillede, et Maal, som han endnu ikke har naaet, --- ja,
det vilde egentlig gøre ham afhængig af en Skæbne (fatum), --- ham, der ifølge sit
Væsens Natur er den første Aarsag til alt!

\textit{γ.} Den sidste Sætning i første Bog (36) hævder et ofte miskendt Element i
\textsc{Spinoza}'s System. Den lyder: „Der existerer intet, af hvis Natur der ikke følger
en Virkning“, og den begrundes ved, at alt, hvad der existerer, udtrykker Guds
Natur eller Væsen paa en vis bestemt Maade (smlgn. 25---26).

Denne Sætning er af særlig Betydning for Forstaaelsen af den Rolle, Selvopholdelsestrangen
hos endelige Væsener spiller som Grundlag for \textsc{Spinoza}'s Psykologi
og Etik (i tredie og fjerde Bog, se dog allerede II, 45 Schol.). Det er som Element
i Substansen, at den enkelte Modus kan siges at være Aarsag til sig selv, have sin
Grund i sig selv, skønt denne Selvstændighed brydes med andre Modi's Indvirken, der
hver for sig ogsaa kan siges at have deres Grund i sig selv. Ethvert Led i Aarsagsrækken
er baade Aarsag og Virkning, skønt man er tilbøjelig til at holde sig til de
følgende Leds Afhængighed af de foregaaende. Brydningen af Aktivitet og Passivitet
vil vise sig at spille en stor Rolle i \textsc{Spinoza}'s Psykologi og Etik.

% --- p. 32 ---
\setcounter{footnote}{0}%
\opage{32}\subsection*{e. Tillæg til første Bog.}
\begin{center}\textit{Teologi og Videnskab.}\end{center}

\textit{α.} Naar \textsc{Spinoza} spurgte sig selv, hvad der hindrede Anerkendelsen af den
store Sammenhæng mellem Tingene i Verden (rerum concatenatio) i at blive almindelig,
fandt han Forklaringen i, at Menneskene i deres Betragtning af Naturen ledes
af en praktisk Trang, ikke af Trang til Videnskab. De er derfor tilbøjelige til at
tro, at alt er indrettet til deres eget Bedste, det, som er Genstand for deres egen
stadige Stræben. Fra først af er de uvidende om Tingenes Aarsager, ogsaa om
Aarsagen til den Trang, de har til at søge deres eget Bedste. De tror sig frie,
naar de ledes af denne Trang, og de tror, at Midler til at tilfredsstille denne Trang,
forsaavidt de ikke selv kan frembringe disse, frembringes af Væsener, der ligner
dem selv. Naar der overgaar dem Ulykker, udleder de det af Gudernes Vrede over
Menneskenes Synder. Og hvis der endnu staar uløste Gaader tilbage, trøster de sig
med, at Gudernes Visdom langt overgaar Menneskenes Forstand. Overfor hele dette
kunstige System (fabrica) vilde Sandhedens Sag have været haabløs, dersom ikke
Matematiken, der ikke spørger om Formaal, men blot om Tingenes Væsen og Egenskaber,
havde givet Menneskene en anden Sandhedsnorm end den, deres blinde
Trang fører dem til at følge.

Den teleologisk-teologiske Opfattelse vender op og ned paa Sagerne. Den gør
det første til det sidste, Aarsag til Virkning, det fuldkomne til det mest ufuldkomne.
Den evige Naturorden er ikke Virkning af Menneskers Trang og Ønsker,
men indeholder Forklaringen af, at Mennesker overhovedet kan føle Trang og nære
Ønsker. Og det fuldkomne er ét med den evige Naturorden, der ikke hviler paa
nogen Forudsætning; det ufuldkomne er det, der forudsætter mange Betingelser.
Den teologiske Opfattelse gør Gud ufuldkommen ved at lade ham stræbe efter
noget, som endnu fattes ham. Og den indfører en hel ny Maade at tænke paa:
istedetfor at bevise noget ved at godtgøre Umuligheden af det modsatte (reductio
ad impossibile) appellerer den til Uvidenheden (reductio ad ignorantiam). Og Guds
ubegribelige Villie bliver Uvidenhedens sidste Tilflugt (asylum ignorantiæ). Den, der
søger efter naturlige Aarsager til det, der vækker Forundring og Betagethed, betragtes
som en ugudelig Kætter af de Personer, i hvilke den store Masse ser Fortolkere af
Naturens Væsen og af Gudernes Villie. Og disse Personer ser jo nok, at hvis Uvidenheden
hævedes, vilde ogsaa den tankeløse Betagethed (stupor) forsvinde, uden
hvilken deres Autoritet ikke kan holdes oppe\footnote[1]{I Ep. 73 skelner \textsc{Spinoza} mellem Religion og Overtro saaledes, at hin grunder sig paa Visdom, denne paa Uvidenhed.}.

Man kan, fortsætter \textsc{Spinoza}, overhovedet ikke anvende Begreber som godt og
ondt, Orden og Forvirring, skønt og hæsligt til Forklaring af Naturen, da saadanne
Begreber skyldes menneskelig Trang, ikke stammer fra nogen i Naturen paaviselig
Sammenhæng. Og naar man klager over de mange Ufuldkommenheder i Naturen,
glemmer man, at det ikke er menneskelig Gavn og Skade, der kan afgive Maalestok
% --- p. 33 ---
\setcounter{footnote}{0}%
\opage{33}for Fuldkommenhed og Ufuldkommenhed. Maalestokken maa søges i selve Tingenes
Natur og Magt (ex sola rerum natura et potentia)\footnote[1]{Af Identiteten af Realitet og Fuldkommenhed følger, at hvad der slet ingen Fuldkommenhed har, ikke kan existere. Derfor kan Djævelen ikke existere! (\textit{Den korte Traktat}, II, 25).}. Spørger man da, hvorfor Gud
ikke har skabt Menneskene saaledes, at de ene lededes af Fornuften, maa Svaret
blive, at Naturens Love er saa omfattende, at alt, hvad en uendelig Forstand kunde
tænke, fra den højeste Grad af Fuldkommenhed til den laveste, kunde finde sin Plads.

Hele denne Tankegang er i høj Grad karakteristisk for \textsc{Spinoza}. Han staar
her som den, der hævder den strenge videnskabelige Metodes Ret, og tillige som
den, der er overbevist om, at der ved fuld Gennemførelse af denne Metode ikke
mistes noget af det, der virkelig har Værdi. Menneskene maa kun ikke maale Tilværelsens
Værdi efter deres egen lille Verdens Behov. Indenfor den menneskelige
Verden har de Værdibegreber, hvis kosmiske Gyldighed han afviser, deres Berettigelse,
saaledes som han forsøger at vise det i fjerde Bog.

Det er, som allerede omtalt, alligevel et Værdibegreb, \textsc{Spinoza} anvender, naar
han identificerer Fuldkommenhed og Realitet. Der indtræder derfor, naar han i
fjerde Bog gaar over til en specielt menneskelig Værdilære, ikke blot en Brydning
mellem Kosmologi og Etik, men ogsaa en Brydning mellem to Værdimaalere. Dertil
vil vi paa sit Sted komme tilbage. Her skal blot bemærkes, at medens \textsc{Spinoza}
fra det kosmologiske Synspunkt bestemmer det Fuldkomne som det, der ingen Forudsætninger
har, maa det fra det specielt menneskelige Synspunkt (som han selv
anlægger i fjerde Bog) omvendt siges, at det Fuldkomne netop er det, der hviler
paa de fleste Forudsætninger. Livet har flere Forudsætninger end det Livløse,
Sjælelivet end det blot organiske Liv, det etiske, intellektuelle og æstetiske Liv end
det elementære Sjæleliv. Kategorilæren viser, at Værdibegreberne er de mest specielle
af alle vore Grundbegreber\footnote[2]{Se herom \textit{Den menneskelige Tanke}, p. 237---241.}.

\textit{β.} I den Strid om \textsc{Spinoza}'s Filosofi, der fremkaldtes ved \textsc{Jacobi}'s „Briefe
über die Lehre des \textsc{Spinoza}“ (1785), især ved de mundtlige Ytringer af \textsc{Lessing},
\textsc{Jacobi} anførte, kom særlig det Spørgsmaal frem, om alt kunde udledes af Tanken
om et højeste Formaal. I en Samtale med \textsc{Jacobi} havde \textsc{Lessing} kaldt sig Spinozist
og erklæret det for en Fordom, at Tanken skulde være det første og højeste; Udstrækning
og Bevægelse er ligesaavel som Tanke grundet i en højere Kraft, der ikke
er udtømt med disse Former. Han havde udtalt sig imod, at „unsere elende Art,
nach Absichten zu handeln“ skulde gælde for den bedste Metode. Foruden af \textsc{Spinoza}
kan \textsc{Lessing} her være paavirket af \textsc{Hume}, hvis „Dialogues on Natural Religion“
var udkomne kort iforvejen. Samtalen mellem \textsc{Lessing} og \textsc{Jacobi} fandt Sted i
Sommeren 1780. Den vakte, da den blev bekendt, stor Opmærksomhed og Forargelse,
idet Oplysningstidens populære Teologi (som man havde troet, ogsaa \textsc{Lessing}
hyldede) omstyrtedes, naar man ikke kunde slutte til Guds Tilværelse fra Naturens
Hensigtsmæssighed\footnote[3]{Om \textsc{Lessing}'s Stilling i Filosofiens Historie se \textit{Den nyere Filosofis Historie\textsuperscript{2}}, II, p. 16---25).}.

% --- p. 34 ---
\setcounter{footnote}{0}%
\opage{34}Ifølge \textsc{Jacobi} var \textsc{Spinoza} Ateist, fordi hans Filosofi nægtede Guds Personlighed.
\textsc{Spinoza} selv havde afvist denne Benævnelse, som han mente skyldtes den
Misforstaaelse, at han gjorde Gud afhængig af Skæbnen. Tværtimod, siger han (Ep.
43), antager jeg, at alt følger med uundgaaelig Nødvendighed af Guds Natur. For
\textsc{Spinoza} udelukker jo Frihed og Nødvendighed ikke hinanden. Og han henviser
desuden til, at efter hans Opfattelse er de etiske Idealer uafhængige af, om vi udleder
dem af Gud som Dommer eller af den guddommelige Naturs evige Nødvendighed.


Det er ikke let at afgøre, hvorledes man religionsfilosofisk skal karakterisere
\textsc{Spinoza}'s Standpunkt. I Almindelighed mener man, at enhver anstændig Filosofi
skal repræsentere en eller anden -isme, og man mener saa tillige, at man en Gang
for alle er i Besiddelse af en Katalog over alle mulige -ismer, i hvilken enhver Filosofi
maa kunne finde sin Plads. --- Hvad det for \textsc{Spinoza} karakteristiske angaar, er
især to Punkter af Betydning. Han hylder ikke den sædvanlige teologiske Forestilling
om Gud og Verden som to forskellige Væsener, men Gud er for ham Verdens
inderste Væsen, der kundgør sig i dens Love. Og han nægter, at Personlighedsbegrebet
kan finde Anvendelse paa Gud, da Personlighed (Forstand og Villie), som vi
ovenfor har set, efter hans Opfattelse forudsætter noget givet, der skal forstaas, og
en Modstand, der skal overvindes for et Formaals Skyld. I Begrebet Panteisme,
som \textsc{Toland} indførte i Begyndelsen af det attende Aarhundrede, og som gerne opfattes
saaledes, at det udtrykker et indre Forhold mellem Gud og Verden og nægter
Guds Personlighed, har man ofte ment at have en Rubrik, hvori \textsc{Spinoza} kunde
passe ind. Men ifølge \textsc{Spinoza} er nu dog Tanken (cogitatio), skønt alt ikke kan
udledes af den, et evigt Attribut ved Substansen, uagtet de specielle Former, vi
kalder Forstand og Villie, ikke hører til „den frembringende Natur“. I dette Attributs
Begreb ligger en lignende Idealisation som den, filosofiske Teister (som f. Ex.
\textsc{Lotze}) foretager, naar de sublimerer Begrebet Personlighed saaledes, at det formenes
at passe paa et Væsen, der skal have sin Grund i sig selv. Grænsen mellem Panteismen
og Teismen er altsaa ikke let at drage\footnote[1]{Se nærmere om Begreberne Panteisme og Teisme min \textit{Religionsfilosofi} § 22 og § 95. --- Paa den der anlagte rent psykologiske Betragtningsmaade af, hvad der ligger til Grund for disse Begreber, er jeg, siden jeg skrev min \textit{Religionsfilosofi}, kommen til at lægge stedse større Vægt.}.

Jeg har ved denne Betragtning, som allerede er gjort gældende i „Den nyere
Filosofis Histori“ (2. Udg. I, p. 316), kun villet vise, hvor vanskeligt det er at bringe % PRINTED AS IS: Histori (Historie)
et dybtgaaende Tankearbejde som \textsc{Spinoza}'s ind under Rubriker, der er fastslaaet
iforvejen. Med Urette nævner derfor \textsc{Huan} mig blandt dem, der fortolker \textsc{Spinoza}'s
Gudsbegreb i teistisk Retning\footnote[2]{\textit{Le dieu de Spinoza}, p. 221.}. \textsc{Spinoza}'s Tænkning har paa det her omtalte
Punkt, som paa mange andre Punkter, en saadan individuel Ejendommelighed, at
den ikke er let at klassificere, men, som \textsc{Huan} træffende siger paa sidste Blad i sin
Bog, egentlig kun kan betegnes som --- Spinozisme!

Grænsen mellem Teologi og Filosofi finder \textsc{Spinoza} deri, at hin bruger psyko
% --- p. 35 ---
\setcounter{footnote}{0}%
\opage{35}logiske Udtryk om Gud, medens denne nægter saadanne Udtryk Berettigelse. Naar
han advarer Filosoferne mod at bruge teologiske Udtryk (Ep. 23), er det dette, han
tænker paa. Men der ligger bag hans egen Fremstilling, især i Slutningsanmærkningen
til første Bog, et endnu mere fundamentalt Problem, nemlig om Berettigelsen
af overhovedet at lægge Værdibegreber til Grund ved Afslutningen af vor Verdensopfattelse.
\textsc{Spinoza} selv har af det teologiske Gudsbegreb beholdt mere, end han
er sig bevidst. I hans Substans skulde efter Definitionen kun være indeholdt Forstaaelighedens
(Rationalitetens) Princip; men den indeslutter --- formedelst Definition
af Fuldkommenhed som Realitet --- tillige et Værdibegreb i sig\footnote[1]{Smgln. min \textit{Religionsfilosofi} § 24 og § 80---89.}. Den Begrundelse,
\textsc{Spinoza} giver heraf (i femte Bog), er ganske vist kun den, at det er Substansen,
der er Aarsag til Erkendelsesglæden. Men saa opdukker Spørgsmaalet om Forholdet
mellem intellektuel Værdi og andre Værdier.

% --- p. 36 ---
\setcounter{footnote}{0}%
\opage{36}\section*{Anden Bog.}
\subsection*{a. Definitioner og Axiomer.}

Som Indledning til anden Bog giver \textsc{Spinoza} nogle Begyndelsesbestemmelser
og Grundsætninger, der er af erkendelsesteoretisk og psykologisk Interesse. Anden
Bogs Indhold kan siges at være Erkendelsens Psykologi og Undersøgelse af Erkendelsens
Gyldighed (Erkendelsesteori).

\textit{α.} Af særlig Vigtighed er Def. 3, ifølge hvilken en Forestilling (idea) er et
Begreb, som Sjælen (mens) danner, fordi den er et tænkende Væsen. Og denne
Definition forklares nærmere, idet Vægten lægges paa, at enhver Forestilling er et
Produkt af aandelig Virken (actio mentis); derfor passer Ordet Begreb (conceptus)
bedre end Ordet Fornemmelse (perceptio), der kunde synes at antyde, at Bevidstheden
forholder sig passiv overfor Genstanden. --- \textsc{Spinoza} berigtiger her sin egen
tidligere Opfattelse, idet han i den korte Traktat lod Bevidstheden staa passiv overfor
sine Genstande og lod det ene Attribut paavirke det andet.

Def. 4 giver en videre Udvikling af disse Bemærkninger, idet den bestemmer
en fuldkommen Tanke (idea adæqvata) som en saadan, der, betragtet i og for sig
og afset fra Genstanden, har alle en sand Tankes indre Egenskaber. De ydre Egenskaber
angaar Overensstemmelsen med Genstanden (ideatum). Hvilke de indre
Egenskaber er, kan ses af Afhandlingen om Forstandens Forbedring, hvor det siges,
at en sand Tankes Form beror paa dens egen Kraft og Natur og viser sig i Sammenhængen
mellem Subjekt og Prædikat, Grund og Følge. Denne indre Sammenhæng
kan føre os ud af Falskhed og Tvivl og er den egentlige Sandhedsnorm. Der
kan intet ydre Tegn paa Sandhed gives; de ydre Egenskaber ved Tanken, dens
Overensstemmelse med Genstanden, er derfor afledede i Forhold til de indre ---
saaledes som det allerede ses af Eth. I, Ax. 6, naar dette Axiom forstaas rigtig. (Se
ovenfor).

Da \textsc{Spinoza} blev spurgt af \textsc{Tschirnhausen}, hvilken Forskel der var mellem en
fuldkommen Tanke (idea adæqvata) og en sand Tanke (idea vera), svarede han (Ep.
60), at der i Grunden ingen Forskel var, kun at Ordet „sand“ hentyder til Tankens
Overensstemmelse med sin Genstand (ideatum), medens Ordet „fuldkommen“ hentyder
til Tankens Natur i og for sig selv.

Fra Tankens Genstand (objectum = ideatum = res percepta) skelner \textsc{Spinoza}
(som det ses af Sætningerne 5, 7, 13 og 15) Tankens Indhold (esse formale ideæ,
% --- p. 37 ---
\setcounter{footnote}{0}%
\opage{37}smlgn. I, 30 Dem.: id, quod in intellectu continetur). Genstandene for Begreberne
Trekant (I, 11 Dem.) eller Cirkel (II, 7 Schol.) er de virkelig (in natura) existerende
Trekanter eller Cirkler; Genstanden for den Tanke, hvori Sjælen bestaar, er det til
Sjælen svarende Legeme. Men disse Begrebers Indhold er givet med deres videnskabelige
Definitioner. Genstandene (Trekanter, Cirkler, Legemer) hører under Udstrækningens
Attribut, men Indholdet af Begreberne hører til Tænkningens Attribut.

Det paafaldende ved denne Opstilling er for det første, at selve Tankevirksomheden,
som \textsc{Spinoza} i Def. 3 og 4, saavel som senere hen (43 Schol.; 49 Schol.),
saa stærkt betoner, ikke nævnes her. Den hører aabenbart ind under Tænkningens
Attribut, ligesaa vel som det Begrebsindhold, den fører til at slaa fast, og der kunde
ogsaa for dens Vedkommende være Grund til at spørge. hvad der svarer til den % PRINTED AS IS: a stop for a comma
under Udstrækningens Attribut. Det vil vise sig, at dette Punkt hænger sammen
med en Uklarhed i \textsc{Spinoza}'s Lære om Forholdet mellem Attributerne. --- For det
andet er det paafaldende, at Triangler og Cirkler staar paa Linie med menneskelige
Legemer og altsaa betragtes som reale Naturgenstande. Der tales saaledes (II, 7
Schol.) om den i Naturen existerende Cirkel. Dette Punkt hænger sammen med
\textsc{Spinoza}'s dogmatiske Erkendelseslære, ifølge hvilken der til hvert videnskabelig
dannet Begreb svarer en Realitet.

\textit{β.} Det Begreb om aandelig Virken (actio mentis), som Def. 3 opstiller, forudsætter
en Definition af Aktivitet, som \textsc{Spinoza} mærkelig nok først giver i tredie
Bogs anden Definition, men som burde have haft sin Plads her. Det er ikke den
eneste Anledning, der er til at give \textsc{Leibniz} Ret, naar han i sit Haandexemplar af
Ethica skrev, at Bevisførelsen ikke var \textsc{Spinoza}'s stærke Side (\textsc{Spinoza} non est magnus
demonstrator).

Definitionen af Aktivitet (III, Def. 2) lyder saaledes: „Jeg siger, at vi er aktive
(agere), naar der sker noget i os eller udenfor os, som vi er fuldstændig Aarsag
til (adæqvata causa), d. v. s. naar der af vor Natur følger noget i os eller udenfor
os, som klart og tydelig kan forstaas ved den alene. Derimod siger jeg, at vi er
passive (pati), naar der sker noget i os eller følger noget af vor Natur, som vi kun
er delvis Aarsag (causa partialis) til“.

I denne Definition ligger, at helt passive er vi aldrig, hvorimod vi meget vel
skulde kunne være absolut aktive, altsaa fuldstændig Aarsag til Noget, som sker i
os eller udenfor os. Nu er det at være fuldstændig Aarsag ifølge den syvende Definition
i første Bog ét med at være fri. Fuldstændig Aktivitet eller Frihed bliver
for \textsc{Spinoza} et Ideal (exemplar), som etisk Stræben er rettet imod (IV, Fortale), og
han beskriver i en Række af Sætninger (IV, 67---73) det frie Menneske som Højdepunktet
af menneskelig Udvikling. Men erfaringsmæssig er, som vi vil se \textsc{Spinoza}
selv indrømme, ethvert Enkeltvæsen (Modus) baade aktivt og passivt. Kun Substansen
er absolut aktiv.

Vi kan dog stadig bruge \textsc{Spinoza}'s Definition af Aktivitet, naar vi ændrer den
saaledes: Vi er des mere aktive, jo mere Aarsagen til vore Tilstande og Handlinger
ligger i os selv, --- des mere passive, jo mere Aarsagen til dem ligger udenfor os.

% --- p. 38 ---
\setcounter{footnote}{0}%
\opage{38}\textit{γ.} For \textsc{Spinoza}'s Psykologi er Ax. 3 af Interesse. Det udsiger, at saadanne
sjælelige Fænomener (modi cogitandi) som Kærlighed, Drift (cupiditas) o. s. v. forudsætter
Forestillinger om det, der elskes, attraas o. s. v., men at der kan gives
Forestillinger, uden at andre sjælelige Fænomener er tilstede.

\textsc{Spinoza} har naturligvis Ret i, at vi ikke kan elske eller attraa uden Forestilling
om, hvad vi elsker eller attraar. Men han overser, at der kan gives elementære
Følelser som Angst, Velbefindende, Nedstemthed og desuden Trang og Stræben,
uden at der er nogen Forestilling tilstede om Genstanden for (Aarsagen til) disse
Tilstande. Det erkender han selv i Anmærkningen til niende Sætning i tredie Bog,
hvor han skelner mellem den blinde Selvopholdelsestrang (appetitus) og Driften (cupiditas),
hvilken sidste forudsætter en Forestilling om en Genstand. Ja, han føjer
endog til, at vi ikke trænger til noget eller vil det, fordi det er godt, men omvendt
anser det for godt, fordi vi trænger til det og vil det. Der fremtræder her en Brydning
mellem Intellektualisme og Voluntarisme, som er karakteristisk for \textsc{Spinoza}'s
Psykologi, og som vi senere vil komme tilbage til. Her skal kun tilføjes, at det er
i højeste Grad tvivlsomt, om \textsc{Spinoza} har Ret i, at der kan gives Forestillinger uden
nogensomhelst Grad af Følelse eller Trang.

Det femte Axiom, at vi ikke kender andre Emner end Legemer og Tankemodi,
har vi allerede drøftet ovenfor.

\subsection*{b. Sætning 1---10.}
\begin{center}\textit{Forholdet mellem Attributerne.}\end{center}

\textit{α.} For ret at opfatte \textsc{Spinoza}'s Lære om Forholdet mellem Attributerne maa
man fastholde, at der hos \textsc{Spinoza}, endnu mere end hos andre Filosofer før det
attende Aarhundrede, mangler den bestemte Sondring mellem forskellige filosofiske
Problemer, som nu efterhaanden trænger igennem. Religion, Verdensanskuelse (Kosmologi),
Erkendelsesteori, Etik, Psykologi --- alle de Bestræbelser og Opgaver, der
betegnes ved disse Ord, slynger sig saaledes i ét hos \textsc{Spinoza}, at det kan være
vanskeligt at udrede dem, og at et Forsøg derpaa let fører til at lægge for moderne
en Tankegang ind hos denne Renaissancetænker. Hvad Forholdet mellem Attributerne
angaar, er det især Forskellen mellem Erkendelsesteori og Psykologi (resp.
Psykofysik), vi faar med at gøre.

Hvad \textsc{Spinoza} i de sex første Sætninger af anden Bog fremsætter om Attributernes
indbyrdes Forhold, er en Konsekvens af den anden og fjerde Definition i
første Bog, en Konsekvens, der allerede er draget i første Bogs tiende Sætning. Naar
hvert Attribut skal udtrykke Substansens Væsen, og naar Substansens Væsen bestaar
i at have Grunden (eller Aarsagen) i sig selv, saa maa ethvert Attribut ogsaa indeholde
Grunden (eller Aarsagen) til alle Emner, der hører ind under det, og man kan
ikke hente nogen Forklaring ud fra andre Attributer.

Men det Spørgsmaal ligger da nær, hvad der har ført \textsc{Spinoza} til at opstille
en Definition af Begrebet Attribut, som medfører saadanne Konsekvenser. Om de
% --- p. 39 ---
\setcounter{footnote}{0}%
\opage{39}Tankemotiver, der har gjort sig gældende paa dette Punkt, er der allerede talt i
Indledningen (se ovenfor). Det ene Motiv er naturvidenskabeligt og skyldes den
Opfattelse af den nye Videnskabs Metode og Principer, som \textsc{Spinoza} har vundet
i Mellemtiden mellem Affattelsen af den korte Traktat og Udarbejdelsen af Ethica,
og som krævede, at ethvert materielt Emne skulde forklares ved et andet materielt
Emne, der ligesom hint kunde paavises i Erfaringen. Det andet Motiv er
erkendelsesteoretisk. Sandhedskriteriet viser sig tilsidst at bestaa i Tankernes indbyrdes
Overensstemmelse og kræver altsaa, at enhver Tanke skal udledes af en
anden Tanke.

Disse to Tankemotiver har hos \textsc{Spinoza} gensidig belyst hinanden. Den principielle
Kontinuitet, Naturvidenskaben fandt eller krævede paa det materielle Omraade,
og som allerede lod \textsc{Descartes} opfatte alt materielt som en eneste Substans,
søgte \textsc{Spinoza} naturlig et Sidestykke til paa det aandelige Omraade. Men af II, 1.
Schol. og 6. Cor. ses, at Tankens Omraade selv har frembudt for ham en lovmæssig
Sammenhæng, der har kunnet være Forbillede for Forskningen paa det materielle
Omraade. I 1. Schol. vises, hvorledes der kan bygges Tankerækker, der efter samme
Lov kan fortsættes i det Uendelige, og i Sætning 2 overføres saa dette ad Analogiens
Vej paa Udstrækningens Attribut. I 6. Cor. hedder det, at paa samme Maade og
med samme Nødvendighed, som Tankerne følger af Tænkningens Attribut, følger og
sluttes Objekterne (res ideatæ) af deres forskellige Attributer. Her sluttes altsaa fra
den indre Sammenhæng indenfor Tænkningens Attribut til en tilsvarende Sammenhæng
indenfor andre Attributer. Det eneste af disse andre Attributer, som Erfaringen
kundgør for os (Ax. 5), er Udstrækning. Overførelsen af Begrebet om Rækkedannelser
efter iboende Love fra Tankens Attribut til Udstrækningens Attribut betød
for \textsc{Spinoza}, der, heri følgende \textsc{Descartes}\footnote[1]{\textit{Principia Philosophiæ}, II, 10---11. (Rum og Legeme er kun forskellige som genus og species, hvilket tydelig viser sig derved, at selv om vi tænker os alle sanselige Egenskaber aftage i Grad, saa at de forsvinder for os, bliver dog Rummet tilbage).}, fandt Legemernes Væsen udtømt i de
rent geometriske Egenskaber (se Def. 1), at dette Begreb ogsaa maatte gælde for
alle Legemer. At alle Legemer er i Hvile eller i hurtigere eller langsommere Bevægelse,
hævder han senere hen i Bogen (i de to Axiomer efter Sætning 13) som en
Erfaringskendsgerning, hvis Forhold til Definitionen af Legemlighed som Udstrækning
han ikke undersøger; i sine sidste Aar ytrede han dog Tvivl om den kartesianske
Definitions Rigtighed (se Ep. 81 og 83). --- At der er en saadan tilsvarende
Sammenhæng indenfor ethvert Attribut, det er, efter den i det Foregaaende fremstillede
Tydning af \textsc{Spinoza}'s Grundbegreber, det, som \textsc{Spinoza} har betegnet ved, at
ethvert Attribut paa sin Vis udtrykker Substansens Væsen.

Endnu et tredie Motiv har sikkert været medvirkende, nemlig den bestemte
Afvisning af al Teleologi, som \textsc{Spinoza} udtaler, ikke blot i Tillægget til første Bog,
men flere Steder i Løbet af Bogen. Hans Kritik af Teologiens Gudsbegreb og hans
Opfattelse af videnskabelig Metode virkede her i samme Retning. Al Teleologi forudsætter,
at en Tanke eller Plan er Aarsag til Forandringer i Naturen; men det
% --- p. 40 ---
\setcounter{footnote}{0}%
\opage{40}vilde forudsætte, dels at der ingen indre lovmæssig Sammenhæng skulde være i
Naturen, dels at der i Guds Væsen skulde være en Forskel mellem Tanke og Magt.

Der har dog sikkert hos \textsc{Spinoza} selv ikke gjort sig nogen skarp Sondring
gældende mellem disse forskellige Tankemotiver, som vi her har udredet hvert for sig.

\textit{β.} I Sætning 7, som udsiger, at Tankernes Orden og Sammenhæng er den
samme som Tingenes Orden og Sammenhæng, udsiges Identitetshypotesen som gældende
for Tænkningens Attribut i Forhold til alle andre Attributer. Ved Tingene
maa nemlig her forstaas alle Emner indenfor alle disse andre Attributer. Hvad der
menes ved Tanker, ses af Beviset, der henviser til fjerde Axiom i første Bog: „Virkningens
Begreb (idea) afhænger af Aarsagens Begreb (idea) og forudsætter den“.
Det er klart, at de Tanker, der her er Tale om, er videnskabelig udformede Tanker
med deres klare og bestemte Indhold --- hvad \textsc{Spinoza} (se ovenfor p. 36) kalder
esse formale idearum\footnote[1]{Formalis betyder hos \textsc{Spinoza}, hvad vi kalder objektiv. Og ligesom vi ikke blot taler om objektive Genstande i Modsætning til blotte Tanker, men ogsaa om objektive Tanker i Modsætning til Ideassociationer, Fantasier og Følelser, saaledes taler \textsc{Spinoza} ogsaa baade om essentia formalis attributorum dei (Eth. II, 40 Schol. 2), altsaa om Genstande, og om esse formale idearum (Eth. II, 5), altsaa om Tankeindhold. --- Med Hensyn til Terminologien maa endnu mærkes, at „objective“ (7 Cor.) betyder, hvad vi vilde kalde „subjektivt“, idet det bruges som Modsætning til formaliter. Se især I, 30 Dem., hvor „objective“ staar i Modsætning til „in Natura“.}. Det er til dette udklarede Tankeindhold, at „Tingene“ svarer,
idet Tingenes indbyrdes Afhængighed svarer til Følgens (Virkningens) Afhængighed
af Grunden (Aarsagen). Man ser her klart, hvorledes Substansen finder Udtryk i
de forskellige Attributer: Forholdet mellem Grund og Følge forudsættes i dem alle,
er det i dem alle identiske (hvor mange de saa end er).

I Anmærkningen til den syvende Sætning anvendes Identitetshypotesen specielt
paa Forholdet mellem de to Attributer, vi kender af Erfaring. „Den tænkende og
den udstrakte Substans“, hedder det, „er en og samme Substans, der nu betragtes
under dette Attribut, nu under hint. Og saaledes er ogsaa en Udstrækningsmodus
og Tanken om denne Modus en og samme Ting, kun udtrykt paa to Maader\dots{}..
Saalænge vi betragter Tingene som Tankemodi, maa vi forklare hele Naturens Orden
eller Aarsagernes Sammenhæng ved Tænkningens Attribut alene, og saalænge vi
betragter Tingene som Udstrækningsmodi, maa hele Naturens Orden forklares ved
Udstrækningens Attribut alene. Og det samme gælder efter min Opfattelse om de
andre Attributer. Gud er Aarsag til Tingene, som de er i sig selv\footnote[2]{Udtrykket „Tingene, som de er i sig selv“, forekommer hos \textsc{Spinoza} i forskellig Sammenhæng og derfor i forskellig Betydning. 1) II, 7. Schol. betegner det Tingene, forsaavidt de betragtes under alle mulige Attributer (hvad menneskelig Tænkning ikke formaar). 2) II, 44 Dem. (med Henvisning til I, 29) betegner det Tingene i deres Nødvendighed i Modsætning til den Maade, hvorpaa de staar for blot Iagttagelse. 3) I det mærkelige Brev (Ep. 6), i hvilket \textsc{Spinoza} drøfter nogle Experimenter af \textsc{Boyle}, siges Begreber som Bevægelse, Hvile og deres Love (disse „kyske Begreber“, notiones castæ) at udtrykke Naturen som den er i sig selv, i Modsætning til Forestillinger (som Farve, Kulde, Varme, Fasthed, Flydenhed o. s. v.), der kun udtrykker deres Forhold til menneskelig Sansning. I denne Betydning bruges Ordet ogsaa af \textsc{Locke}, der (Essay II, 8, 23) siger om „de primære Egenskaber“ (Figur, Tal, Bevægelse o. s. v.), at de er i Tingene, hvad enten vi opfatter dem eller ikke, saa at de giver os „the idea of the thing, as it is in itself“. Denne Betydning stammer fra \textsc{Galilei}, \textsc{Hobbes} og \textsc{Descartes}. 4) Endelig bruges Udtrykket „res in se consideratæ“ i Fortalen til fjerde Bog (se ogsaa IV, 73 Schol.) til at betegne Modsætningen til de Værdibegreber, man er tilbøjelig til at anvende paa Tingene.}, fordi han bestaar
af uendelig mange Attributer“.

% --- p. 41 ---
\setcounter{footnote}{0}%
\opage{41}Saaledes som Sætningen her er opstillet, er den af erkendelsesteoretisk Art og
angaar Forholdet mellem videnskabelig udformede Tanker (objektive Begreber) og
andre Sider af Tilværelsen. Det ses ikke blot af Henvisningen til \textsc{Kabbala}'s Lære
om Gud, hvis Forstand og det, han forstaar, er et og det samme; det ses ogsaa af
\textsc{Spinoza}'s Exempel: Forholdet mellem den i Naturen existerende Cirkel og Cirkelens
Begreb.

Men som det ses i det følgende (Sætning 11---13 og III, 2. Schol.), betragter
\textsc{Spinoza} Forholdet mellem Sjæl og Legeme som et specielt Tilfælde af Forholdet
mellem Erkendelsen (det vil sige: de videnskabelige Begreber) og dens Genstande,
og gaar saaledes uden at mærke det over fra Erkendelsesteori til Psykologi (resp.
Psykofysik). Forholdet mellem det matematiske Begreb om en Cirkel og den i Naturen
existerende Cirkel er jo dog et andet end Forholdet mellem Begrebet som
psykisk Produkt og den fysiologiske Proces, det paa en eller anden Maade er knyttet
til, og som maa høre ind under Udstrækningens Attribut ligesaavel som „den existerende
Cirkel“. --- Endnu tydeligere bliver Forvexlingen, naar vi spørger om, hvorledes
det forholder sig med den Tankevirksomhed, der fører til at danne videnskabelige
(objektive, „formale“) Begreber. Svarer der ikke ogsaa til den noget indenfor
Udstrækningens Attribut? Og allertydeligst bliver Sagen, naar vi spørger om, hvorledes
det forholder sig med urigtige Begrebsdannelser og selvmodsigende Forestillinger.
Man kan tænke sig, at der ogsaa til disse svarer visse fysiologiske (hjernefysiologiske)
Processer. Men \textsc{Spinoza} antager ganske sikkert ikke, at der til Forestillingen „firkantet
Cirkel“ svarer nogetsomhelst, som „existerer i Naturen“.

Det er \textsc{Spinoza}'s Dogmatisme, der har ført ham til at sammenblande to Problemer.
Han er som alle dogmatiske Filosofer og Naturforskere vis paa, at der til
vore efter streng Metode dannede Begreber maa svare absolute Realiteter (se ovenfor
p. 16 f). Og hans strenge Opfattelse af den materielle Naturs indre Sammenhæng
gjorde ham det umuligt at antage en Virken af det ene Attribut paa det andet;
sjælelige Tilstande blev derfor for ham et andet Udtryk for visse materielle Tilstande,
og omvendt. Disse to Forhold har han identificeret. Og snart er det da hans dogmatiske
Erkendelsesteori, snart hans realistiske Psykofysik, der har Ordet. --- Vi vil
endnu flere Gange faa Lejlighed til at omtale dette Spørgsmaal.

\textit{γ.} I Sætning 8---10 indskærpes, at Begrebet Substans ikke kan bruges om et
Menneske, saaledes som Skolastiken og tildels endnu \textsc{Descartes} paa Grund af et
vagere Substansbegreb havde lært. Mennesket er et Enkeltvæsen, en Modus, og er
som saadant Led i en Række af Enkeltvæsener. Det har sin Aarsag i Gud, ikke
forsaavidt Gud betragtes som et uendeligt Væsen, men forsaavidt Guds Væsen lægger
sig for Dagen i Enkeltvæsenernes lovmæssige Sammenhæng. I de guddommelige
Attributer indeholdes Muligheder for uendelig mange Enkeltvæsener, ligesom
der i en Cirkel er Mulighed for et uendeligt Antal hinanden skærende Chorder, hvis
% --- p. 42 ---
\setcounter{footnote}{0}%
\opage{42}Stykker danner ligestore Rektangler, medens i hvert enkelt Tilfælde kun et endeligt
Antal kan være givet. I denne Sammenligning, som findes i Anmærkningen til
den ottende Sætning, svarer Cirkelen til Substansen, Chorderne til Modi.

\subsection*{c. Sætning 11---15.}
\begin{center}\textit{Om Forholdet mellem Sjæl og Legeme.}\end{center}

\textit{α.} Det almindelige Forhold mellem Tænkning og Udstrækning, der i det foregaaende
er bestemt som Konsekvens af Forholdet mellem Attributer i det hele, anvendes
nu særlig paa Forholdet mellem Sjæl og Legeme. Disse forholder sig altsaa,
som Tanken (idea) forholder sig til sin Genstand (res). Legemet (corpus) er den
Enkeltting (res singularis), som svarer til Bevidstheden (mens).

Naar det i Sætning 13 hedder: „Genstanden for den Tanke, der udgør den
menneskelige Bevidsthed (mens), er et Legeme, en vis Udstrækningsmodus, og ikke
andet“, maa Meningen være, at Menneskets Ejendommelighed bestaar deri, at der
til dets Tanke svarer en Udstrækningsmodus, og ikke en Modus af et af de andre
Attributer. I det Bevis, der gives for denne Sætning, henvises til Ax. 5, som slaar
fast, at vi ikke mærker eller opfatter andre Enkeltting end Legemer og Tankemodi.

--- Nu antager \textsc{Spinoza} jo dog, at der gives uendelig mange Attributer, der alle staar
i proportionalt Forhold til hverandre. Hvad der gælder for menneskelig Bevidsthed
i Forhold til Udstrækningens Attribut, maa altsaa konsekvent gælde for den i Forhold
til hvilketsomhelst andet Attribut, skønt her det ene Led i Forholdet er x.
(Smlgn. ovenfor p. 22---24).

Hvad der gælder for Mennesket, maa (13 Schol.) ogsaa gælde for alle Enkeltvæsener
indenfor vor Erfaring, i Kraft af det almindelige Forhold mellem Tænkningens
og Udstrækningens Attribut. Alle Enkeltvæsener maa være besjælede (animata),
skønt i forskellige Grader. Forskellen mellem de forskellige sjælelige Væsener svarer
til Forskellen mellem de forskellige Legemer, som navnlig frembyder højst forskellige
Forhold mellem Aktivitet og Passivitet. Allerede hvad selve Mennesket angaar,
er der store Grundforskelle i denne Henseende. Jo mere aktivt og selvstændigt
Legemet er, d. v. s. jo mere dets Funktioner afhænger af det selv, des fuldkomnere
vil ogsaa Sjælens Erkendelse være. Dunkelhed og Klarhed i Tankelivet svarer til
Passivitet og Aktivitet --- baade for Sjælens og Legemets Vedkommende. Ved de
dunkle Forestillinger tænker \textsc{Spinoza}, saavidt man kan se, paa Livsfornemmelser
og elementære Lyst- og Ulystfølelser og blotte Forestillingsassociationer, maaske ogsaa
paa den dunkle Trang, der ikke er sig sit Maal bevidst. Ifølge Def. 3 betegner
Ordet idea egentlig kun aktivt udformede Tanker; men der skelnes i Ax. 3 mellem
saadanne Tanker og andre Tankemodi (amor, cupiditas og andre affectus animi).
Som allerede bemærket (se ovenfor p. 38) fremtræder der her en Vanskelighed for
\textsc{Spinoza}'s Psykologi, idet disse passive Elementer stedse skal forudsætte de aktive,
men ikke omvendt.

\textit{β.} Efter Sætning 13 indskyder \textsc{Spinoza} en Række Laanesætninger (Lemmata)
% --- p. 43 ---
\setcounter{footnote}{0}%
\opage{43}fra Fysiken, naturligvis saaledes som han forstod sin Tids Fysik. Efter først at
have slaaet fast, at intet Legeme er Substans, opstilles Inertisætningen (Lemma 3),
som allerede er omtalt ovenfor (p. 9 f). Dernæst hævdes, at af Bevægelse følger Bevægelse,
af Hvile kun Hvile. --- Her ses tydelig den Hovedvanskelighed i den kartesianske
Fysik, som først blev hævet af \textsc{Leibniz}, nemlig, hvorledes en hvilende
Ting kan komme i Bevægelse, eller en Ting i Bevægelse komme i Hvile, naar Hvile
og Bevægelse er to helt forskellige Tilstande. Allerede \textsc{Hobbes} havde iøvrigt, i \textsc{Galilei}'s
Spor, antydet Muligheden af en Løsning gennem Begrebet Bevægelsestendens
(conatus), der ikke forsvinder, fordi den hæmmes ved en modsat Bevægelsestendens\footnote[1]{\textit{Den nyere Filosofis Historie\textsuperscript{2}}, I, p. 268. --- Ogsaa \textsc{Gassendi} korrigerede paa denne Maade \textsc{Descartes} (ib. p. 254).}.


Rent dogmatisk fremstilles derefter, hvorledes de simple Elementer, af hvilke
Legemerne bestaar, ved deres forskellige Lejringsforhold bevirker Forskellen mellem
Tilstandsformerne, mellem faste, bløde og flydende Legemer. Flere Legemer kan igen
forbindes til et Hele gennem fælles Virken, et Hele, der kan bestaa, selv om de
enkelte Dele skifter. Af saadanne Helheder kan der igen opstaa mere sammensatte
individuelle Helheder, og tilsidst udgør hele Naturen et stort Individ, hvis Dele, de
enkelte Legemer, forandres paa mange Maader uden Forandring i Naturen som Helhed.
--- Hvad \textsc{Spinoza} her kalder „hele Naturen“, er aabenbart det samme, som
paa andre Steder (se ovenfor p. 29) kaldes „hele Universets Beskaffenhed“ (Verdensfysiognomiet),
og som gennem det konstante Forhold mellem Hvile og Bevægelse
har sin Grund i Udstrækningens Attribut.

At de enkelte Forandringer i Verden (variatio rerum) ikke kan udledes af
Udstrækningens Attribut, og at \textsc{Descartes} med Urette definerer Materie som blot
Udstrækning, indrømmer \textsc{Spinoza} i det sidste Brev (Ep. 83), der haves fra hans
Haand, og han bebuder en nærmere Drøftelse, som han ikke kom til at foretage.

Endelig slaas det fast (i Postulaters Form), at det menneskelige Legeme er
sammensat af mindre Helheder, der er indbyrdes forskellige og hver for sig sammensatte,
--- at Indvirken fra andre Legemer kan efterlade sig Spor (vestigia), der
bestaar deri, at en af det paavirkede Legemes flydende Dele forandrer en blød Del
ved at ændre dens Lejring, --- og at vort Legeme stadig trænger til at genfødes ved
Paavirkning fra andre Legemer. Og fra denne sammensatte Karakter af det menneskelige
Legeme sluttes der (Sætning 15), i Kraft af Proportionaliteten mellem Sjæl
og Legeme, at den Tanke, som udgør den menneskelige Bevidstheds Væsen, er
sammensat af saare mange Tanker.

\textit{γ.} Sin Teori om Forholdet mellem Sjæl og Legeme anvender \textsc{Spinoza} i de
specielle Dele af Ethica saaledes, at han snart fra den Kundskab, vi har (eller som
han mener, vi har) om Legemet, slutter til, hvad der foregaar i Sjælen, snart fra
det, der vides at foregaa i Sjælen, slutter til, hvad der foregaar i Legemet. Begge
Dele er berettigede, da jo Forholdet mellem Sjæl og Legeme ifølge \textsc{Spinoza} er et
rent logisk eller matematisk Funktionsforhold, uden at noget af Leddene en Gang
% --- p. 44 ---
\setcounter{footnote}{0}%
\opage{44}for alle kan erklæres at være den uafhængige Variable. Alle Attributer er samtidige
Udtryk for Substansens Væsen (simul natura. I, 17 Schol., III, 2 Schol.). Man
har med Urette fundet en Inkonsekvens i, at \textsc{Spinoza} snart gaar den ene Vej, snart
den anden. Han gaar fra det legemlige til det aandelige i Sætning 16---19; den
modsatte Vej gaar han f. Ex. III, 12 Dem. og især V, 1\footnote[1]{Det har sin Interesse at sammenligne \textsc{Avenarius}' Lære om, hvad han kalder de to Vitalrækker (den fysiologiske og den psykologiske) med \textsc{Spinoza}'s Lære om Attributerne. Se \textit{Moderne Filosofer}, p. 101 f. --- \textsc{Avenarius} erklærer paa Forhaand den fysiologiske Vitalrække for den uafhængige Variable.}.

Angaaende Forholdet mellem de to Rækker af Emner bruger \textsc{Spinoza} enten
Udtrykket „antyde“ eller „angive“ (indicare) (II, 16 Cor. 2; III, 14 Dem.) eller Udtrykket
„forudsætte“ (involvere) (II, 16; III, 29 Dem.). Det første Udtryk betegner
Forholdet som symptomatisk; det andet, som er korrektere, betegner det som et
rent logisk Slutningsforhold, et Forhold mellem en Erkendelsesgrund og det, der
udledes af den. Dog bruger \textsc{Spinoza} ogsaa ofte det Udtryk, at Sjælen „opfatter“
(percipit) Legemet, et Udtryk, der hænger sammen med den ovenfor omtalte Forvexling
af Erkendelsesteori og Psykofysik, idet Forholdet mellem Sjæl og Legeme
betragtes som et Forhold mellem Subjekt og Objekt.

Maaske er det ikke overflødigt at forklare Forskellen mellem de to Synspunkter
ved et Exempel. Fysikeren slutter fra en Farvefornemmelse til en vis Brydningsgrad
af det hvide Lys. Det er et erkendelsesteoretisk Forhold. Fysiologen slutter
fra den samme Farvefornemmelse til visse Processer i Sanseorgan og Hjerne. Dette
er et psykofysisk Forhold. Ved enhver Sansefornemmelse gør begge Synspunkter
sig gældende. Og ligesaa ved enhver Tanke, enhver Følelse, enhver Villen. Da
\textsc{Spinoza} tænkte Begrebet Substans, skulde derved udtrykkes et erkendelsesteoretisk
Princip; men samtidig foregik der noget i hans Hjerne (blot vi vidste, hvad det var!).
En heroisk Beslutning (hvad \textsc{Spinoza} kalder mentis decretum som species fortitudinis,
se III, 2 Schol. og 59 Cor.) angaar forskellige Forhold i og udenfor det besluttende
Subjekt, og herpaa beror dens etiske Betydning; men fra et andet Synspunkt
er en saadan Beslutning et fysiologisk Fænomen (ifølge \textsc{Loeb} et rent kemisk Fænomen),
der tilsidst maa kunne udledes „af Lovene for Bevægelse og Hvile“. Hvilket
Synspunkt, man anlægger, beror naturligvis paa Sammenhængen i det enkelte Tilfælde.


Hvis man vilde sige, at Tanken (Bevidstheden) selv kunde være en materiel
Proces, vilde \textsc{Spinoza} svare, som han har svaret i et Brev (Ep. 4): „Selv om jeg
indrømmer dette, hvad jeg ingenlunde gør, kan det dog ikke nægtes, at Udstrækning,
som Udstrækning betragtet, ikke er Tanke“\footnote[2]{I Afhandlingen om Forstandens Forbedring anfører \textsc{Spinoza} blandt Exempler paa uklare og falske Forestillinger den Antagelse, at der gives Legemer, af hvilke Tanken opstaar ved simpel Sammensætning (dari corpora, ex quorum sola compositione fiat intellectus).}. --- I Anledning af denne Bemærkning
kan endnu tilføjes, at det erkendelsesteoretiske Synspunkt maa finde Anvendelse,
selv om man holder sig til den fysiologiske Side af Menneskets Væsen. Problemet
bliver da, hvorfor og hvorledes man kan betragte en vis Hjernetilstand som Tegn
paa noget, der maaske har en helt anden fysisk Beskaffenhed end den, maaske for%
% --- p. 45 ---
\setcounter{footnote}{0}%
\opage{45}holdende sig til den som en optisk Proces til en kemisk. \textsc{Spinoza} giver selv i det
nærmest følgende Antydninger i den Retning.

\subsection*{d. Sætning 16---36.}
\begin{center}\textit{Forestillingernes Psykologi.}\end{center}

\textit{α.} Læren om Forestillingerne indleder \textsc{Spinoza} ud fra Udstrækningens Attribut.
Vort Legeme staar i stadig Vexelvirkning med andre Legemer, og dets Tilstand er
derfor i hvert Øjeblik ikke blot bestemt ved dets egen Natur, men ogsaa af Beskaffenheden
af de andre Legemer, der paavirker det. I Sjælens Opfattelse (perceptio)
af sit eget Legeme indeholdes forsaavidt ogsaa en Opfattelse af andre Legemer,
men betinget ved den Maade, hvorpaa det egne Legeme er blevet paavirket af dem.
Den betragter dem, saa at sige, fra det egne Legemes Standpunkt. Dens Erkendelse
af andre Legemer kan derfor ikke blive fuldkommen (adæqvata) (25), men staar
som en Konsekvens uden Præmisser, --- og det samme gælder Sjælens Erkendelse
af det egne Legeme, da dette stedse er bestemt ved Forholdet til andre Legemer (28).
--- Hvad \textsc{Spinoza} her fremsætter, er, hvad man nu kalder Sansekvaliteternes Subjektivitet.
Han fremsætter den paa en mindre dogmatisk Maade, end man efter hans
Attributlære skulde vente.

\textit{β.} Ofte opfatter vi paa én Gang flere andre Legemer, og naar vi saa senere
igen opfatter et af dem, kommer vi til at mindes de andre. Her virker nemlig de
Spor (vestigia)\footnote[1]{Disse Spor er materielle (III, Postulat 5). Undertiden kaldes de materielle Processer, som svarer til Erindringer og Fantasier, for Billeder (imagines), da Tingene staar som nærværende. \textsc{Spinoza} gør en Undskyldning for dette Udtryk, fordi hine materielle Processer ikke gengiver Tingenes Udseende (II, 17 Schol.). I III, Postulat 2 skelnes mellem Spor og Billeder, og disse betragtes som Følge af hine (dog som materielle!). Smlgn. ogsaa V, 11---14. --- Men i II, 49 Schol. bruges „Billede“ psykologisk.}, som den første Opfattelse har efterladt (18). --- Det er den saakaldte
Berøringsassociation, \textsc{Spinoza} her omtaler. Andetsteds (Tract. Theol. polit. cap. 4)
anfører han Lighedsassociationen som en selvstændig Associationsform, og i sin
Forklaring af Medlidenhed (III, 27) forudsætter han i hvert Tilfælde umiddelbar
Genkendelse som Betingelse. Naar han IV, 18 Schol. siger, at et Ord kan fremkalde
Forestillingen om en Genstand, skønt der ingen Lighed er mellem dem, ligger heri
egentlig den Forudsætning, at Forestillingsassociationen gaar lettere, naar der er
Lighed mellem Leddene.

Forestillingsassociationen foregaar uvilkaarlig, og vi har fra først af ingen Grund
til at tvivle om Gyldigheden af de Forestillinger, der opdukker hos os. Hvad vi
forestiller os, betragter vi foreløbig som gyldigt, indtil der kommer andre Forestillinger,
som udelukker Gyldigheden. I sig selv indeholder vore Forestillinger ingen
Vildfarelse; vi mangler kun en Forestilling, der kunde udelukke Gyldigheden af de
foreliggende Forestillinger (17). Ligesom \textsc{Spinoza} har kaldt Sansefornemmelserne
Konsekvenser uden Præmisser, saaledes kunde han efter sin Opfattelse ogsaa have
brugt denne Betegnelse om uvilkaarlig opstaaende Forestillinger.

% --- p. 46 ---
\setcounter{footnote}{0}%
\opage{46}Denne interessante Betragtningsmaade har senere den engelske Psykologi udført
videre, og dels sat den i Forbindelse med en uvilkaarlig Tendens til at fortsætte
en Forestillingsrække, indtil der kommer bestemte Hindringer (\textsc{Hume}), dels med den
uvilkaarlige Bevægelsestrang, der gør, at vi uden Kritik handler efter uvilkaarlig
opdukkende Forestillinger, indtil denne Handlen medfører Skuffelser eller Lidelser
(\textsc{Bain}).

\textit{γ.} Meget skarpt skelner \textsc{Spinoza} derpaa mellem den blotte Erindring eller
Forestillingsassociation og den Forbindelse af vore Forestillinger, som finder Sted
efter Forstandens Regler (secundum ordinem intellectus), og som fører til Forstaaelse
af Tingene ud fra deres første Aarsager, uafhængig af de for hvert Individ forskellige
Tilstandsændringer (18 Schol.). Den sande Erkendelse maa omfatte de Love,
der gælder for alle Legemer og for alle Sjæle. Den menneskelige Bevidsthed har
kun en ufuldkommen Viden om sig selv, om sit Legeme og om de ydre Ting, saalænge
Mennesket udefra, nemlig ved tilfældigt Sammentræf, bestemmes til sine Forestillinger
om Tingene. Anderledes naar Forestillingerne bestemmes ved en Undersøgelse
af Tingenes Overensstemmelser, Forskelligheder eller Modsætninger (29 Schol.).

I Substansen (Gud) er der en fuldkommen Viden (idea seu cognitio), fordi den
paa én Gang virker under alle Attributer og i alle Modi, ikke blot under et enkelt
Attribut eller i en enkelt Modus (19---32). Usandhed skyldes Begrænsning og Ufuldstændighed
og deraf følgende Forvirring. En uklar Fornemmelse eller Forestilling
(idea confusa) er, som vi har set, en Konsekvens uden Præmisser (28 Dem.); kan
Præmisserne skaffes tilveje, hører Uklarheden op. Men disse ufuldkomne og dunkle
Forestillinger opstaar med samme Nødvendighed som de fuldkomne, klare og sande
Forestillinger --- og for Gud, for hvem denne Nødvendighed staar klart, existerer
derfor ingen dunkle Tanker. Dunkelheden er kun tilstede, forsaavidt Forestillingerne
henføres til en enkelt og begrænset Bevidsthed (36).

Man har fundet en Vanskelighed i, at \textsc{Spinoza} tillægger den menneskelige Bevidsthed
en Forestilling om sig selv, saa at den kan være Objekt for sig selv, ligesom
Legemet er Objekt for den (21). \textsc{Spinoza} har selv givet den fornødne Forklaring,
idet han siger, at denne Forestilling om vor Bevidsthed (idea mentis) --- der, da Bevidsthed
selv er en Forestilling, er en Forestilling om en Forestilling (idea ideæ) --- ikke er
andet end Formen for Bevidsthed eller for Forestilling, betragtet som Tankemodus,
bortset fra Forholdet til Legemet (21 Schol.). Dette svarer til, hvad der siges i Afhandlingen
om Forstandens Forbedring, at Tænkningens Metode er en reflexiv Erkenden
(cognitio reflexiva), ved hvilken vi bliver os vore Tankers Form og Fremgangsmaade
bevidst. Kun den, der har sande Tanker, kan vide, hvad Sandhed er. --- Naar \textsc{Spinoza}
siger, at denne Forestilling om os selv forholder sig til vor Bevidsthed (vor
Sjæl), paa samme Maade (eodem modo), som denne forholder sig til Legemet, saa
betyder naturligvis „samme Maade“ kun Analogi, ikke Identitet. Der er intet i Vejen
for, at der indenfor et og samme Attribut kan gives Forhold, der er analoge med
Forholdet mellem to Attributer. --- Efter Attributernes psykofysiske Betydning vil
% --- p. 47 ---
\setcounter{footnote}{0}%
\opage{47}Opgaven for denne Selvbevidstheds Vedkommende, ganske ligesom for alle andre
Bevidsthedsfænomeners Vedkommende, være at finde, hvilken Hjerneproces der
foregaar, naar den optræder.

\subsection*{Sætning 37---47.}
\begin{center}\textit{Erkendelsesteori.}\end{center}

\textit{α.} Det er Tidens Naturvidenskab, der giver \textsc{Spinoza} Forstaaelse af, hvad Videnskab
overhovedet er. Han gaar ud fra, at alle Legemer stemmer overens i visse
Egenskaber, navnlig deri, at de er underkastede Lovene for Bevægelse og Hvile og
tilsidst alle hører ind under et eneste Attribut, Udstrækningens Attribut. Det er de
almindelige matematiske og mekaniske Love, han tænker paa. Deres Indhold er
noget, som i lige Grad fremtræder i Tilværelsen som Helhed og i dens enkelte Dele
(37). --- Begrebet Helhed er hos \textsc{Spinoza} afledet i Forhold til Begrebet Lovmæssighed.
Som vi har set (ovenfor p. 29 og 43), opfatter han Universet som en individuel
Helhed, og denne Opfattelse begrundes paa, at der i Universet bestaar et konstant
Forhold mellem Hvile og Bevægelse trods alle Forandringer i det enkelte\footnote[1]{Udtrykket „det, som i lige Grad er i Delen og i det Hele“ (id, quod æqve in parte et in toto est) har sin Historie. \textsc{Spinoza} kan have det fra \textsc{Bruno}, som lærer, at „Verdenssjælen“ (der er Principet for Naturens Lovmæssighed og ogsaa kaldes „forma universalis“) er „tutta in tutto et qualsivogla parte del tutto“ (\textit{Op. ital.} ed. Lagarde p. 203, 242). (Ligesaa \textit{De Immenso} IV, 15). Men han kan ogsaa have det fra Skolastikerne. Det gaar tilbage til \textsc{Augustinus} og \textsc{Athanasios}, der sikkert har det fra \textsc{Plotinos}. Se herom min \textit{Religionsfilosofi}, Note 15. En Antydning findes allerede hos \textsc{Platon} (\textit{Parmenides} p. 131 B). --- \textsc{Hume} (\textit{Treatise} I, 4, 5) omtaler Udtrykket som skolastisk og finder det „shocking when crudely propos'd“. Han tænker vel paa den mytologiske og dogmatiske Anvendelse af Begrebet. Men dets Historie er interessant, fordi den viser, at Tanken om Tilværelsens Lovmæssighed trænger sig frem, saasnart Mennesker begynder at tænke. Denne Tanke gennemløber forskellige Former. Hos \textsc{Spinoza} faar den en afgjort rationel Form, skønt den fremsættes dogmatisk.}.

Den Erkendelse, der angaar Tilværelsens almindelige Love, er altsaa ikke betinget
ved noget Enkelt i dettes Forskellighed fra andre enkelte Ting. Eller, som
\textsc{Spinoza} udtrykker det i sit metafysisk-teologiske Sprog, Guds Væsen fremtrædrr % PRINTED AS IS: fremtrædrr
paa samme Maade, hvad enten han betragtes som udgørende Menneskets Væsen
eller som udgørende en anden Tings Væsen. Der er jo her Tale om det, som Menneskets
Legeme har fælles med alle andre Legemer, og derfor kan Mennesket her
have en fuldkommen (adækvat) Erkendelse (37---39). Vi staar her overfor, hvad
\textsc{Spinoza} i Afhandlingen om Forstandens Forbedring har kaldt „de faste og evige
Ting og de i dem indeholdte Love“, der svarer til, hvad Ethica kalder uendelige
Modi (af anden Grad); og sidste Grundlag er her Udstrækningens Attribut.

Paa dette Punkt falder Erkendelsesteori og Psykologi sammen. Havde Sjælelivet
ikke andet Indhold end den videnskabelige Loverkendelses, vilde der derfor
ikke være nogen Modsætning mellem den erkendelsesteoretiske og den psykofysiske
Tydning af Attributlæren. Men \textsc{Spinoza} anerkender selv, at der gives andet sjæleligt
Indhold end den videnskabelige Loverkendelse. Dels hævder han, som vi skal
se, at Loverkendelsen ikke er den højeste Form for Videnskab, fordi den ikke med%
% --- p. 48 ---
\setcounter{footnote}{0}%
\opage{48}tager de individuelle Fænomener i deres Ejendommelighed. Dels hævder han Vildfarelsens
Realitet, altsaa den Kendsgerning, at der gives Forestillinger, som strider
mod, hvad videnskabelig Erkendelse godtgør. Dels endelig antager han Følelser og
Drifter, som har en vis Selvstændighed i Forhold til den blotte Tanke. Og til alle
disse irrationelle Elementer maa der (da de dog alle hører ind under „Tænkningens“
Attribut) svare Elementer indenfor Udstrækningens Attribut. Her viser den psykofysiske
Tydning sig i sin Modsætning til den erkendelsesteoretiske. --- Men vi vender
tilbage til Erkendelsesteorien.

Det er de almene Love (notiones communes), som danner Grundlaget for al
Tænken over Naturen. Fundamenta ratiocinii nostri kaldes de (40 Schol. 1). De ser
bort fra alle Tidsforskelligheder og kan derfor siges at vise os Tingene under Evighedens
Synspunkt (sub qvadam æternitatis specie, 44 Cor. 2). Ved Evighed mener
\textsc{Spinoza} jo ikke bestandig Bestaaen, men en Existens, der med Nødvendighed følger
af Tingenes Væsen (I, Def. 8). Og denne Tingenes Nødvendighed er (ifølge I, 16) ét
med Guds evige Naturs Nødvendighed (44 Cor. 2).

Udtrykket notiones communes stammer fra Stoikerne, og det betegner oprindelig
Begreber, der er fælles for alle Mennesker\footnote[1]{Smlgn. \textsc{Paul Barth}: Die Stoa. (Frommanns Klassiker der Philosophie). 1903. p. 72 f.}. Ogsaa \textsc{Spinoza} antager, at der
gives visse Begreber, som er fælles for alle Mennesker, og som (delvis) opfattes af
alle „fuldkomment, d. v. s. klart og tydelig“; men dette (mulige) Fællesskab grundes
hos ham ved Slutning derfra, at disse Begreber kun angaar, hvad der gælder „i lige
Grad for det Hele og for enhver Del“ (38 med Cor.).

Om den Vej, ad hvilken der naas til Erkendelse af de almene Love, som vi maa
støtte os til i vor Forstaaelse af Tilværelsen, vil \textsc{Spinoza} ikke tale nærmere, men
henviser til Afhandlingen om Forstandens Forbedring, som han jo desværre ikke
fik fuldendt. Han har dog paa et tidligere Sted i Ethica (29 Schol.) givet en lille
Antydning, idet han i Modsætning til den „udvortes“ Opfattelse, som kun skyldes
Sammentræf af tilfældige Omstændigheder, stiller den Erkendelse, der kan kaldes
„indre“, fordi den beror paa Sammenligning og derved fører til Forstaaelse af Overensstemmelser,
Forskelle og Modsætninger; ved en saadan Fremgangsmaade kommer
man til en klar og tydelig Opfattelse af Tingene. (Smlgn. ovenfor p. 46). Det er jo
ogsaa paa Muligheden af at ordne vore Tanker i Rækker efter bestemte Love, at han
begrunder Attributets Begreb. (Se ovenfor p. 8 og p. 39). Af Interesse er i denne
Sammenhæng ogsaa hans Forhandling med den berømte Kemiker \textsc{Boyle} (med \textsc{Oldenburg}
som Mellemled). Overfor \textsc{Boyle} hævder han, at intet Experiment kan
godtgøre Berettigelsen af den mekaniske Naturopfattelse, der fører alle sanselige
Egenskaber tilbage til Bevægelse og Hvile og Lovene for disse (Ep. 6). Det synes
at være \textsc{Spinoza}'s Mening, at ethvert Experiment kun kan afgive et Exempel paa
disse Love, men at disse staar som den stadige Forudsætning, der ifølge \textsc{Spinoza} er
en evig og nødvendig Sandhed. Trods den dogmatiske Tankegang frembyder hans
Stilling her dog en Analogi til \textsc{Kant}'s. Han har staaet overfor sin Tids Naturvidenskab
paa lignende Maade, som \textsc{Kant} stod overfor \textsc{Newton}'s Principia, betragtende
% --- p. 49 ---
\setcounter{footnote}{0}%
\opage{49}den som en aandelig Kendsgerning, hvis Grundlag skulde findes ad Analysens Vej.
Analogien ligger saa meget nærmere, som der paa dette Punkt endnu er et dogmatisk
Element hos selve Kriticismens Grundlægger\footnote[1]{Smlgn. \textit{Totalitet som Kategori}, p. 36 f.}.

Fra Begreberne om almene, gennemgaaende Love (notiones communes) skelner
\textsc{Spinoza} saadanne Begreber eller Ord, der kun har almen Betydning, fordi der har
manglet Evne til paa tydelig Maade at sammenfatte forskellige Emner. Herhen
hører saadanne Ord som „Ting“, „Væsen“, „Noget“, og desuden alle Artsbegreber
(notiones universales) som „Menneske“, „Hest“, „Hund“ o. s. v. (40 Schol. 2)\footnote[2]{Det har sin Interesse at sammenligne \textsc{Spinoza}'s Kritik af Almenbegreberne med \textsc{Leibniz}'. For \textsc{Spinoza} var Hovedpunktet, at Almenbegreberne ikke angiver virkelige Love for Forholdet mellem de forskellige Kendemærker eller Egenskaber. Naar \textsc{Leibniz} tvivler om Almenbegrebernes Gyldighed, er det, fordi han er overbevist om, at „Naturen har Forskelle og Ligheder, som vi ikke kender“. Derfor maa vore Bestemmelser af de naturlige Arter stedse være foreløbige og betingede ved vore Kundskaber til given Tid. Desuden kan Overgangen fra Art til Art være kontinuerlig, hvorfor Artsbegreber kun kan opstilles rent vedtægtsmæssig. (\textit{Nouveaux Essais}. III, 6, 14. 23. 28).}.
Denne Skelnen minder om, hvorledes \textsc{Bruno} og \textsc{Baco} tolkede \textsc{Platon}'s „Ideer“
som Love, ikke som blotte Almenbegreber, hvad de oprindelig var.

Til de blotte Almenbegreber regner \textsc{Spinoza} ogsaa Talbegrebet. „Vi opfatter“,
siger han (Ep. 50), „ikke Tingene under Tallets Begreb, uden efter at have ført dem
tilbage til en fælles Art“. Jeg kan sige, at jeg har to Mønter, naar jeg har en Daler
og en Skilling. Og da nu Artsbegreber, som vi har set, ifølge \textsc{Spinoza} er Tegn paa
ufuldkommen Tænkeevne, kan han ikke tillægge Talbegrebet nogen fremragende
Plads i Tænkningen. Det kan ikke anvendes paa Substansen. „Maal (mensura),
Tid og Tal er kun Tanke- eller rettere Forestillingsformer (imaginandi modi)“
(Ep. 12). Det er kun Aritmetiken, ikke Geometrien, \textsc{Spinoza} har Betænkelighed ved.
Den Forskel, han her gør, udspringer vel af, at Kontinuiteten fremtræder tydeligere
paa det geometriske end paa det aritmetiske Omraade; Tallet er Udtryk for Diskontinuitet.
Men af \textsc{Descartes}' analytiske Geometri (som fandtes i hans Bibliotek)
kunde han have lært, at geometriske Former kan udtrykkes ved aritmetiske
Formler (Cirkelen f. Ex., der jo ifølge \textsc{Spinoza} „existerer i Naturen“, ved Ligningen
r\textsuperscript{2} = y\textsuperscript{2} + x\textsuperscript{2}), saa at Tal og Udstrækning kan følges ad, idet Formlerne ændres eftersom
Linierne og deres Forhold ændres, og omvendt. Maaske har dog \textsc{Descartes}'
analytiske Geometri paa anden Maade faaet Indflydelse paa \textsc{Spinoza}'s Tænkning,
idet Forholdet mellem de forskellige Attributer netop er af samme Art, som ifølge
\textsc{Descartes} Forholdet mellem Ligning og Kurve\footnote[3]{Se den interessante Betragtning hos L. \textsc{Bbunschwicg}: \textit{Les Étapes de la philosophie mathématique}. p. 144 f.}. Men den kvalitative Diskontinuitet, % PRINTED AS IS: BBUNSCHWICG (BRUNSCHWICG, p. 53)
som overhovedet medfører saa store Vanskeligheder for \textsc{Spinoza}'s System, har han
ikke kunnet anerkende som tankemæssig og henviser den derfor til den blotte
„Imagination“.

Det var først \textsc{Leibniz}, som saa, at Tallets Opstaaen skyldes Evnen til at fortsætte
en Række efter samme Lov, efter hvilken den var begyndt, og at denne 
Fort%
% --- p. 50 ---
\setcounter{footnote}{0}%
\opage{50}sættelse stedse stod i vor Magt. Det er Gentagelsens positive Betydning, der her
fremtræder og betegner Grænsen mellem Logik og Matematik. En ny Fase i Erkendelsesproblemets
Historie begynder med, at \textsc{Kant} optager dette Synspunkt som
afgørende, hvis rationel Erfaringsvidenskab skal ses i sin Begrundelse og Berettigelse\footnote[1]{Smlgn. \textit{Totalitet som Kategori}, p. 34---37.}.

En anden, og mere berettiget Brug af sin kritiske Holdning overfor Artsbegreber
gør \textsc{Spinoza} i sin Kritik af Antagelsen af en „fri Villie“ i den Betydning af
Ordet, at enhver skulde have samme Evne til Selvbeherskelse. Denne Antagelse
skyldes, mener han, et abstrakt Begreb om Mennesker og overser den store Forskellighed
mellem Menneskers Evner og mellem de Veje, ad hvilke de enkelte maa
arbejde sig frem til Fuldkommenhed. (Ep. 19 og 23).

\textit{β.} Det er ikke \textsc{Spinoza}'s Mening, at kun den materielle Natur er underkastet
almene Love. I Læren om Attributernes indbyrdes Forhold ligger den Forudsætning
skjult, at det aandelige ikke kan udledes af det materielle og derfor maa betragtes
som en særegen Side, et særligt Attribut ved Substansen. Hvad der er fælles for
de forskellige Attributer, er jo netop Forholdet mellem Grund og Følge, det Forhold,
som indenfor de forskellige Attributer fremtræder i forskellige Former. Alle
Enkeltting (ligegyldigt, under hvilket Attribut de henføres) er underkastede en Verdensorden,
der bestemmer deres Existens til given Tid og Sted, og som tilsidst har sin
Grund i den guddommelige Naturs Nødvendighed. (II, 45 Schol. og V, 29 Schol.).
Hvad de Love, der gælder under Tænkningens Attribut, angaar, kunde man hos
\textsc{Spinoza} --- formedelst den Forvexling af erkendelsesteoretiske og psykologiske
Synspunkter, som ovenfor er paavist --- ligesaavel vente en Henvisning til logisk-matematiske
som til empirisk-psykologiske Sammenhænge. Medens man, hvor
Talen er om Muligheden af stedse fortsatte Rækkedannelser under Tænkningens
Attribut (Eth. II, 1), nærmest vilde tænke paa logisk-matematiske Begrebsrækker
(se ovenf. p. 9 Note), saa henviser \textsc{Spinoza} selv, hvor han (Tr. theol. polit. cap. 4)
skal nævne Exempler paa Lovmæssighed i Naturen, Loven for Forestillingsassociation
som Korrelat til Bevægelseslovene paa det materielle Omraade. Og i tredie og
fjerde Bog, samt Begyndelsen af femte, viser han, hvorledes der formedelst Forestillingsassociation
kan finde Forskydninger Sted paa Følelses- og Villieslivets Omraade,
og han slutter fra disse Fænomener til korrelate materielle Processer (V, 1).
Og paa den anden Side finder han paa det aandelige Omraade et Sidestykke til den
paa det materielle Omraade gældende Inertisætning (III, 6) i Selvopholdelsestrangen,
der spiller saa stor en Rolle i hans Etik. Det var overhovedet hans Overbevisning,
at „hvad enten man betragter Naturen under det ene eller det andet Attribut, saa
er det en og samme Tingenes Orden eller Sammenhæng, man har for sig“ (III, 2
med Henvisning til II, 7 Schol.). Og paa dette Grundlag bygger han i Ethica's tre
sidste Bøger sin Lære om det menneskelige Livs højeste Udvikling.

\textit{γ.} Hvad Erkendelsens faktiske Udvikling angaar, skelner \textsc{Spinoza} mellem tre
Trin (II, 40 Schol. 2). Paa det første Trin betinges Erkendelsen ved enkelte spredte
Iagttagelser, der staar som Konklusioner uden Præmisser. (Smlgn. ovenfor p. 45).
% --- p. 51 ---
\setcounter{footnote}{0}%
\opage{51}I Afhandlingen om Forstandens Forbedring gives der en nærmere Karakteristik af
dem. Deres spredte og tilfældige Optræden (sensationes fortuitæ et solutæ) fremhæves
og sættes i Forbindelse med, at de ikke udspringer af Bevidsthedens egen
Energi, men skyldes ydre Aarsager. Spredtheden og Passiviteten hænger altsaa
nøje sammen. Det er, føjer \textsc{Spinoza} til, ligegyldigt, hvorledes man tænker sig de
sanselige Forestillinger og deres Opstaaen\footnote[1]{Ogsaa \textsc{Hume} (\textit{Treatise} I, 3, 5) lader det [erkendelsesteoretisk] staa hen, hvorledes Sansefornemmelserne opstaar. Hvis \textsc{Kant} havde indtaget en lignende forsigtig Stilling, vilde Begrebet Ding an sich neppe være blevet dannet.}; det afgørende er deres vage og passive
Karakter. Med disse spredte og passive Iagttagelser sammenstiller \textsc{Spinoza} de Forestillinger,
som skyldes Beretning og Overlevering, og som formes i Lighed med,
hvad vi selv har set og hørt. Dette første Erkendelsestrin kalder \textsc{Spinoza} ubestemt
Erfaring (expeenritia vaga), Mening (opinio) eller Indbildning (imaginatio). Det % PRINTED AS IS: expeenritia
andet Trin er det nys beskrevne videnskabelige Standpunkt, der fører til Erkendelse
af almene Love og støtter sig til kritisk Sammenligning af Iagttagelser og Beretninger
(29 Schol.). Foruden Naturvidenskaben og Psykologien vilde \textsc{Spinoza}
hertil kunne regne den historiske Kritik, saaledes som han selv har anvendt den
paa de bibelske Bøger i sin teologisk-politiske Afhandling. Dette andet Trin kalder
\textsc{Spinoza} Fornuft (ratio). Allerede her erkender vi Tingene „under Evighedens Synspunkt“,
fordi vi støtter os til det, som er uafhængigt af Tider, Steder og individuelle
Forskelligheder (44 Cor. 2). Det tredie Trin betegnes som anskuende Viden (scientia
intuitiva).

Den anskuende Viden staar baade i Modsætning til den Kundskab, der bygger
paa Iagttagelse eller Overlevering og derfor kun angaar enkelte Emner (singularia),
og til den rationelle Erkendelse, der paaviser almene Love. Den forener begges
Ejendommelighed i sig, idet den er en Erkendelse af en enkelt Tings Væsen gennem
Forstaaelse af den lovmæssige Sammenhæng, indenfor hvilken Tingen staar som
et ejendommeligt Led. Det Exempel, hvormed \textsc{Spinoza} oplyser denne Beskrivelse
af anskuende Viden, er hentet fra simple aritmetiske Proportioner. Vi ser umiddelbart,
at 6 er det Tal, der forholder sig til 3, som 2 til 1. Til Grund ligger ganske
vist en Viden om Tallenes Natur og om proportionale Forhold; men denne Viden
anvendes ikke med fuld Bevidsthed. I dette Exempel (II, 40 Schol. 2) betragter
\textsc{Spinoza} aabenbart Tallet 6 som et Individ, en res singularis, der i Kraft af Proportionalitetens
Lov indordner sig paa en bestemt Maade i en hel Sammenhæng.

Denne tredie og højeste Art af Erkendelse kan kun naas af den, der har gennemløbet
de to første Grader. Indsigt i Tilværelsens Lovmæssighed er en Betingelse
for dens Opstaaen. Herved skiller Ethica's Intuitionsbegreb sig fra det mere mystiske
Begreb om Intuition i den korte Traktat. I femte Bog af Ethica siges Intuitionen
at bestaa i, at der ud af en enkelt Tings eget Væsen (ex ipsa essentia rei cujuscunqve
singularis) vindes Indsigt i dens Afhængighed af Tilværelsens Substans,
baade hvad Væsen og hvad Existens angaar, og der tilføjes, at der i første Bog kun
var vist, hvorledes Tingene i Almindelighed var afhængige af Substansen (V, 36 Schol.).
% --- p. 52 ---
\setcounter{footnote}{0}%
\opage{52}Altsaa hører Tankegangen i første Bog ind under den anden Erkendelsesart (ratio),
som vel viser os Tingene under Evighedens Synspunkt, men paa abstrakt Maade
(I, 25). Det Enkeltvæsen, som \textsc{Spinoza} i femte Bog vil vise i dets Sammenhæng
med Substansen, er den menneskelige Aand, som gennem Erkendelse af Tilværelsens
Love staar som Led i den hele Sammenhæng, der gennem den uendelige Forstand
og Tænkningens Attribut har sin sidste Grund i Substansen.

Det er maaske karakteristisk, at nogle af de Exempler, \textsc{Spinoza} i Afhandlingen
om Forstandens Forbedring anfører paa anskuende Viden, mangler i Ethica. Saaledes
det Exempel, at naar jeg erkender noget, véd jeg deraf, hvad det er at erkende.
Dette er en Analyse, eller, om man vil, en umiddelbar Slutning. I den
nævnte Afhandling defineres Metode som Erkendelse af (Bevidsthed om) Erkendelse
(cognitio reflexiva), en Tanke om en Tanke (idea ideæ). (Se ogsaa Ethica II, 20---21,
smlgn. ovenfor p. 46). I Ethica hævdes ogsaa udtrykkelig, at den, der har en sand
Tanke, tillige maa vide, at han har en sand Tanke og ikke kan tvivle om dens
Sandhed, da en Tanke ikke er noget stumt, som et Billede (pictura), man kan stirre
paa, men er en Virksomhed, er selve det at forstaa. Enhver sand Tanke indeholder
Sandhedsnormen i sig (II, 43). Naar \textsc{Spinoza} i Ethica ikke nævner dette som Exempel
paa Intuition, er det altsaa ikke, fordi han ikke anerkender Rigtigheden
deraf. Det er snarere, fordi han mener, at det ikke er den Art Intuition, som han
kalder den tredie og højeste Grad af Erkendelse. Naar den umiddelbare Erkendelse
af, hvad der gaar for sig i en Erkendelsesakt, i Afhandlingen om Forstandens Forbedring
kaldes Intuition, saa er det, hvad jeg har kaldt analytisk Intuition, den
Bevidsthed, der ytrer sig i enhver umiddelbar Slutning, enhver simpel Overgang fra
Led til Led i en Bevisførelse\footnote[1]{\textit{Om Begrebet Intuition} (Vid. Selsk. Oversigt 1914), p. 77---80).}. Derimod lægger \textsc{Spinoza} i Ethica Vægt paa, at en
Enkeltting (res singularis) ved intuitiv Viden erkendes i sin Sammenhæng med Tilværelsen
i det hele. Der er her ikke Tale om simpel Analyse, umiddelbar Tankeovergang,
men om en Helhedsopfattelse, en Konstruktion, der danner Afslutningen
af et langt Tankearbejde og bygger paa den gennem dette Tankearbejde vundne
Loverkendelse. Dette er, hvad jeg har kaldt syntetisk Intuition\footnote[2]{\textit{Om Begrebet Intuition}, p. 80---82.}. Jeg antager, som
ovenfor antydet, at \textsc{Spinoza} i sit Talexempel betragter et Tal (her Sextallet) som et
Individ, der ved en paa Grundlag af Loven for Talrækkens Dannelse vunden Intuition
faar sin Plads bestemt i Forhold til andre Tal. Kun saa passer Exemplet paa
det, som det aabenbart er hans Hensigt at udtrykke. I femte Bog er det selve det
tænkende Menneske, lad os sige \textsc{Spinoza} selv, hvis Personlighed ses som Led i
Verdenssammenhængen, som evigt Aandselement (æternus cogitandi modus) (V, 40
Schol.).

Analytisk Intuition maa allerede være medvirkende ved Overgangen fra det
første Erkendelsestrin til det andet. Den anvendes ved den omhyggelige Sammenligning,
der ikke bliver staaende ved spredt Iagttagelse, men gennem Undersøgelse
af Tingenes Ligheder og Forskelligheder finder det Grundlag for vor Tænken, der
% --- p. 53 ---
\setcounter{footnote}{0}%
\opage{53}bestaar i Begreber om almene Forhold (II, 28---29; 44 Cor.). Den analytiske Intuition
er saaledes et Middel, en Forberedelse. Den syntetiske Intuition derimod, som
\textsc{Spinoza} opstiller i Ethica, er en Helhedsopfattelse, der staar som Erkendelsens sidste
Maal og store Ideal\footnote[1]{\textsc{Joachim} (\textit{A Study of the Ethics of Spinoza}, 1901, p. 183) har kun Ret i, at der allerede forudsættes Intuition ved det andet Erkendelsestrin, naar man derved tænker paa analytisk Intuition. \textsc{Huan} (\textit{Le dieu de Spinoza}, p. 16---19) har klart set dette. --- Hvad jeg kalder analytisk Intuition, beskrives allerede af \textsc{Descartes} (\textit{Regulæ}, Kap. 8), og efter ham af \textsc{Locke} (\textit{Essay} IV, 2, 1 og 10, 3) og \textsc{Hume} (\textit{Treatise} I, 3, 1---3). \textsc{Descartes} bruger lignende Exempler som \textsc{Spinoza} i Afhandlingen om Forstandens Forbedring, og \textsc{Spinoza} kan godt have kendt Regulæ, der i Afskrift vides at være kendt i \textsc{Spinoza}'s Kreds (se herom \textsc{Brunschwicg}: \textit{Les Étapes de la philosophie mathématique}, p. 141).}. Det maa forstaas om den, hvad \textsc{Spinoza} siger om intuitiv
Viden i Afhandlingen om Forstandens Forbedring: „Hvad jeg hidtil har kunnet
forstaa ved denne Art Erkendelse, er meget lidet (perpauca)“. Det er jo ogsaa klart,
at en helt gennemsigtig, rationel Sammenhæng mellem det individuelle og det universelle
stiller en uendelig Opgave for Tanken. \textsc{Leibniz}, der havde saa klart et Blik
for det individuelle og nuancerede og for de utallige Overgange, saa klart dette.
Det er umuligt at kende en Tings fulde Individualitet, da ethvert Enkeltvæsen staar
i en Uendelighed af Vexelvirkninger med alt andet i Verden, og da der er en Uendelighed
af smaa Forskelle, der ikke kendes af os (Nouveaux Essais, III, 3, 6). Det
er dog af stor Interesse, at netop \textsc{Spinoza} har opstillet den syntetiske Intuition som
Ideal, han, som saa ofte er bleven betragtet som Abstraktionernes Mand.

Det matematiske Exempel, \textsc{Spinoza} bruger til at belyse den „syntetiske Intuition“,
viser, at han meget vel anerkender Muligheden af saadan Intuition paa specielle
Omraader, særlig paa det matematiske Omraade. Han vilde altsaa ikke være
uenig med \textsc{Zeuthen}, som hævder, at syntetisk Intuition (en Helhedserkendelse) er
mulig i Matematiken, idet han dog tilføjer, at dette kommer af, at den matematiske
Erkendelse selv er symbolsk. „Den omfatter nemlig i Virkeligheden kun et begrænset
Antal Enkeltheder, som lader sig samle til et Hele, som man dernæst ved
Brug af dertil egnede Symboler kan lade optræde som Enkelthed enten i en videregaaende
matematisk Undersøgelse eller i en af Matematikens Anvendelser“\footnote[2]{\textsc{Zeuthen}: \textit{Om Anvendelse af Regning og Ræsonnement} (Vid. Selsk. Overs. 1914, p. 275).}. Som
jeg har antydet i min Afhandling om Intuition, vil en Helhedserkendelse ogsaa
kunne naas paa specielle Erfaringsomraader, f. Ex. det psykologiske og det historiske.
Men der faar den en afgjort hypotetisk Karakter.

Der er en Forskel mellem den Maade, paa hvilken den intuitive Viden omtales
her i anden Bog og senere i femte Bog. Sagen er, at i femte Bog samvirker
Følelses- og Villieslivet med Tankelivet. Den Erkendelsesglæde, som tilsidst gav
Livet Værdi for \textsc{Spinoza}, naar sit Højdepunkt, idet det viser sig, at det er Tilværelsens
inderste Grund, der gør Erkendelse mulig, ligesom Erkendelsen selv fra sin
Side naar sin Fuldendelse i Tanken om denne inderste Grund og dens Fremtræden
i Attributer og i Modi. Derved faar \textsc{Spinoza}'s syntetiske Intuition et religiøst Præg
og gaar ind under, hvad jeg har kaldt praktisk Intuition. Den faar tillige et kunst%
% --- p. 54 ---
\setcounter{footnote}{0}%
\opage{54}nerisk Præg, da den forudsætter Evne til at se det individuelle i dets hele konkrete
Ejendommelighed og dog ligesom gennemstraalet af Tilværelsens store Love.

Hvad jeg i min nævnte Afhandling har kaldt konkret Intuition, hvorunder
hører de Anskuelsesbilleder, der uvilkaarlig danner sig i Sansning, Erindring og
Fantasi, svarer til \textsc{Spinoza}'s første Erkendelsestrin (imaginatio). Det er paa Grundlag
af den, at analytisk Intuition kan føre til Loverkendelse\footnote[1]{Et interessant Exempel paa en Overgangsform mellem konkret og analytisk Intuition frembyder den indiske Matematiker \textsc{Apastamba}'s Forudsætning, at to Triangler, der er lige store med en og samme tredie, er lige store. Det er Anskuen (Fladeopfattelse), der raader; men det betegner en begyndende Analyse, at hin Forudsætning udskilles fra den hele Synsoplevelse. \textsc{Zeuthen}: \textit{Hvorledes Matematiken i Tiden fra Platon til Euklid blev rationel Videnskab} (Vid. Selsk. Skr. 1917), p. 55 f. --- I det af \textsc{Edgar} \textsc{Rubin} (\textit{Synsoplevede Figurer}. 1915) paaviste Forhold mellem Fladeopfattelse og Konturopfattelse haves et andet smukt Exempel paa konkret Intuitions Overgang til Analyse.} og derigennem muligvis
bane Vej til syntetisk Intuition. For \textsc{Spinoza} var ikke de almene Love Videnskabens
sidste og højeste Opgave. Det sidste Maal var Forstaaelse af det enkelte
og individuelle ved Hjælp af Loverkendelsen\footnote[2]{Se om dette Punkt \textit{Den menneskelige Tanke}, p. 222---227. --- \textit{Totalitet som Kategori}, p. 63---66.}. ---

I epigrammatisk Korthed har \textsc{Spinoza} angivet Hovedtrinnene for menneskelig
Erkendelse, begyndende med Indsamling af Emner, opstigende til Videnskab og afsluttende
med den anskuende Viden, der staar paa Grænsen af Videnskab, Religion
og Kunst, og der tillige --- hvad dog \textsc{Spinoza} neppe vilde indrømme --- i hvert enkelt
Tilfælde maa bære et forskelligt Præg efter Tænkerens individuelle Personlighed
og Livserfaring. Han vilde neppe heller have indrømmet, at et afrundet Helhedsbillede,
som det, der skal betegne Erkendelsens Højdepunkt, videnskabelig set
maa betragtes som hypotetisk.

\subsection*{f. Sætning 48---49.}
\begin{center}\textit{Villiesbegrebet.}\end{center}

\textit{α.} Sin Erkendelsespsykologi og Erkendelsesteori afslutter \textsc{Spinoza} ved paany
at indskærpe, at man ikke maa forvexle Tanke, Billede og Ord. En virkelig Tanke,
der er konsekvent udviklet, er ikke noget hvilende og passivt, men Frugt af Aandsvirksomhed
og selv stadig virkende, udfoldende sine Konsekvenser videre. Naar
man mener, at én kan tro eller indse noget og dog ikke ville det, saa forstaar man
ved Tro eller Indsigt et passivt Forhold til Billede eller Ord („et stumt Billede paa
en Tavle“), ikke en aandelig Virksomhed. Den Virksomhed, ved hvilken en Tanke
udfolder sig i sine Konsekvenser, er ét med en Villen, og omvendt: Villen (volitio)
er Tankekonsekvens. Det er umuligt at have Tanken om en Triangel og ikke
hævde (affirmare), at Vinklernes Sum i det er lig to rette. --- Tillige indskærpes, at
man knn kun have fuld Vished om det, der er fuldt og klart indset og bevist. % PRINTED AS IS: knn (kan)
Vished er ikke det samme som en Fraværelse af Tvivl, der kan skyldes den tilfældige
Omstændighed, at der ikke er indtraadt nogen Anledning til Uro i Sindet.

Ved Ordet Villie (voluntas) forstaar \textsc{Spinoza} i denne Sammenhæng kun Evnen
% --- p. 55 ---
\setcounter{footnote}{0}%
\opage{55}til at bekræfte eller benægte. Bekræftelse er ganske simpelt en Konsekvens af en
Tanke, man har. \textsc{Spinoza} erklærer udtrykkelig, at han ved Villie ikke forstaar
Drift (cupiditas), som ytrer sig i Trang til noget eller Afsky for noget.

Hvad der ligger \textsc{Spinoza} paa Hjerte i hele denne Tankegang, er at indskærpe
Tankelivets aktive Karakter, dets Udspring af Menneskets inderste Væsen, som igen
er en enkelt Modus af den i alt virkende Substans. Det er i Selvvirksomhedens
Begreb, at for ham Forstand og Villie er ét. Hans Standpunkt her minder om
\textsc{Sokrates}'. Ogsaa \textsc{Sokrates} hævdede jo, at den, der véd, ikke kan andet end
handle efter sin Viden; gør kan det ikke, er det for \textsc{Sokrates} et Bevis for, at han % PRINTED AS IS: kan (han)
i Virkeligheden ingen Viden havde.

Naar \textsc{Descartes} lærte, at Villien kan omfatte mere end Forstanden, indrømmer
\textsc{Spinoza} kun dette, forsaavidt man ved Forstand mener Evne til klar og tydelig
Tænken. Han indrømmer naturligvis, at der gives dunkle Tanker; men de beror
kun paa Ufuldstændighed. Og han indrømmer, at der gives Vildfarelse; men den
er noget negativt, beroende paa Fraværelse af Tilfælde, som strider mod, hvad man
foreløbig har slaaet sig til Ro med.

Naar \textsc{Spinoza} erklærer Villen og Tænken for ét, tænker han paa saadanne
Karakterer, hos hvem Overvejelsen formaar at give en klar Forestilling om de forskellige
Muligheder, saa at Beslutningen bliver en simpel Konsekvens af Forholdet
mellem disse Muligheder, --- og hos hvem saa dernæst den Tanke, om hvilken
Sindet i Beslutningen har koncentreret sig, fastholdes i sin Konsekvens, indtil den
er virkeliggjort. Men selv saadanne Karakterer har gennemgaaet en Udvikling, der
har ført gennem Dunkelheder og Vildfarelser, og hvis det, som menes med Ordet
Villie, ikke er noget, der først bliver til paa Fuldendelsens Trin, maa der være
andre Villiesformer end den, som er ét med klar Tænken, og som meget vel kan
kaldes aktive, da de kan have deres egentlige Aarsag i Menneskets eget Væsen.
Saadanne Villiesformer er det, der betegnes ved Ord som Trang, Stræben, Attraa,
Drift, Begær. De kan ofte staa i Modsætning til den paa klar Eftertanke grundede
Villen --- som naar \textsc{Goethe} siger om sin Ven:

\begin{quote}
Noch ist, bei tiefer Neigung für das Wahre,\\
ihm Irrthum eine Leidenschaft.
\end{quote}

I Indrømmelsen af, at Villien under visse Forhold og i en vis Betydning kan
strække sig videre end Forstanden, ligger Antydningen af et fyldigere Villiesbegreb
end det, der paa dette Sted hævdes af \textsc{Spinoza}. Ensidigheden ved dette er, at
Dunkelhed og Vildfarelse kun opfattes som noget negativt, skønt de er reale, positive
Tilstande. Fraværelse af Tvivl betyder uvilkaarlig Tro, Tillid eller Tilslutning
og udelukker ikke Klarhed og fuld Konsekvens indenfor den Forestillingskreds, der
staar til Raadighed. At \textsc{Spinoza} ikke har set saadanne Tilstandes positive Karakter,
hænger sammen med, at han ikke skelner mellem to Slags positive Domme,
--- primære, der bliver til ved umiddelbar Besindelse (Analyse) overfor givne Emner,
og sekundære, der afviser en Nægtelse eller en Tvivl angaaende saadanne Emner.
% --- p. 56 ---
\setcounter{footnote}{0}%
\opage{56}Han miskender særlig den umiddelbare Anskuens, særlig Fantasiens Betydning og
Selvstændighed, naar han paastaar (49 Schol.), at det at forestille sig en vinget
Hest ikke er andet end at udsige Bevingethed om Hesten\footnote[1]{Se om de her berørte Forhold \textit{Den menneskelige Tanke}, p. 53---57; 78 f. --- \textit{Totalitet som Kategori}, p. 8.}.

Endelig kan Identitet af Forstand og Villie i og for sig ligesaa godt betyde
Reduktion af Forstand til Villie, som omvendt. Hvad det er \textsc{Spinoza} om at gøre,
er Hævdelse af den aandelige Aktivitet. Til denne Side af Sagen vil jeg senere (ved
III, 9) komme tilbage.

\textit{β.} Medens \textsc{Spinoza} her har gjort Villen til ét med Tænken, opfatter han i
tredie Bog (9 med Schol.) Villen som Udtryk for en Stræben efter Selvopholdelse
(conatus perseverandi in suo esse), der, naar den mærkes, kaldes Trang (appetitus),
og naar den er forbunden med Forestillingen om sit Maal, kaldes Drift (cupiditas).
Og i de Definitioner, der danner Tillægget til tredie Bog, siges Drift (cupiditas) at
være selve Menneskets Væsen, forsaavidt det ud fra en indre Tilstand bestemmes til at
foretage en vis Handling. Endelig gives Ordet Drift en saa udstrakt Betydning, at
det kan betegne al menneskelig Stræben (conatus), Tilskyndelse (impetus), Trang
(appetitus) og Villen (volitio).

Det er aabenbart to forskellige Villiesbegreber, der opstilles i Slutningen af
anden Bog og i Begyndelsen af tredie Bog, altsaa i stor Nærhed af hinanden.

Man har paa forskellig Maade søgt at forklare denne Forskel. \textsc{Sigwart} formoder
(i sin Oversættelse af den korte Traktat p. 208), at Selvopholdelsestrangen
lægger sig for Dagen i selve de Bekræftelser og Benægtelser, som omtales i II, 49.
Han lægger altsaa Fremstillingen i tredie Bog til Grund som den mere fundamentale.
Men han indrømmer, at \textsc{Spinoza} selv ikke siger dette med rene Ord, saa nær det
maatte ligge ham. \textsc{Baensch} gaar (i sin Oversættelse af Ethica p. XVI) den modsatte
Vej, idet han i tredie Bog netop finder Højdepunktet af den Intellektualisme,
som allerede gør sig gældende i anden Bog: Følelse og Trang, der ifølge II, Ax. 3
blot skulde være afhængige af Forestillingerne, gøres i tredie Bog til ét med disse,
og Selvopholdelsestrangen skal saa ikke være andet end Selvbekræftelsen af den
Tanke (idea), der udgør Menneskets Væsen. Men \textsc{Baensch} maa indrømme, at Intellektualismen
saa dog paa sit Højdepunkt slaar over i sin Modsætning, idet det
(III, 9 Schol.) hævdes, at vi ikke tragter efter noget, ikke vil det, fordi det er godt,
men omvendt betragter det som godt, fordi vi vil det. Her sættes jo udtrykkelig
Trang og Villen som mere fundamentale end Erkendelsen. \textsc{Tönnies} har (Studien
zur Entwickelungsgeschichte des \textsc{Spinoza}. Vierteljahrsschrift für wissenschaftliche
Philosophie VI) søgt at løse Vanskeligheden ved den Antagelse, at der har været et
Tidsmellemrum mellem anden og tredie Bog i deres nuværende Form, og at tredie
Bog først er afsluttet, efterat \textsc{Hobbes}' Psykologi har øvet Indflydelse paa \textsc{Spinoza}.
Hvad der taler for denne Forklaring, er, at \textsc{Hobbes} netop opfatter Villien som Trang
eller Drift. I Elements of Law (1640) kaldes (I, 12, 2) den sidste Drift, der ytrer sig
under Overvejelsen, Villen. Trang, Frygt og Haab kan ikke være villede (voluntary)

% --- p. 57 ---
\setcounter{footnote}{0}%
\opage{57}because they are the will. De samme Begrebsbestemmelser findes i De Homine
(1658) Kap. 11 og i Leviathan (1651). Deres historiske Forudsætninger vil vi komme
til at omtale ved tredie Bog. --- Der er meget, som taler for \textsc{Tönnies}' Hypotese.
Men det er og bliver dog paafaldende, at \textsc{Spinoza} ved den endelige Sammenarbejden
af sine Udkast ikke er bleven opmærksom paa Uoverensstemmelsen. Forklaringen
kunde ligge i, at han uvilkaarlig har identificeret den Trang og Drift, der ikke kan
gøres til ét med ren Tænken, med den Villen, som ledes af dunkle eller falske Forestillinger.
I sin Indrømmelse overfor \textsc{Descartes} har han jo, som allerede bemærket,
antydet en Udvidelse af det Villiesbegreb, han opstiller II, 49. Men dertil kommer,
hvad Drøftelsen af de tre sidste Bøger vil vise, at der gennem hele hans psykologiske
og etiske Opfattelse gaar en Brydning mellem en intellektualistisk og en
voluntaristisk Opfattelse.

\textit{γ.} I en Slutningsbemærkning afviser \textsc{Spinoza} først Antagelsen af Villiens Indifferens
(„den frie Villie“), idet han dels hævder, at en Forestilling ikke er et hvilende
Billede, men en Tankevirksomhed, der ikke kan forliges med Indifferens, dels
direkte afviser Muligheden af absolute Ligevægtstilstande. Et Menneske, der er lige
sulten og lige tørstig og derfor ikke kan bestemme sig til enten at spise eller at
drikke, vilde \textsc{Spinoza} --- med Hentydning til \textsc{Buridan}'s berømte Æsel\footnote[1]{Hvad dette Dyr angaar, har allerede \textsc{Schopenhauer} forgæves søgt efter det i de Skrifter, der haves af \textsc{Buridan} (en Skolastiker fra første Del af det 14. Aarhundrede). Senere har \textsc{Siebeck} (\textit{Beiträge zur Entwickelungsgeschichte der neueren Psychologie}. Giessen 1891) været paa Jagt efter det uden Held. \textsc{Siebeck} formoder, at \textsc{Buridan} i mundtlig Fremstilling har brugt Exemplet til at vise Forskellen mellem Menneske og Dyr: Mennesket vil kunne finde en Differens, hvor Dyret ingen finder. --- \textsc{Spinoza}'s Tankegang overfor Æslet gaar netop i denne Retning. --- Som \textsc{Schopenhauer} paaviser (\textit{Die beiden Grundprobleme der Ethik\textsuperscript{2}}, p. 58 f) er Exemplet ældre end \textsc{Buridan}. Det forekommer allerede hos \textsc{Dante} og antydes hos \textsc{Aristoteles}. Det nye hos \textsc{Buridan} er, at det er et Æsel, der klemmes mellem to ligestore Muligheder.} --- snarere
kalde et Æsel end et Menneske.

Tilsidst viser \textsc{Spinoza}, hvorledes Overbevisningen om alle Tings Nødvendighed
lærer os at stræbe efter at naa den højest mulige Forstaaelse af Tilværelsen, fører
os til Ro og Resignation overfor Skæbnens Svingninger, indskærper Medfølelse og
giver Anvisning til den rette Statsstyrelse.

% --- p. 58 ---
\setcounter{footnote}{0}%
\opage{58}\section*{Tredie Bog.}
\subsection*{a. Fortale, Definitioner, Axiomer og Postulater.}

I Fortalen udtales i kraftige Ord Overbevisningen om, at der paa det aandelige
Omraade raader bestemte Love, ligesom paa det materielle Omraade. Naturen
er stedse den samme, dens Magt og Virkeevne overalt ens; det er de samme Love,
efter hvilke alt sker og alle Omsætninger fra en Form til en anden foregaar. Derfor
er det meningsløst at laste Naturen, og det er ogsaa meningsløst at udlede
menneskelig Svaghed og Ustadighed af Syndighed eller le eller græde eller harmes
over Menneskenaturen, som den nu engang er. Hvad det gælder om, er at lære
de menneskelige Følelsers Natur at kende og paa Grundlag af denne Erkendelse at
vise, hvad Sindet selv kan gøre for at harmonisere (moderari) dem, selv om det
ikke, som \textsc{Descartes} mente, har absolut Herredømme over dem. Men bortset herfra
er Forstaaelse af menneskelige Følelser ligesaa værdifuld som Forstaaelse af den
materielle Naturs Fænomener, og derfor vil \textsc{Spinoza} her anvende samme Metode
som ellers og betragte menneskelige Handlinger og Tilskyndelser, som om det gjaldt
Linier, Flader eller Legemer.

I den anden Definition omtales den vigtige, allerede omtalte Bestemmelse af
Begreberne Aktivitet og Passivitet, der, som det vil vise sig, er af afgørende Betydning
for Forstaaelsen af Følelsernes Udvikling. Den finder strax Anvendelse ved
selve Definitionen af, hvad Følelse er. Idet \textsc{Spinoza} her tager sit Udgangspunkt fra
den legemlige Side, definerer han Følelse som en legemlig Tilstand, ved hvilken
Legemets Virkekraft (agendi potentia) forøges eller formindskes, og til hvilken der
svarer en vis Forestilling (idea).

Med den systematisk-deduktive Fremstilling af Følelsernes Psykologi, som nu
begynder, maa stadig sammenholdes den mere rent beskrivende Fremstilling, som
gives i Tillægget til tredie Bog, og som ofte er meget oplysende, skønt den ikke
altid stemmer med Fremstillingen i selve Bogen. Det beskrivende Tillæg ender
f. Ex. med en „almindelig Definition af Følelser“, hvorved disse opfattes som passive
Tilstande og gøres til ét med „dunkle Forestillinger“ (confusæ ideæ). Her
tages altsaa ikke Hensyn til, at Følelser ifølge Def. 3 kan være baade aktive og
passive, alt efter som de har deres Udspring i selve Menneskets Natur eller skyldes
ydre Aarsager. Og der tages ikke Hensyn til, at der i Sætning 58---59 udtrykkelig
vises, hvorledes en Følelse kan faa en aktiv Karakter.

% --- p. 59 ---
\setcounter{footnote}{0}%
\opage{59}Baade i Def. 3 og i Tillægget gøres Følelse til ét med Forestilling (idea). Der
ligger her, ligesom ved Villiesbegrebet i Slutningen af anden Bog, en intellektualistisk
Opfattelse til Grund. Det er først i det Parti af Bogen, der begynder med Sætning 4,
at en mere realistisk og voluntaristisk Tankegang gør sig gældende.

Begrebet Forestilling (idea) har i det hele en skæbnesvanger Indflydelse hos
\textsc{Spinoza}. Derved, at Sjælen fra først af defineres som den til Legemet svarende
Forestilling (idea corporis), muliggøres den Sammenblanding af Erkendelsesteori og
Psykologi, som er paavist i det foregaaende. Og formedelst den aktive Karakter af
Forestillingerne, som \textsc{Spinoza} med Rette betoner, hævdes der (II, 48---49) en absolut
Identitet mellem Erkendelse og Villie, skønt der kun er partiel Identitet mellem dem,
da der gives anden sjælelig Selvvirksomhed end den, der fremtræder som Erkendelse.
Hvad særlig Følelsernes Psykologi angaar, bliver det overset, at der er Forskel
mellem fornemmelsesbetinget og forestillingsbetinget Følelse (elementære og ideelle
Følelser), --- og at selve de Følelser, der allermest er betingede ved Forestillinger,
lige saa lidt er absolut identiske med disse, som Villiesytringerne er det.

\subsection*{b. Sætning 1---3.}
\begin{center}\textit{Aandelig Aktivitet og Passivitet.}\end{center}

Vi forholder os passive, naar vi kun har ufuldstændige Forestillinger (ideæ
inadæqvatæ), d. v. s. saadanne, der ikke kan forstaas ud fra vor egen Natur alene,
men forudsætter Paavirkning udefra. Men for ikke at blive misforstaaet, indskærper
\textsc{Spinoza} meget energisk og tydelig den Opfattelse af Forholdet mellem Sjæl og Legeme,
der efter hans Mening følger af det i de to første Bøger hævdede Forhold
mellem de forskellige Attributer: III, 2 Schol. begrundes uden videre ud fra II, 7.
En sjælelig og en legemlig Tilstand eller Virksomhed, f. Ex. en Beslutning og den dertil
svarende legemlige Tilstand, er en og samme Ting og gør sig samtidig (simul
natura) gældende. Der maa i Kraft heraf kræves baade en psykologisk og en fysiologisk
Forklaring af de passive eller aktive Tilstande, Mennesker kan være i, og af
de Forandringer, de kan undergaa.

Ved i Konsekvens af sin Attributlære at hævde Samtidigheden af de sjælelige
og de tilsvarende legemlige Tilstande, kommer \textsc{Spinoza}, uden at være sig det bevidst,
indirekte til at pege paa den eneste mulige Maade, paa hvilken Problemet
om Forholdet mellem Sjæl og Legeme kunde tænkes at faa en virkelig experimental
Afgørelse. Dersom det nemlig kunde godtgøres, at en sjælelig Tilstand, en Sansefornemmelse
f. Ex., først indtraadte, efter at den tilsvarende legemlige Tilstand
var opstaaet, --- eller at (ved Handling) en Bevægelse først indlededes i Hjernen,
efter at en Beslutning var taget, vilde den populære Vexelvirkningslære være bevist.
Vi vilde da i det ene Tilfælde faa en Række af materielle Tilstande, der paa
et vist Punkt afbrødes og fortsattes af en kortere eller længere Række af sjælelige
Tilstande, --- og i det andet Tilfælde en Række af sjælelige Tilstande, der paa et vist
Punkt afbrødes og fortsattes af en Række af materielle Tilstande. Det er nu ikke
% --- p. 60 ---
\setcounter{footnote}{0}%
\opage{60}let at se, hvorledes man skulde komme til at foretage et saadant Experiment. En
anden Sag er det, at interessante Forsøg har godtgjort de materielle Processers
kontinuerlige Forløb (i som udenfor Organismen) og deres Overensstemmelse med,
hvad man nu kalder Loven om Energiens Bestaaen, ogsaa hvor sjælelige Processer
samtidig kunde konstateres, og at man fremdeles har paavist en Proportionalitet
mellem det organiske og det psykiske. Men skønt \textsc{Spinoza} ikke paa disse Punkter
kunde vide, hvad vi nu véd eller mener at vide, er det dog netop, hvad Anmærkningen
til III, 2 viser, de materielle Fænomeners Kontinuitet og Proportionaliteten
mellem det psykiske og det organiske, han her beraaber sig paa. Og den almindelige
Attributlære, han tilsidst beraaber sig paa, er, som jeg har søgt at vise, selv fremgaaet
som Konsekvens af den „nye Videnskabs“ Metode. Han beraaber sig udtrykkelig
(i III, 2 Schol.) paa Erfaringsvidenskabens Metode. Man har, siger han
her, endnu ikke paavist nogen Grænse for, hvad der indenfor Organismen kan forklares
efter den legemlige Naturs Love, og at beraabe sig paa Sjælens Indgriben til
Forklaring af legemlige Processer er kun at indrømme sin Uvidenhed. Dette Ræsonnement
er det samme, som han andensteds (i Slutningsbemærkningen til første
Bog) anvender overfor den teleologiske Naturforklaring og (i Ep. 75) overfor Mirakeltroen.
Spiritualisme, Teologi og Mirakeltro hviler for ham paa samme Princip.
Han betragter i denne Sammenhæng sin Attributlære (Identitetshypotesen) som en
Arbejdshypotese, ligesom Tilfældet ogsaa mere og mere er hos moderne Tilhængere
af denne Lære. I Grunden kan der ikke anvendes anden Arbejdshypotese ved
Behandlingen af Forholdet mellem Sjæl og Legeme end Attributhypotesen, der
bygger paa Muligheden af at slutte fra den ene Kreds af Fænomener til den anden,
skønt et Aarsagsforhold (der tillige er et Tidsforhold) ikke antages. --- Udtryk som
Parallelisme eller Duplicisme, der saa ofte benyttes om den spinozistiske Hypotese,
skyldes en vildledende Analogi\footnote[1]{Derimod passer de nævnte Udtryk (og den Polemik, de plejer at give Anledning til) paa \textsc{Chr}. \textsc{Wolff}'s Lære. Smlgn. \textit{Den nyere Filosofis Historie\textsuperscript{2}}, I, p. 524 (Note 76).}. Søger man en Analogi, der virkelig passer, kan
man benytte \textsc{Fechner}'s Billede med Forholdet mellem den konvexe og den konkave
Side af en og samme Kurve, eller \textsc{Brunschwicg}'s Analogi (som ifølge ham er mere
end Analogi) med Forholdet mellem Algebra og Geometri i den analytiske Geometri.

\subsection*{c. Sætning 4---12.}
\begin{center}\textit{Selvopholdelsestrangen.}\end{center}

\textit{α.} Som Indledning til Følelsernes Psykologi fremsætter \textsc{Spinoza} den Lære, at
ethvert Væsen har en Tendens til at bevare og opholde sig selv. Hvad der skal
kunne ophæve en Tings Existens, maa være noget andet end den selv. Begrundelsen
af denne Sætning søges i, at ethvert Enkeltvæsen (Modus) er en særlig Form
(Maade) for Substansens Væsen og Virken og derfor paa sin Vis har Del i Substansens
Evne til at bestaa ved sig selv og forstaas ved sig selv. Selvopholdelsestendensen
er ikke en underordnet Egenskab, men er ét med det enkelte Væsens inderste Natur.
% --- p. 61 ---
\setcounter{footnote}{0}%
\opage{61}\textsc{Spinoza} henviser her til den tidligere (I, 36) opstillede Sætning, at der intet gives,
af hvis Natur der ikke udspringer Virkninger, en Sætning, som netop begrundes
ved, at alt, hvad der existerer, udtrykker Guds Væsen paa sin bestemte Maade.
Det er den positive Side af Forholdet mellem Substans og Modus, \textsc{Spinoza} her
drager frem. Selve det positive (qvid positivum), hvorved Tingene bestemmes til
at virke, udspringer af Gud (I, 26 Dem.). For det meste fremhæver \textsc{Spinoza} den
negative Side af hint Forhold, i Henhold til den Forudsætning, at enhver Begrænsning,
enhver nærmere Bestemmelse er en Negation (omnis determinatio negatio,
Ep. 50). Vi kender kun Enkeltvæsener af Erfaring, men naar de nu engang existerer,
maa Substansen lægge sig for Dagen i dem, og derfor er deres Evne eller
Stræben (potentia sive conatus) efter Selvopholdelse ét med deres givne, aktuelle
Natur (7 Dem.).

Alligevel optræder dette Begreb ret uforberedt hos \textsc{Spinoza}. Han gør ikke,
som \textsc{Galilei}, \textsc{Gassendi}, \textsc{Hobbes} og senere \textsc{Leibniz}, nogen naturfilosofisk Brug af
dette Begreb. En saadan Brug vilde have kunnet føre ham ud over den kartesianske
Bevægelseslære, hvis Uholdbarhed han (i sit sidste Brev til \textsc{Tschirnhausen})
indrømmede. I de fysiske Læresætninger i anden Bog giver han en rent mekanisk
Bevægelseslære i kartesiansk Aand.

Det er aabenbart ad rent psykologisk Vej, han er ført til at give Begrebet
Selvopholdelse en grundlæggende Betydning i Aandsvidenskaben, og her har han
\textsc{Bruno}\footnote[1]{Af Interesse for Forholdet mellem \textsc{Spinoza} og \textsc{Bruno} er det, at den korte Traktat (I, 5) betragter den enkelte Tings Stræben efter Selvopholdelse som Udtryk for et specielt Forsyn. En ganske lignende Tanke findes hos \textsc{Bruno}. (\textit{Den nyere Filosofis Historie\textsuperscript{2}}, I, p. 128).} og \textsc{Hobbes}, og gennem dem \textsc{Telesio} til Forgængere. Gennem \textsc{Telesio}
fører, som \textsc{Dilthey}\footnote[2]{\textit{Weltanschauung und Analyse des Menschen seit Renaissance und Reformation}. (1914). p. 433. 463.} har paavist, hele denne Tankegang tilbage til Stoicismen.

\textit{β.} Hvad Mennesket angaar, ytrer Selvopholdelsestrangen sig, baade forsaavidt
det har klare og tydelige Forestillinger, og forsaavidt det har dunkle Forestillinger.
Derfor foretages der nu en Udvidelse af det i II, 48---49 opstillede Villiesbegreb.
Selvopholdelsestrangen kaldes Villie (voluntas), forsaavidt den blot henføres til Sjælen,
Trang (appetitus), forsaavidt den henføres baade til Legemet og Sjælen, og Drift
(cupiditas), forsaavidt vi bliver os Trangen [og dens Genstand] bevidst. Hvad vi
anser for godt, beror paa Selvopholdelsestrangen, altsaa paa Villie (i videste Betydning
af Ordet). Vi vil derfor ikke noget, fordi det er godt, men vi anser noget for
godt, fordi vi vil det (9 Schol. og 39 Schol.). Vi kan ville det, vi i Grunden ikke
vil, og vende os bort fra (nolle), hvad vi i Grunden vil, naar f. Ex. Frygt for et
større Onde faar os til at ville et mindre Onde eller opgive, hvad der ellers vilde
være et Gode (39 Schol.).

Uden igen at komme ind paa Forholdet mellem det her fremtrædende Villiesbegreb
og det i Slutningen af anden Bog opstillede, skal jeg kun bemærke, at selv
om, som det hlev udtalt i II, 49, Forstand og Villie var absolut ét og det samme, % PRINTED AS IS: hlev
saa vilde man med samme Ret kunne sige, at vi erkender det gode, fordi vi vil det,
% --- p. 62 ---
\setcounter{footnote}{0}%
\opage{62}og at vi vil det, fordi vi erkender det. Naar A = B, gælder det jo baade at A er B,
og at B er A. Naar vi ser paa Villieslivets Udvikling, kan den første Sætning gælde
paa mere primitive Trin, den sidste paa højere Trin: vi begynder med at anse for
godt, hvad vi vil (har Lyst og Trang til), og vi kan ende med at ville, hvad vi har
lært at betragte som et Gode. Den mest fundamentale Sætning er og bliver dog
den, at vor Villen er Udtryk for vor Trang. Denne Sætning har sin Gyldighed helt
op i Aandslivets højeste Regioner. Vi vil kun Videnskab og Kunst, fordi vi har
Trang til dem. Kun ved denne Betragtning bliver det muligt at bygge Bro mellem
\textsc{Spinoza}'s to Villiesbegreber. Men en saadan Betragtning har \textsc{Spinoza} ikke anstillet,
idet det genetiske Synspunkt ikke fuldstændig kommer til sin Ret hos ham.

Den betydningsfulde Forskel mellem de to Villiesbegreber ses af IV, 7---8 og 14,
hvor det indskærpes, at Erkendelsen af godt og ondt kun er en Magt i os, fordi
den er en Følelse (affectus), ikke fordi den er sand Erkendelse. Der er altsaa en
Følelse eller Trang med i Spillet, overalt hvor Erkendelsen udfører et indre Arbejde
(ikke blot selv er et Arbejde).

\textit{γ.} Forskellen mellem Trang (appetitus) og Drift (cupiditas) bestaar i, at vi i
Driften er os bevidst, hvad vi stræber efter, medens den blotte Trang er blind (III,
9 Schol.). Denne Forskel tillægger \textsc{Spinoza} dog ikke afgørende Betydning. I Tillægget
til tredie Bog (Kap. 1 Expl.) siger han, at hvad enten Mennesket er sig sin
Trang bevidst eller ikke, er det dog den samme Trang, og han kan derfor bruge
Ordet Drift (cupiditas) til at betegne al Stræben (conatus), Trang (appetitus), Tilskyndelse
(impetus) og Villen (volitiones).

Det er et stort Spørgsmaal, om \textsc{Spinoza} har Ret heri. Bevidstheden om Maalet
for vor Trang kan blive af afgørende Betydning. Naar Trangen ikke længere er
blind, kan Retningen blive bestemtere; der kan søges Midler, som ikke umiddelbart
frembyder sig; der kan vækkes Opmærksomhed for Hindringer. Er blot en enkelt
Forestilling, den om Maalet for Trangen, givet, kan den fremkalde andre Forestillinger,
endog meget fjerntliggende, og Handlingen kan saaledes bestemmes ved en
større Kreds af Erfaringer. Det er Indledningen til, hvad vi rent psykologisk kalder
en højere Udvikling af Villieslivet, at Bevidstheden om Maal og Midler klares. Med
Blindheden følger let Ubændighed og Uro, medens Klarhed over Maal og Midler
giver den Begrænsning, der gør Mesterskabet muligt. Denne Udvikling fra Dunkelhed
til Klarhed gaar ikke for sig uden Fare. Den umiddelbare Trang, Reflexbevægelsens
og Instinktets uvilkaarlige Higen kan svækkes ved Eftertanke og Spekulation.
Men netop denne Fare viser, at Forholdet mellem bevidst og ubevidst Trang ikke
er saa ligegyldig, som den anførte spinozistiske Sætning udtaler. Og \textsc{Spinoza}'s egen
glimrende Fremstilling af Erkendelsens Betydning for Følelseslivet i de følgende Afsnit
af tredie Bog maa ogsaa gælde for Villieslivet, da, som vi vil faa at se, alle
Følelser af Lyst og Ulyst for \textsc{Spinoza} er Symptomer paa Trang\footnote[1]{Om det snævrere og det videre Villiesbegreb og om Betydningen af Bevidsthedsgrader for Villieslivet se min Afhandling \textit{Begrebet Villie} (i tredie Række af Mindre Arbejder). --- Medens \textsc{Schopenhauer} udtaler sig i lignende Retning som \textsc{Spinoza}, antager \textsc{Fichte} en aldeles afgørende Forskel mellem ubevidst og bevidst Trang. --- I nyeste Tid har \textsc{Henri} \textsc{Bergson} betragtet Instinktets Afløsning af Eftertanken som en Art Syndefald. Smlgn. \textit{Henri Bergson's Filosofi}, p. 16 ff; 39 f.}.

% --- p. 63 ---
\setcounter{footnote}{0}%
\opage{63}I sit interessante Skrift „Psychologie des emotionalen Denkens“ (1908) p. 566---
574 nægter \textsc{Heinrich} \textsc{Maier}, at der gives forestilingsløse Villiesytringer. Han betragter % PRINTED AS IS: forestilingsløse
Villie som Sjælelivets inderste Væsen, men spontane og reflexe Bevægelser
antager han slet ikke for psykiske, og i Instinktet er der, paastaar han, altid Forestilling
om Maalet tilstede. Det er nu ogsaa saa, at spontane og reflexe Bevægelser
kan foregaa ubevidst. Men der kan være forbundet Bevidsthedsytringer med dem,
dels idet de mærkes, medens de udføres, dels ogsaa saaledes, at der føles Lyst, naar
de faar frit Afløb, og Ulyst, naar de hæmmes, dels endog saaledes at der mærkes
en Uro, svingende mellem Lyst og Ulyst, inden de kan udføres. I alt dette kundgør
der sig en Tendens eller en Trang, som kun er forstaaelig formedelst Vigtigheden
for Selvopholdelsen af de Bevægelser, den fører til. Det er dette, \textsc{Spinoza}
har for Øje ved sit Begreb om Trang (appetitus). --- Hvad det andet Punkt angaar,
er det en vidtgaaende Antropomorfisme, det vil her sige en Tydning ud fra det
voksne og reflekterende Menneskes Standpunkt, naar \textsc{Maier} lægger en Formaalsforestilling
ind i enhver Instinkthandling. Vel føjer han til, at denne Formaalsforestilling
kun optræder som Led i et Komplex, ikke analyseres ud af dette. Men netop
denne Forskel kan være af stor Betydning. Et psykisk Element kan ikke uden
videre siges at være det samme, før og efter at det er blevet analyseret ud af det
Komplex, i hvilket det først optraadte. Det ikke-differentierede nærmer sig til det
ubevidste, hvad de enkelte Elementer angaar. Der er i hvert Tilfælde en hel Række
af Overgange. Og det er altid betænkeligt uden videre at springe dem over og i
de endnu udifferentierede Tilstande at lægge ind, hvad man har fundet i de mere
differentierede. Forskellen er kvalitativ.


Efter denne psykologiske Exkurs vender vi tilbage til \textsc{Spinoza}.


\textit{δ.} Af Selvopholdelsestrangen følger (11---13) en Trang til at fastholde eller
foretrække, hvad der forøger Virkekraften i aandelig og legemlig Henseende. Der
kan, med Hensyn til Virkekraft eller Existenskraft (potentia agendi --- vis existendi),
der er ét med Fuldkommenhed eller Realitet (11 med Schol. og Slutningsstykket af
tredie Bog), snart erfares en Overgang til større, snart til mindre Aktivitet eller
Fuldkommenhed. I første Tilfælde opstaar Glæde (Lyst) (lætitia), i andet Tilfælde
Tilfælde Sorg (Ulyst) (tristitia). I Tillægget til tredie Bog (Kap. 2 og 3) betoner \textsc{Spinoza} % PRINTED AS IS: Tilfælde doubled at the line turn
meget stærkt, at Glæde og Sorg bestaar i selve Overgangene (actus transeundi)
fra mindre til større, eller fra større til mindre Fuldkommenhed, --- ikke i selve
den Fuldkommenhed eller Ufuldkommenhed, hvortil der gaas over. Det afgørende
er altsaa Forholdet til den foregaaende Tilstand. Man maa jo besidde en vis Fuldkommenhed
(Handle- eller Existenskraft) for at kunne gaa over til mindre Fuldkommenhed,
og omvendt maa der være en vis Ufuldkommenhed tilstede, for at der
kan gaas over til større Fuldkommenhed. Hvad \textsc{Spinoza} her indskærper, er, hvad
jeg i min Psykologi kalder „Forholdsloven paa Følelsernes Omraade“. I Slutningsstykket
til tredie Bog indskærpes paany, at det er selve Overgangen, det kommer
% --- p. 64 ---
\setcounter{footnote}{0}%
\opage{64}an paa. Men her er det ikke Forskellen mellem Overgang og Tilstand, han vil
indpræge. Han hævder derimod her (noget, hvori han gør flere moderne Psykologer
til Skamme), at der ikke behøver at finde en udtrykkelig Sammenligning Sted
mellem den foregaaende og den følgende Tilstand; det er et aldeles umiddelbart
Forhold mellem to Tilstande, der betinger Glæden og Sorgen. \textsc{Spinoza}'s Finhed
som psykologisk Iagttager træder maaske intetsteds mere frem end ved disse Bemærkninger.


Ved Siden af Glæde og Sorg stilles Driften som tredie Grundfølelse (affectus
primarius). Af disse tre kan alle andre Følelser afledes (11 Schol.). Senere hen
(57 Dem.) siges Glæden og Sorgen at være ét med Trang eller Drift, forsaavidt
denne fremmes eller hæmmes ved ydre Aarsager. Og dette er konsekvent indenfor
\textsc{Spinoza}'s hele Filosofi, da Mennesket, som ethvert Enkeltvæsen, staar som en særlig
Form for Substansen, repræsenterer ligesom en Del af dennes Magt: det kommer
da an paa, hvorledes Forholdet er mellem den indefra bestemte Virken og de ydre
Forhold, d. v. s. mellem den enkelte Modus og de andre Modi, der ligesom den er
særlige Former for Substansens Magt.

Naar dette Forhold er saaledes, at Selvhævdelsen begunstiges derved, er der
ifølge \textsc{Spinoza} en Tendens til at dvæle i Forestillingen ved alt, hvad der kan opretholde
og udvikle det (12). Og den samme Tendens viser sig ogsaa, naar Mennesket
føler Glæde ved og Kærlighed til et andet Væsen; da vil Forestillingerne om, hvad
der tjener til dette Væsens Bevarelse og Udvikling, fastholdes (25). --- Dette er en
Side af det (af \textsc{Hume}, \textsc{Lotze} o. fl. beskrevne) Fænomen, som kan kaldes Følelsens
Expansion. Men \textsc{Spinoza} omtaler kun Fænomenet, forsaavidt det er Udtryk for
Selvhævdelse. Han viser (26 Schol.), hvorledes det kan føre til Overvurdering af
sig selv og af dem, man elsker, og til Undervurdering af andre. Men Erfaring viser,
at saadan Expansion ogsaa kan finde Sted ved Følelser (som Frygt og Medlidenhed),
der ikke fremmer Selvhævdelse. Forholdet er ikke saa simpelt, som \textsc{Spinoza} mente.
Sjælelivets enkelte Elementer kan, i Kraft af en psykisk Inertilov, have Tendens til
at brede sig paa Bekostning af Sjælelivet i dets Helhed. Her er mere Mulighed for
Disharmoni, end \textsc{Spinoza} antog\footnote[1]{13 Cor. siger \textsc{Spinoza}, at Sjælen har Uvillie mod at forestille sig det, der formindsker eller hæmmer dens og Legemets Kraft. --- Men formedelst Tendensen til Expansion kan der opstaa en forunderlig Trang til netop at dvæle ved det hæmmende (fordybe sig i Sorg f. Ex.).}.

Inden vi gaar over til \textsc{Spinoza}'s Fremstilling af de specielle Følelser og Drifter,
trænger den Bemærkning sig frem, at Uhensigtsmæssigheden af \textsc{Spinoza}'s „geometriske
Metode“ intetsteds viser sig saa tydelig som i det nu følgende Afsnit. De forskellige
Fænomener kommer til alle at staa paa Linie med hverandre, skønt de falder i
forskellige Grupper, og skønt der indenfor hver enkelt Gruppe er Forskel mellem
fundamentale og afledede Fænomener. Dog antyder \textsc{Spinoza} lejlighedsvis (især i
Tillægget) visse Inddelingspunkter, og i Henhold til dem kan følgende Oversigt gives
over Følelsernes Psykologi hos \textsc{Spinoza}:
% --- p. 65 ---
\setcounter{footnote}{0}%
\opage{65}I. Passive Følelser (13---57).

\textit{α.} Forestillingers Indflydelse paa Følelse (13---18).

\textit{β.} Følelser kan bestemmes 1) ved Forholdet til andre (19---29) eller 2) ved
Forholdet til os selv (30---35).


\textit{γ.} Følelser kan have Karakteren af Drift (36---59).
II. Aktive Følelser (58---59).

\subsection*{d. Sætning 13---57 (og de tilsvarende Afsnit af Tillægget).}
\begin{center}\textit{Passive Følelser.}\end{center}

\textit{α.} To specielle Følelser afleder \textsc{Spinoza} strax af Grundfølelserne (som jo er
Glæde, Sorg og Drift), nemlig Kærlighed (amor), som er Glæde (lætitia) forbunden
med Forestillingen om dens Aarsag, og Had (odium), som er Sorg (tristitia) forbunden
med Forestillingen om dens Aarsag. I begge indeholdes Trang eller Drift,
hist til at fastholde og bevare, her til at fjærne og ophæve. Allerede her viser
Forestillingernes Indflydelse paa Følelsen sig tydelig. Det er ved at forbindes med
Forestillingen om sin Aarsag, at Glæde bliver Kærlighed, Sorg Had. Naar noget,
der i og for sig vilde vække Sorg, tillige har Sider, fra hvilke det kan glæde os,
vil det paa én Gang kunne fremkalde Kærlighed og Had; Sindet vil da bølge frem
og tilbage mellem de to Følelser. Denne Bølgen i Sindet (animi fluctuatio) svarer
paa Følelsens Omraade til Tvivl (imaginationis fluctuatio) paa Forestillingernes Omraade.
Og der er i Grunden, siger \textsc{Spinoza} (17 Schol.), kun en Gradsforskel mellem
de to Fænomener. --- Ved denne sidste Bemærkning kommer \textsc{Spinoza} til at indrømme,
at der ikke gives Forestillinger uden Følelsesvirkning. Desuden maa bemærkes
en vis Uoverensstemmelse mellem 17 (Sætningen og Beviset) og det tilhørende
Scholium. I Sætningen og i Beviset er der Tale om Samtidighed af de to Følelser,
men i Anmærkningen er der Tale om en Bølgen, altsaa om Skiften\footnote[1]{At Ordet Bølgen (fluctuatio) maa forstaas som en Skiften, ses af et Sted i anden Bog, til hvilken \textsc{Spinoza} her henviser. Det hedder nemlig der (II, 44 Schol.), at naar en Forestilling tidligere til forskellige Tider har været forbunden med to forskellige Forestillinger, vil den senere snart fremkalde den ene af disse, snart den anden, og denne Tilstand kaldes Forestillingsbølgen (imaginationis fluctuatio).}. Det er naturligvis
vanskeligt for Iagttagelsen at skelne mellem en hurtig Overgang og en Samtidighed.


\textit{β.} 1) En hel Række af Følelser bestemmes ved, om det, vi elsker eller hader,
bestaar eller forgaar, fremmes eller hæmmes. Hvad der paa udvortes Maade eller
ved Lighed hænger sammen med Kærlighedens eller Hadets Genstand, vil faa Del i
vor Følelse overfor denne. Ad denne Vej opstaar Medlidenhed, Velvillie, Harme,
Misundelse, Undseelse, --- alle ifølge samme Princip, kun i forskellige Modifikationer
og Komplikationer. \textsc{Spinoza} fremhæver særlig (32 Schol.), hvorledes Medlidenhed,
Misundelse og Ærgerrighed udspringer af samme Egenskab ved den menneskelige
Natur, nemlig af den Maade, paa hvilken Følelse af Glæde eller Sorg ændres gennem
Indflydelsen af Forestillinger om deres Aarsager eller om Ting, der paa en eller
% --- p. 66 ---
\setcounter{footnote}{0}%
\opage{66}anden Maade staar i Forbindelse med disse Aarsager. --- Medlidenhed forklarer \textsc{Spinoza}
dog ikke blot ved Forestillingsassociation --- nemlig derved, at vi, naar vi forestiller
os et Væsen, der ligner os selv, behersket af en vis Følelse, selv faar denne
Følelse; --- men ogsaa ved en reflektorisk Tendens til Efterligning af andres Tilstande
og Bevægelser. „Den, der flygter, fordi han ser andre flygte, eller frygter, fordi han
ser andre frygte, eller den, der trækker sin Haand tilbage, naar han ser en anden
brænde sin Haand, --- om ham siger vi, at han efterligner andres Følelser“ (Tillægget
Kap. 33, Note). Der antydes her et Udgangspunkt til Forklaring af sympatiske
Følelser, som senere blev videre gennemført af \textsc{Adam} \textsc{Smith}\footnote[1]{Se \textit{Den nyere Filosofis Historie\textsuperscript{2}}, I, p. 450.}.

2) I Tillægget hedder det, efterat den nysomtalte Gruppe af Følelser er gennemgaaet:
„Dette er de Følelser af Glæde og Sorg, som betinges ved Forestillingen om
en ydre Ting som umiddelbar eller middelbar Aarsag. Jeg gaar nu over til andre
Følelser, der bestemmes ved Forestillingen om en indre Ting som Aarsag“. (Kap. 24
Expl.). Denne indre Ting er os selv. --- I selve tredie Bog er denne Forskel ikke
gjort, idet dér Kærlighed til sig selv og Kærlighed til et andet Væsen omtales paa
én Gang (se Sætning 25). Følger vi Tillægget, der er mere systematisk, kommer vi
nu til en Række af Følelser, der er bestemte ved Forestillingen om os selv og vore
Kaar. Herhen hører Selvtilfredshed (acquiescentia in se ipso), Ydmyghed, Anger,
Stolthed, Nedslaaethed, Undseelse. --- \textsc{Spinoza} mener, at Ydmyghed og Nedslaaethed
er meget sjeldne, da den menneskelige Natur har en uvilkaarlig Trang til at dvæle
ved eller fremhæve alt, hvad der kan styrke dens Virkekraft og holde den oppe.
Meget ofte viser det sig desuden, at de, der synes at være i højeste Grad nedslaaede
og ydmyge, er saare ærgerrige og misundelige. (Tillægget Kap. 29, Expl.).

Gennem hele denne Lære om, hvorledes Forestillinger om ydre eller indre,
umiddelbare eller middelbare Aarsager til Følelser kan bevirke Kvalitets- og Retningsændringer
i Følelseslivet, er \textsc{Spinoza} bleven den egentlige Grundlægger af Motivforskydningens
Teori, der senere udførtes videre af den engelske Psykologi i det attende
og nittende Aarhundrede, og som er saa vigtig til Forstaaelse af sjælelig Udvikling.
Der er sandsynligvis en historisk Forbindelse mellem \textsc{Spinoza} og de engelske Psykologer,
idet \textsc{Hartley}, den første i disses Række, viser Spor af spinozistisk Paavirkning\footnote[2]{Se \textit{Den nyere Filosofis Historie\textsuperscript{2}}, I, p. 455. --- I samme Værk II, p. 372 har jeg med Urette henført en meget klar Fremstilling af Teorien til \textsc{James Mill}; den skyldes \textsc{Mackintosh}.}.


\textit{γ.} Efter Omtalen af de to Grupper af Følelser, der skyldes Forholdet til andre
og til os selv, hedder det i Tillægget (Kap. 31, Expl.): „Hermed har jeg sluttet Beskrivelsen
af de Følelser, der hører under Glæde og Sorg; jeg gaar nu over til de
Følelser, jeg regner under Drift (cupiditas)“. Heller ikke denne Inddeling benyttes
i selve tredie Bog. Og desuden lærer \textsc{Spinoza} selv (Sætning 57, Dem., se ovenfor
p. 64), at Glæde og Sorg egentlig er ét med Selvopholdelsestrang eller Selvopholdelsesdrift,
idet de opstaar, alt efter som denne fremmes eller hæmmes.

Det første Exempel paa Driftfølelse er Begær (desiderium), der (Tillægget Kap. 32)
% --- p. 67 ---
\setcounter{footnote}{0}%
\opage{67}beskrives som Trang eller Drift, der hæmmes ved Forestillingen om noget, som udelukker
Existensen af dens Genstand. Begæret er altsaa en Drift, der ytrer sig under
stærk Modstand. Det forudsætter mere sammensatte og indbyrdes stridende Forestillingsforbindelser
end den simple Drift, der blot er bestemt ved en enkelt Forestilling.
I selve tredie Bog kaldes Begæret en ved en elsket Genstands Fraværelse
bestemt Sorg (36 Schol.), og i Tillægget indrømmes det ogsaa, at det udgør Modsætningen
til den Glæde, der føles ved en hadet Genstands Fraværelse; naar den
alligevel regnes til Driftfølelse, er det paa Grund af det aktive Elements Fremhersken.
--- Der foreligger her en Gradsforskel, bestemt ved Forholdet mellem Følelsen og
Villieselementet i Begæret.

I Sætning 37---38 indskærpes, at en Drift vil røre sig saa meget des stærkere,
jo større den Nedgang i Kraft, altsaa jo større den Sorg har været, som Driften afløser,
og hvorledes det samme gælder Glæde, der afløser Sorg eller Had. Det gælder
ogsaa omvendt, at Sorg og Had bliver stærkere, naar Glæde og Kærlighed gaar forud.
--- \textsc{Spinoza} omtaler ikke, at der her er en Grænse for Forhøjelsen ved Kontrast.
Nedgangen kan nemlig blive saa stor, at ingen Kraft derefter kan udløses, altsaa
ingen Glæde opstaa. I modsat Retning melder Grænsen sig ikke saa let; dog kan
der gives Tilfælde, hvor overvættes Glæde og Betagethed i den Grad forbruger al
Kraft, at der efter deres Ophør kun er Mulighed for Sløvhed, ikke for Lidelse og
Sorg. \textsc{Spinoza} var for meget Optimist til at gaa ind paa denne Side af Sagen.

Had øges ved gensidigt Had, men kan ophæves ved Kærlighed (Sætning 43).
Og selv om der føles Glæde ved, at det, vi hader, lider eller gaar under, vil denne
Glæde dog ikke være uden et Element af Sorg, naar det, der gaar under, er et Væsen
i Lighed med os selv (47).

Jo mere vi tænker os det Væsen, vi elsker eller hader, som isoleret eller „frit“,
des større vil Kærligheden eller Hadet blive, fordi der da ingen andre Forestillinger
og Følelser kan opstaa, som kunde begrænse eller hæmme Genstandens Indflydelse
paa vor Følelse (49).

Efter endnu at have fremsat nogle Sætninger om individuelle Forskelligheder
og komplexe Fænomener indenfor Følelseslivet, bemærker \textsc{Spinoza} (57 Schol.): „Alt
dette angaar de Følelser, der opstaar hos Mennesket som passivt Væsen (qvatenus
patitur). Jeg har nu tilbage at tilføje nogle faa Ting om de Følelser, Mennesket har
som aktivt (qvatenus agit). % PRINTED AS IS: the quote is not closed

\subsection*{e. Sætning 58---59.}
\begin{center}\textit{Aktive Følelser.}\end{center}

\textit{α.} Ved aktive Følelser maatte \textsc{Spinoza}, efter sin almindelige Definition af Aktivitet
og Passivitet (III, Def. 2), forstaa alle Følelser, hvis hele Aarsag (eller dog
største Delen af deres Betingelser) ligger i vor egen Natur. Men da det nu er en
Forudsætning hos ham, at vi kun som klart tænkende eller forstaaende Væsener er
i Tilstande, der helt har deres Grund i os selv, saa bliver for ham kun de Følelser
% --- p. 68 ---
\setcounter{footnote}{0}%
\opage{68}at kalde aktive, der er bestemte ved klare Tanker. Han sætter i Beviset for 59 udtrykkelig
Evne til at forstaa (potentia intelligendi) som ét med Evne til at virke
(potentia agendi), og kan derved beraabe sig paa III, 1. Men III, 9 havde han hævdet,
at Selvopholdelsestrangen ytrer sig, baade forsaavidt vi har dunkle Forestillinger og
forsaavidt vi har klare Forestillinger. Det vil sige, vi har en Trang til at fastholde
det Stade, vi nu engang staar paa, hvad enten det er rationelt eller ikke. Endnu i
Beviset for 58 fastholdes dette; det er først i Beviset for 59, at den snævrere Definition
af aandelig Aktivitet træder i Kraft. --- Intetsteds ser man saa tydelig som her,
hvorledes Intellektualisme og Voluntarisme brydes hos \textsc{Spinoza}.

I \textsc{Leibniz}' Monadelære, ifølge hvilken Universet bestaar af psykiske Væsener
paa alle mulige Trin af Dunkelhed og Klarhed, og dog hvert for sig udfoldende sig
efter en indre Lov, er \textsc{Spinoza}'s Selvopholdelsesbegreb kommet mere til sin Ret. Og
senere, i det nittende Aarhundrede, har man fra forskellige Sider hævdet Vigtigheden af,
at Mennesket lever ud af sig selv, hvilket Trin af Sandhedserkendelse det saa end staar
paa. Man havde gjort Erfaringer om det menneskelige Sjælelivs indviklede og sammensatte
Natur og havde ikke længere Rationalismens Tro paa, at Menneskets Væsen
var udtømt. ved „Fornuften“. Man vil for enhver Pris have Natur, Oprindelighed, % PRINTED AS IS: a stray stop
fordi det er en nødvendig Betingelse for at komme videre. Og man finder Mulighed
for Selvvirksomhed paa ethvert Trin af Udvikling, idet der kan røre sig en
Stræben efter at frigøre sig fra alt, hvad der paa udvortes og overfladisk Maade er
optaget udefra. Da, og kun da, vindes der Klarhed over, hvorledes Livet rører sig
hos de enkelte. En saadan Synsmaade kom i Danmark frem i \textsc{Poul} \textsc{Møller}'s og
\textsc{Kierkegaard}'s Begreb om personlig Sandhed\footnote[1]{\textit{Danske Filosofer}, p. 123---127; 125.}. I dette Begreb er der gjort Alvor af
den Sætning hos \textsc{Spinoza}, at Mennesket stræber efter at hævde sig, ikke blot, naar
det ledes af klare, men ogsaa naar det ledes af dunkle Tanker. \textsc{Spinoza} selv havde
jo ogsaa den Overbevisning, at naar man blot arbejdede konsekvent med de Tanker,
man har, vil der banes Vej for Klarheden og Sandheden. (Se ovenfor p. 17).

\textit{β.} Ikke alle Følelser kan faa den aktive Karakter. Sorg er, som Udtryk for
Hæmning og Formindskelse af Virkekraft, paa Forhaand udelukket. Kun Glæde og
Drift kan være betingede ved klar Erkendelse. Klar Erkendelse skyldes Selvvirksomhed
og volder Erkendelsesglæde og derved Trang til fortsat Selvvirksomhed (58
Dem.), --- en Kærlighed til aandelig Frihed (amor libertatis), som \textsc{Spinoza} i en senere
Sammenhæng kalder den (V, 10 Schol.).

Al Virken, der udspringer af aktiv Følelse, henregner \textsc{Spinoza} under Aandskraft
(fortitudo), som er Livsmod (animositas), naar den, f. Ex. i Selvbeherskelse og
Aandsnærværelse, gaar ud paa, hvad der er godt for Individet selv, og er Højmod
(generositas), naar den, f. Ex. i Mildhed og Beskedenhed, gaar ud paa, hvad der er
godt for andre (59 Sehol.). % PRINTED AS IS: Sehol.

Det er en Følge af den skarpe Modsætning, \textsc{Spinoza} mener at finde mellem
Aktivitet og Passivitet, og derfor mellem aktive og passive Følelser, at han nægter,
at Sorg kan være en aktiv Følelse. Der rører sig jo dog ogsaa i Sorgen en Trang,
% --- p. 69 ---
\setcounter{footnote}{0}%
\opage{69}nemlig til trods alt at fastholde det tabtes Værdi, eller rettere Sorgen selv er et Udtryk
for en Brydning mellem Trangen til at fastholde noget værdifuldt og Erfaringen
om Tabet af det. Derfor maa Menneskets inderste Natur ogsaa kunne ytre sig i
Sorgen. Og i hvert Tilfælde kan Sorgen blive Element i en mere sammensat Følelse,
der kan have Aktivitetens Karakter, selv om hvert enkelt Element i den skulde være
passivt; Aktiviteten viser sig da i den Magt, hvorved Elementerne forenes trods deres
forskellige, maaske modsatte Karakter. \textsc{Spinoza} har ogsaa set dette (58 Schol.). Men
han gaar meget let henover de komplexe Følelser; han lader sig nøje med at bemærke
(59 Schol.), at der er overordentlig mange forskellige Sammensætninger af
Grundfølelser mulige. Han har ikke fulgt den Vej, \textsc{Giordano} \textsc{Bruno} gik i sit begejstrede
Skrift Dell' heroici furori. \textsc{Bbuno} og \textsc{Shakespeare} er de første, der har % PRINTED AS IS: BBUNO
formaaet at fremstille komplexe Følelser i deres Ejendommelighed og i deres Værdi.
\textsc{Spinoza} selv har i „Livsmod“ og „Højmod“ beskrevet to komplexe Følelser (Totalfølelser)
af stor Værdi; men han har ikke villet indrømme, at „passive“ Elementer
indgik i dem\footnote[1]{Smlgn. \textit{Den store Humor}, p. 39 f; 154---159.}. ---

\textsc{Spinoza}'s Fremstilling af Følelsernes Psykologi er med Rette berømt. Han har
stræbt at behandle denne Del af Psykologien med videnskabelig Objektivitet. Der
er i det foregaaende dvælet ved de Inkonsekvenser, som fremkommer ved forskellige
Synspunkters Brydning, og ved den Intellektualisme, der atter og atter gør sig
gældende trods betydningsfulde Tilløb i anden Retning. Men alle de vigtigste
Spørgsmaal, Følelsernes Psykologi maa opkaste, kommer frem hos ham og er ogsaa
fremhævet i det foregaaende. Saaledes Forholdet mellem Følelse og Trang,
Indflydelse af Forestilling paa Følelse, Motivforskydning, Kontrastvirkning i dens
Forskel fra Sammenligning, Expansion o. s. v.

Følelsens Psykologi staar hos \textsc{Spinoza} som Indledning til en Etik. Mennesket
kan blive ufrit, overvældes af Følelser, der fremkaldes udefra. Spørgsmaalet er da,
hvorledes aandelig Frihed og Kraft bliver mulig. Allerede i tredie Bog er der givet
flere Antydninger i denne Retning. Naar Had kan overvindes ved Kærlighed (43.
47), --- naar nye Forestillingsforbindelser kan rive Følelsen bort fra dens oprindelige
Genstand (48---49), --- og naar der endog gives Følelser, som ikke er af passiv
Natur, men er Udtryk for Aandskraft (58---59), saa er der Muligheder for Udvikling
i Retning af Aktivitet og Frihed. Paa dette Punkt tager fjerde Bog Spørgsmaalet op.

% --- p. 70 ---
\setcounter{footnote}{0}%
\opage{70}\section*{Fjerde Bog.}
\subsection*{a. Fortale, Definitioner og Axiomer.}

\textit{α.} I Fortalen til fjerde Bog gøres Overgangen til Etiken, som har givet hele
Værket dens Navn. For en Etik er Værdibegreber af en eller anden Art nødvendige.
Og dog begynder \textsc{Spinoza} med at gentage sin Polemik fra første Bogs Slutning mod
at anvende Værdibegreber paa Naturen. Men Mennesker og deres Handlinger hører
jo dog ogsaa --- efter \textsc{Spinoza}'s energiske Udtalelser i Slutningen af første Bog og i
Fortalen til tredie Bog --- med til Naturen! Der kræves da en Forklaring af, hvorledes
\textsc{Spinoza} selv alligevel kan mene sig berettiget til at anstille etisk Vurdering
af menneskelige Egenskaber og Handlinger.

Naar Værdibegrebers Anvendelse i Naturopfattelsen afvises af \textsc{Spinoza}, er Begrundelsen
den, at saadanne Begreber dannes ved, at vi sammenligner Fænomenerne
i Forhold til Almenbegreber eller i Forhold til Formaal. Vi kalder Tingene gode
eller onde, fuldkomne eller ufuldkomne, alt efter som de stemmer med de almindelige
Typer, vi har dannet os, eller med det Formaal, vi mener, at de skal tjene.
Men saadan Sammenligning og dens Resultater har aldeles intet at gøre med Tingene,
som de er „i sig selv“. Det uendelige Væsen, som vi kalder Gud eller Naturen,
virker med samme Nødvendighed som den, ved hvilken den existerer. Der er intet
højere Begreb og intet Formaal, efter hvilke det skulde kunne maales.

Alligevel, erklærer \textsc{Spinoza} pludselig, maa vi fastholde Begreberne om Fuldkommenhed
og Ufuldkommenhed, godt og ondt, ti vi ønsker (cupimus) at danne
et anskueligt Menneskeideal (exemplar vitæ humanæ, qvod intueamur), i Forhold
til hvilket vi kan bedømme menneskelige Handlinger!

Spørges der, hvorfor vi ønsker dette, svares der, at vor Virkekraft er begrænset;
den kan stadig forøges eller formindskes, og der gives --- som det eneste Axiom i
fjerde Bog hævder --- for ethvert Enkeltvæsens Vedkommende stedse noget, som er
mægtigere end det; --- derfor stræber vi naturlig efter den størst mulige Virkekraft
og dermed efter den størst mulige Frihed (i Betydning af Evne til at have Aarsagen
til sine Tilstande og Handlinger i sig selv). Paa sit Højdepunkt gør denne Magt
eller Frihed (som \textsc{Spinoza} viser i femte Bog) Mennesket til ét med Substansen (Gud
eller Naturen), det eneste der bestaar ved sig selv og forstaas ved sig selv. \textsc{Spinoza}
kan forsaavidt sige, at han stadig fastholder den i anden Bog (Def. 6) givne Defini
% --- p. 71 ---
\setcounter{footnote}{0}%
\opage{71}tion af Fuldkommenhed som ét med Realitet. Men jeg har allerede ovenfor (p. 33
---35) haft Lejlighed til at udtale, at dette kosmologiske Fuldkommenheds- eller
Realitetsbegreb allerede er farvet af en Værdibestemmelse, skyldes en Vurdering,
en Idealisering. Kun derved bliver det muligt, at ogsaa det etiske Fuldkommenhedsbegreb
kan finde en Plads indenfor \textsc{Spinoza}'s System. Etisk Stræben bliver
ifølge \textsc{Spinoza} kun mulig derved, at der i hvert Enkeltvæsen vibrerer en Straale af
Substans, og at denne Straale kan naa stedse højere Klarhedsgrader.

Egentlig kunde \textsc{Spinoza} da meget vel have sagt, at det er Substansen (Gud
eller Naturen), der er det højeste etiske Forbillede. Og der bliver derved en Analogi
mellem hans Filosofi og hans Opfattelse af positiv Religion. De positive Religioners
Gudsforestillinger har for ham kun Betydning paa det praktiske Omraade.
Gud er i Religionen „et Forbillede paa det sande Liv“ (veræ vitæ exemplar)
(Tr. theol. polit. Kap. 14), medens Religionen eller Teologien ikke kan lære os noget
om Guds Væsen i sig selv. Det er karakteristisk, at Udtrykket „Forbillede“ (exemplar)
i Religionen kan bruges om Gud, i Filosofien om det sande Menneske. Den
kosmologiske Betragtning danner en Overgang mellem begge Tankeverdener.

\textit{β.} Men der ligger stadig en Vanskelighed i den Begrundelse, \textsc{Spinoza} giver af,
at Værdibegreber ikke kan anvendec paa selve Tilværelsen. Begrundelsen skulde jo % PRINTED AS IS: anvendec
ligge i, at de er Tankeformer, Begreber, vi danner derved, at vi sammenligner Tingene
indbyrdes (cogitandi modi seu notiones, qvas formamus ex eo, qvod res ad
invicem comparamus). (Se ogsaa Ep. 21). Nu gælder det dog om alle vore Begreber,
ogsaa de rent teoretiske, at de dannes ved Sammenligning. Al Videnskab beror,
ifølge \textsc{Spinoza}'s egen Erkendelsesteori, paa Sammenligning. Den anden Erkendelsesart
(ratio), den egentlige Videnskab i Modsætning til den tilfældige Erfaring, bestaar
jo i, at vi ikke betragter Tingene isolerede, men forstaar deres „Overensstemmelser,
Forskelligheder og Modsætninger“; kun derved bliver klar og tydelig Erkendelse til
(II, 29 Schol., se ovenfor p. 43). Alle Emner, ikke blot Værdier, bliver kun gennem
Sammenligning Genstand for klare og tydelige Begreber. Selve Begrebet Sandhed
eller Virkelighed kan kun anvendes formedelst en Maalestok (norma veritatis), der
ifølge \textsc{Spinoza} er given i det rent logiske Forhold mellem Grund og Følge (se ovenfor
p. 6---7). Ligesom de teoretiske Begreber har deres Norm i dette Forhold, saaledes
har Værdibegreberne for \textsc{Spinoza} deres Norm i Begrebet om den højeste Kraft
eller Frihed. Som Sammenligningsbegreber betragtede staar altsaa Værdibegreberne
paa Linie med alle andre Begreber. Forskellen er, at Værdinormen (Grundværdien)
er mere speciel og konkret end Sandhedsnormen, idet den udspringer af den menneskelige
Naturs Tarv og Trang. I denne specielle Form kan Værdinormen ikke
overføres paa Tilværelsen. Derved adskiller ifølge \textsc{Spinoza} Filosofi sig fra Teologi.
Da Teologerne betragter Gud som et idealt Menneske, er det konsekvent, at de overfører
Værdibegreber paa ham og tillægger ham Ønsker, Glæde og Mishag overfor,
hvad der sker i Verden (Ep. 23 og 54).

Selv fra det konkrete menneskelige Synspunkt har Værdibegreberne ifølge
\textsc{Spinoza} kun begrænset Gyldighed, nemlig saalænge Mennesket ikke har naaet den
% --- p. 72 ---
\setcounter{footnote}{0}%
\opage{72}fuldkomne Frihed, men endnu maa kæmpe for at hævde sig i sin ejendommelige
Væren. Kun saalænge behøver det et Forbillede, og kun saalænge har det Brug for
Begreber om godt og ondt. Henimod Slutningen af fjerde Bog opstiller \textsc{Spinoza} den
Sætning, at, hvis Menneskene fødtes frie, hvilket forudsætter, at de kun havde fuldstændige
(adækvate) Tanker, vilde de ikke danne noget Begreb om godt og ondt (68).

Paa Højdepunktet af aandelig Udvikling kommer vi altsaa „hinsides godt og
ondt“. Forsaavidt kan det etiske, paa \textsc{Spinoza}'s Standpunkt, siges kun at være
Middel og Forberedelse til at virkeliggøre Menneskevæsenets Ideal, som for \textsc{Spinoza}
bestod i den Frihed, der er ét med Erkendelse af vor Aands Sammenhæng med
hele Naturen. I Afhandlingen om Forstandens Forbedring, hvor \textsc{Spinoza} beretter
om det Valg, han traf mellem de forskellige Maal, Mennesker kan sætte sig, nævner
han Moralfilosofi som et af Midlerne til at naa det Maal, der for ham er det sande.
Dette stemmer ganske med Tankegangen i Ethica\footnote[1]{Det er i Virkeligheden ogsaa en misforstaaet Idealisme, der vil sætte det etiske som Livets højeste Aandsform. Smlgn. herom min \textit{Etik} IV, 2. 7 (4. Udg. p. 88. 103).}. --- Det er ikke rigtigt at sige med
\textsc{Wundt} (Ethik, II, p. 113), at Kosmologi og Erkendelsesteori for \textsc{Spinoza} kun var
Hjælpemiddel og Forberedelse for at naa til en etisk Verdensanskuelse. Det etiske
er for \textsc{Spinoza} et Gennemgangsstadium. Derimod kan man sige, at alle de Tanker,
han har nedlagt i sit System, for ham betegner Vej til det, der for ham paa én
Gang var Tankens og Livstrangens højeste Maal, selv om det efter hans Opfattelse
vilde være meningsløst at betragte det som hele Tilværelsens Maal, da Tilværelsen
i sig selv er evig og uforanderlig. Det er jo netop denne Evighed og Uforanderlighed,
det for Mennesket gælder at faa Del i.

Et Problem rører \textsc{Spinoza} ikke ved (uden lige i Begyndelsen af den erkendelsesteoretiske
Afhandling), nemlig, om der gives en Værdinorm, som kan og maa anerkendes
af alle Mennesker, ligesom der gives en Sandhedsnorm, for hvilken alle
tilsidst maa bøje sig. Det er paa dette Punkt, den egentlige Brod i etisk Henseende
mere og mere viser sig at ligge.

\textit{γ.} Den sidste Vanskelighed, hvad Værdibegrebets Stilling indenfor \textsc{Spinoza}'s
System angaar, bliver den, hvorledes der overhovedet kan gives andet end det evige
og uforanderlige. Hvorfor er der Forskelle og Grader, hvad „Realitet eller Fuldkommenhed“
angaar, saa at der kan blive Tale om Overgang fra en Tilstand til en
anden, som i Definitionen af Glæde og Sorg i tredie Bog, og saa at der kan stille
sig den Opgave at naa fra lavere Trin til højere? --- Det er med andre Ord Tids- og
Kvalitetsforskellighederne, der volder Vanskeligheder. Hvorfor existerer der overhovedet
andet end Substansen, det, der har sin Grund i sig selv og derfor ene kan
kaldes fuldkommen frit? Hvorfor koster det saa haardt et Arbejde at naa op til det
Stade, hvor hverken Tilskyndelser eller Værdibegreber skal gælde?

\textsc{Spinoza}'s Svar er givet i Slutningsordene til første Bog: „Der har ikke været
Mangel paa Stof til at frembringe alt, lige fra den højeste til den laveste Grad af
Fuldkommenhed. Naturens Love var saa omfattende, at de kunde rumme alt, hvad
en uendelig Tanke kunde indeholde“. Men dette Svar vilde kun afgøre Sagen, der
% --- p. 73 ---
\setcounter{footnote}{0}%
\opage{73}som \textsc{Spinoza}'s Grundbegreber kunde omfatte, ikke blot det konstante, men ogsaa
det skiftende, ikke blot det fuldendte, men ogsaa det stræbende, ikke blot det kontinuerlige,
men ogsaa det diskontinuerlige. Det er Tilværelsens Uforanderlighed i
dens inderste Grund, der her ytrer sig i sine Konsekvenser og volder Vanskeligheder,
baade for Etikens og Erkendelsesteoriens Vedkommende. Der røber sig her
en Forudsætning, som gaar langt tilbage i Tænkningens Historie, og hvis Haardnakkethed
maaske bedst ses deraf, at den endnu gør sig gældende hos den kritiske
Filosofis Grundlægger i hans dogmatiske Fastholden af „Ding an sich's“ Konstans\footnote[1]{Se herom \textit{Kontinuiteten i Kant's filosofiske Udviklingsgang} (Vid. Selsk. Skr. 1893), p. 34. --- \textit{Den menneskelige Tanke}, p. 327 f.}.

Det maa dog straks tilføjes, at denne Vanskelighed vel træder særlig skarpt
frem i \textsc{Spinoza}'s System, men at den gør sig gældende i al Videnskab. Videnskaben
kan føre os fra Tids- og Kvalitetsforskelligheder til almene Love, som maaske endog
kan formuleres i matematisk Lignings Form; men den kan ikke af disse Love og
Ligninger udlede, at der existerer netop disse eller hine Kvalitetsforskelligheder og
Tidsforskelligheder i Verden. Det er netop en Hovedinteresse ved \textsc{Spinoza}'s System,
at det saa klart fører os til at se dette Forhold mellem de Grundbegreber, i hvilke
vi tænker Tilværelsen, og de faktisk givne Emner, der udgør Tilværelsens Indhold.
\textsc{Spinoza} har ikke selv været sig dette Forhold bevidst. Hans Filosofi er, som jeg
ovenfor har udtrykt det, et Forsøg paa at tænke Tilværelsen ved Hjælp af de formale
Kategorier alene.

\subsection*{b. Sætning 1---18.}
\begin{center}\textit{Følelsernes Styrkeforhold.}\end{center}

\textit{α.} En Følelses Styrke, dens Evne til at holde sig, bestemmes ikke ved vor
egen Selvopholdelsestrang alene, men beror paa dennes Forhold til de ydre Aarsager,
vi paavirkes af. Jeg er kun en Del af den hele Natur og undergaar derfor mange
Forandringer, der ikke kan forklares ved mit eget Væsen alene. Og endnu mindre
kan hele Naturens Orden forklares ud fra mit Væsen. Men der er Mulighed for, at
en enkelt Følelse kan faa Herredømmet over alle de andre Følelser, naar Forholdet
mellem indre og ydre Aarsager begunstiger den (2---6). Herved gør Forskellen
mellem sand og falsk Erkendelse sig gældende, dog ikke direkte. Ti som paavist i
II, 32---35 er Vildfarelse ikke noget positivt, men skyldes Fraværelsen af udfyldende
og berigtigende Forestillinger, og derfor ophæver sand Erkendelse ikke i og for sig
det positive i den Tilstand, hvortil Vildfarelsen hører (1). Jeg har samme Sansebillede
af Solen, hvad enten jeg mener, at dens Afstand fra mig er lang eller kort.
Gaar den sande Erkendelse op for mig, ændres nogle af de Forestillinger, der knyttede
sig til Sansebilledet, men ikke dette selv. Vil man indvende, at Frygten for
en Fare falder bort, naar det viser sig, at der i Virkeligheden ingen Fare var, saa
svarer \textsc{Spinoza}, at Frygten ogsaa falder bort, naar Efterretningen om, at der ingen
Fare er, er falsk. Det er altsaa ikke det sande som saadant (qvatenus verum), der
% --- p. 74 ---
\setcounter{footnote}{0}%
\opage{74}ophæver Frygten, men den Styrke, med hvilken de Forestillinger, der udelukker de
indbildte Farer, optræder (1 Schol.).

\textit{β.} Naar \textsc{Spinoza} her taler om Forestillingers Styrke, saaledes at denne synes
at være forskellig fra den logiske Modsætning (Udelukkelsesforholdet) mellem dem
og andre Forestillinger, maa det antages, at denne Styrke beror paa de Følelser,
der er knyttede til de udelukkende Forestillinger. Ti i 7 hævdes, at en Følelse ikke
fortrænges uden ved en anden og stærkere Følelse. Beviset føres ud fra Inertisætningen.
Hverken en aandelig eller en legemlig Tilstand ophæves uden formedelst
en anden og stærkere Tilstand.

Den vigtige Sætning, \textsc{Spinoza} her opstiller, forekommer ogsaa hos \textsc{Platon},
naar han (som i Phaidon og i Staten) stiller Erkendelsestrangen som Modvægt mod
Trang til Sansenydelse eller Ære --- samme Anvendelse af Sætningen, som vi vil se
\textsc{Spinoza} gøre i femte Bog. \textsc{Spinoza} har kunnet komme til Sætningen ad sin egen
Erfarings Vej. Han beretter i Begyndelsen af Afhandlingen om Forstandens Forbedring,
hvorledes han i Erkendelsen fandt et Gode, der kunde stille alle andre
Goder i Skygge. Han erfarede, at han i Eftertankens Stunder (modo possem serio
deliberare) var befriet fra alt, hvad der ellers kunde hilde. Muligvis er han
dog her paavirket af \textsc{Baco}, der siger, at ligesom man paa Jagten bruger en Falk
til at fange andre Fugle, saaledes kan en Følelse bruges til at hæmme eller regulere
andre Følelser (De Augmentis Scientiarum. VII, 3). I flere Træk minder \textsc{Spinoza}'s
Etik om \textsc{Baco}'s Fremstilling. Saaledes inddeler \textsc{Baco} Etiken i Læren om
Forbilledet (exemplar) og i Læren om Arbejdet paa at føre til Overensstemmelse
med Forbilledet (cultura animi) (De Augm. VII, 1). Og denne sidste Lære inddeles
igen i Læren om Følelserne (de affectibus) og Læren om Midlerne til at regulere
dem (de remediis) (De Augm. VII, 3). Baade i Skema og i Udtryk fremtræder
noget beslægtet hos \textsc{Spinoza}. --- Ogsaa efter \textsc{Spinoza}'s Tid er den baade i psykologisk
og i pædagogisk Henseende vigtige Sætning bleven hævdet. Saaledes af
\textsc{Rousseau}, der i Emile (Livre IV) udtrykker den saaledes: On n'a de prise sur les
passions que par des passions. \textsc{Kant} hævder den (i tredie Afsnit af Afhandlingen
om negative Størrelser og i Indledningen til „Tugendlehre“\footnote[1]{Mærk den for \textsc{Kant} skæbnesvangre Distinktion mellem Frihed som abstrakt Evne (facultas) og som aktuel Styrke (robur).}, uden at agte paa den
Modsigelse, han derved kommer i til sit Frihedsbegreb), og \textsc{Schiller} anvender den
i det ottende Brev om æstetisk Opdragelse. I min Religionsfilosofi (§ 116) har jeg
vist, hvilken Betydning den har for det religiøse Problem.

Sætningen trænger dog til et Korrektiv. Ti foruden ved en anden, stærkere
Følelses Optræden kan en Følelse, Trang eller Drift ogsaa forsvinde af Mangel paa
Næring ved de Forhold, under hvilke Livet føres, ligesom et Organ kan forsvinde
af Mangel paa Brug og --- som Følge deraf --- af Næring. Og naar en Følelse fortrænges
ved en ny Følelse, er det vel egentlig ogsaa, fordi denne sidste lægger Beslag
paa al aandelig og legemlig Energi.

Af Sætning 7 følger nu, at skal Erkendelse af godt og ondt blive en Magt i
% --- p. 75 ---
\setcounter{footnote}{0}%
\opage{75}Sindet, maa den optræde som Følelse (Trang eller Drift) (14). Men ifølge III, 9
Schol. er det jo en Trang eller Drift, der fører os til at betragte noget som godt
eller ondt. Erkendelsen af godt og ondt er altsaa ikke andet end Bevidsthed om
Glæde eller Sorg, altsaa om den Maade, paa hvilken vor Virkekraft øges eller
mindskes (8). Hvad det gælder om, er, at denne Erkendelse kan holde sig, uagtet
det nærværende paavirker os stærkere end det fraværende og det fremtidige, og
uagtet ogsaa Forskellen mellem det i vore Øjne tilfældige, mulige og nødvendige
øver stor Indflydelse paa os (9---13). Hvor Virkningen af, hvad der i Øjeblikket
paavirker os, staar i skarp Modsætning til vor hidtil vundne Erkendelse af godt og
ondt, opstaar der en Strid i vort Indre. Prædikeren har da Ret i, at hvo som øger
sin Kundskab, øger sin Smerte. Dette er dog ingen Grund til at prise Uvidenheden;
men det indeholder Opfordring til at undersøge Betingelserne for aandelig Styrke
og aandelig Afmagt (17 Schol.).

\textit{γ.} I Sætning 18 hævdes, at den Drift, der opstaar af Glæde, under ellers lige
Forhold er stærkere end den, der opstaar af Sorg. Begrundelsen ligger i, at Sorg
(Ulyst) er Symptom paa Tilbagegang i Virkekraft.

I Anmærkningen til 18 kommer \textsc{Spinoza} selv til at indrømme det uhensigtsmæssige
i den „geometriske“ Fremstillingsmaade, navnlig fordi den fører til stor
Vidtløftighed og hindrer Overblikket. Derfor vil han, før han i det enkelte viser,
hvorledes Mennesket kan udvikle sig fra at være et passivt Bytte for Sindsbevægelser
til at besidde Selvbeherskelse og Frihed, give en Oversigt over, hvad Fornuften
paaviser som Betingelsen for at naa dette Maal.

Det gælder at virke efter sin egen Naturs Love og derigennem naa til større
Fuldkommenhed, et Maal, der maa søges for dets egen Skyld, ikke som Middel til
noget andet. De, som helt mister Selvopholdelsestrangen og berøver sig selv Livet,
er --- maaske uden selv at mærke det --- overvældet af ydre Aarsager (causæ latentes
externæ. 20 Schol.).

Blandt Midlerne til Selvopholdelse er især Samvirken med andre Mennesker
af stor Betydning. Naar flere forener deres Kræfter om fælles Formaal, bliver de
ligesom én Sjæl og ét Legeme. Mennesker, som ledes af Fornuften, vil derfor ikke
ønske noget for sig selv, som de ikke ogsaa ønsker andre; de vil være retfærdige,
trofaste og hæderlige. --- Man ser allerede her Forskellen mellem \textsc{Spinoza}'s og
\textsc{Hobbes}' Etik. For \textsc{Spinoza} er social Etik en positiv og direkte Fortsættelse af individuel
Etik, medens den for \textsc{Hobbes} overvejende staar i et negativt eller indirekte
Forhold til denne. I et Brev (Ep. 50) har \textsc{Spinoza} selv udtalt sig om denne Forskel.

\subsection*{c. Sætning 19---28.}
\begin{center}\textit{Den etiske Karakter og det højeste Gode.}\end{center}

\textit{α.} Selvopholdelsestrangen afgør, hvad der anses for godt og ondt, og enhver
higer derfor ifølge sin egen Naturs Love efter, hvad han anser for godt, og afskyr,
hvad han anser for ondt. Hin Trang er ét med Menneskets Væsen. Og Erkendelsen
% --- p. 76 ---
\setcounter{footnote}{0}%
\opage{76}af godt og ondt er, som vi har set, ikke andet end Bevidsthed om Glæde og Sorg,
d. v. s. Forøgelse og Formindskelse af Evnen til at hævde os selv. --- Man ser her,
hvorledes Erkendelse, Følelse og Villie tilsidst gaar i ét for \textsc{Spinoza}, dog saaledes,
at det er en Villen og Virken, der er vort inderste Væsen og i Erkendelse og Følelse
bliver sig selv bevidst. Her ligger Muligheden for en rigere og dybere Psykologi
end den udprægede intellektualistiske Tendens, \textsc{Spinoza} medbragte til sin Filosoferen,
kunde lade komme til Udvikling.

Men naar Mennesket handler i Blinde eller ledes af ufuldstændige (uadækvate)
Forestillinger, handler det ifølge \textsc{Spinoza} ikke ud af sig selv; det er da ikke aktivt,
ikke frit. Der er derfor intet, om hvilket vi med Sikkerhed kan vide, at det er godt,
undtagen Erkendelse. \textsc{Spinoza} taler her, som man ser af det selvbiografiske Fragment
i Afhandlingen om Forstandens Forbedring, af egen Erfaring. Ydre Ting, siger
han i Tillægget til fjerde Bog, er kun gode for os, forsaavidt de er Betingelser for
Delagtighed i det Aandens Liv, der bestaar i Erkendelse (ut mentis vita fruamur,
qvæ intelligentia definitur. IV App. Kap. 5).

\textit{β.} \textsc{Spinoza} har, som tidligere omtalt, ikke mærket til nogen Vanskelighed ved
Forholdet mellem de to Udgangspunkter, der betinger hans psykologiske og etiske
Tankegang, --- det ene, ifølge hvilket Selvopholdelsestrang, --- det andet, ifølge hyilket % PRINTED AS IS: hyilket
Erkendelsestrangen udgør Menneskets Væsen. Modsætningen er der, men den behøvede
ikke at være uovervindelig. Selvopholdelsestrangen kan føre til at betragte
Erkendelsen som et Gode, først maaske som et afledet Gode, et Middel, senere,
maaske efter en hel psykologisk Proces, som et Formaal, et Gode i sig selv, saa at
det kan siges (29 Dem.), at det ikke er for noget andet Formaals Skyld, at vi
stræber efter Erkendelse. Naar \textsc{Spinoza} begrunder denne Sætning ved, at vi overhovedet
ikke for noget Formaals Skyld stræber at opholde os selv (25), saa ligger
her den Forudsætning til Grund, at Erkendelse og Selvopholdelse er ét. Og dog
viser Erfaring, hvad ogsaa \textsc{Spinoza}'s egen Erfaring havde vist ham, at de ikke uden
videre falder sammen. Gennem tredie Bogs Paavisning af Motivforskydningens Mulighed
havde \textsc{Spinoza} skaffet sig et Grundlag til Forstaaelse af, hvorledes det, der
først har Værdi som Middel, senere kan blive Formaal i og for sig selv. Det er
forunderligt, at han ikke gør Brug af denne Teori her i fjerde Bog. Hans sunde
Blik som Psykolog svigter ham her, og Grunden dertil maa søges i, at han efter
sin Filosofis Grundbegreber ikke kunde tillægge Vorden, Overgang og Udvikling
nogen afgørende Betydning. Han betragter Selvopholdelsestrang og Erkendelsestrang
som identiske, istedetfor at holde sig til, at denne ifølge naturlige Love kan
udvikle sig af hin og maaske paa sit Højdepunkt samle al Aaandskraft i sig. --- I % PRINTED AS IS: Aaandskraft
femte Bog vil denne Ejendommelighed ved hans Tankegang træde endnu skarpere
frem.

Ser vi tilbage paa Forholdet mellem Erkendelse og Villie hos \textsc{Spinoza}, finder
vi tre Synsmaader, der brydes hos ham: 1) Menneskets Væsen er Tanke (III, 1---3,
cfr. II, 40); 2) Tanke og Trang er identiske (IV, 25---26, cfr. II, 49); 3) Tanken er
afhængig af Trangen (III, 9 Schol.). En genetisk Psykologi vilde forsøge paa at
% --- p. 77 ---
\setcounter{footnote}{0}%
\opage{77}vise, hvorledes der kan foregaa en Udvikling fra 3) gennem 2) til 1). Men det
dogmatiske og det genetiske Synspunkt brydes hos \textsc{Spinoza} paa en Maade, der ikke
har været ham selv bevidst.

\textit{γ.} Den sande Erkendelse, der for \textsc{Spinoza} staar som det højeste Gode, vil i
en Etik, som ikke blot skal gælde for Videnskabsmænd (eller rettere for Videnskabsmænd,
forsaavidt de blot er Videnskabsmænd), være en ensidig Bestemmelse.
Det Aandsliv (mentis vita), \textsc{Spinoza} taler om, vil for os kunne fremtræde under
flere Former end den intellektuelle. Stemnings- og Fantasilivet er Former for Aandsliv,
der ligesaavel som det rene Tankeliv kan fyldestgøre aandelig Trang, og kan
gøre det uden at bryde Tankelivets Normer. Naar vi vil lægge os \textsc{Spinoza}'s Etik
til Rette, maa vi derfor istedetfor Erkendelse eller Forstaaelse sætte aandelig Kultur
overhovedet. Derved opstaar der ganske vist nye Problemer, idet der kan opstaa
Brydninger mellem de forskellige Kulturomraader; men det er netop et Vidnesbyrd
om den Rigdom, det menneskelige Aandsliv kan sidde inde med\footnote[1]{Smlgn. min \textit{Etik}, Kap. 27---35 og min Afhandling \textit{Aandelig Kultur} (Mindre Skrifter, Tredie Række).}.

\subsection*{d. Sætning 29---37 (og Tillægget Kap. 9---29).}
\begin{center}\textit{Samfundslivets Betydning.}\end{center}

\textit{α.} Saalænge Menneskenes Følelser væsentlig bestemmes udefra, er det tilfældigt,
om de stemmer overens, og saa er det ogsaa tilfældigt, om det enkelte
Menneske stemmer overens med sig selv. Der er snart Mulighed for Modsætning
og Strid, ti det, forskellige Mennesker stræber efter, kan være noget, som kun en
eneste kan besidde. Naar Per og Poul elsker en og samme Ting, er de jo enige om
dens Værdi, men i Virkeligheden staar de i Modsætning til hinanden, ti hvis Per
besidder Tingen, vil han føle Glæde, men Poul vil føle Sorg, og omvendt\footnote[2]{Allerede \textsc{Aristoteles} bemærker, at det, at alle vil ét, ikke er nok til Endrægtighed, naar enhver vil Besiddelsen for sig alene (\textit{Eth. Nic.} IX, 6). --- \textsc{Kant} har en lignende Betragtning (\textit{Kritik der praktischen Vernunft} § 4 Anm.) og anfører Frantz I's Ord: „Hvad min Broder Karl vil, vil jeg ogsaa!“ --- De vilde nemlig begge have Milano.} (34 Schol.).
Skal der være Endrægtighed og virkeligt Samliv, maa Følelsen være knyttet til noget,
som kan være et fælles Gode, saa at den enes Stræben ikke hindrer den andens,
men de kan hjælpe og støtte hinanden. Og dette vil være Tilfældet, jo mere Menneskene
ledes af aktive Følelser, altsaa saadanne, som er betingede ved Erkendelse
og Fornuft. Fornuften er et aandeligt Gode, ved hvis Erhvervelse den ene ikke
staar den anden ivejen. Her er det intet Tilfælde, naar Mennesker stemmer overens.
Uden Evne til at glædes ved aandelige Goder vilde Mennesket ikke være Menneske.
Og et Gode som Erkendelse vil den enkelte ikke beholde for sig selv, men stræbe
at meddele til alle. Ogsaa hvad materielle Goder angaar, vil det være af Betydning,
at Kræfterne forenes; man kan da bedre sikre sig mod Fare og skaffe sig, hvad der
kræves til at opholde Livet.

Baade paa den materielle og den aandelige Kulturs Omraade er altsaa fælles
% --- p. 78 ---
\setcounter{footnote}{0}%
\opage{78}Stræben mulig, naar Mennesket ledes af Fornuften, --- af det, der udgør Menneskets
egentlige Væsen.

\textit{β.} \textsc{Spinoza} er dog ikke saa optimistisk, at han overser de store Hindringer
for et saadant Fællesskab. Der kræves Kunst og Aarvaagenhed for at lede Menneskene
i denne Retning. Der kræves en Opdragelse til Frihed. Og der er ingen
Maade, paa hvilken et Menneske bedre kan vise sin Kløgt og sin Begavelse end
ved at paavirke Mennesker saaledes, at de kommer til at leve under Tankens egen
Myndighed (Tillæg, Kap. 9).

Menneskene er ustadige. Faa er de, som lader sig lede af Fornuften. De
fleste er misundelige og hævnlystne. Men man gaar alligevel en gal Vej, naar man
skælder ud paa Menneskenes Fejl istedetfor at vejlede til Erhvervelse af værdifulde
Karakteregenskaber, og naar man knækker Modet istedetfor at styrke det. Man
skal med roligt Sind udholde Menneskenes Fornærmelser og arbejde paa alt, hvad
der kan fremme Endrægtighed og Venskab, og det er ét med at arbejde for, hvad
Retfærdighed, Billighed og Hæderlighed kræver.

Staten er en Støtte for denne Stræben efter at underordne det skiftende og
stridende under fælles faste Regler. Dens Virken beror paa to psykologiske Kendsgerninger.
Den første er, at en Følelse kun overvindes ved en anden Følelse, den
anden, at Frygt for selv at lide et Onde kan afholde fra at paaføre andre et Onde.
Til disse Kendsgerninger støtter Staten sig ved sine Straffetrusler, der kan udrette,
hvad Fornuften ikke formaar. Kun i Staten, ikke i den rene Naturtilstand, bliver der
derfor Tale om Brøde og Fortjeneste, om Uret og Ret (37 Schol. 2). Ud over det,
der kan fremtvinges ved Trusler, har Staten ingen Myndighed, og derfor ingen Ret,
og der bliver altsaa et Omraade tilbage, paa hvilket Mennesket under Livet i Staten
bevarer Kraft og Ret til at følge sin Selvopholdelsestrang. Det er af stor Betydning,
at Staten kender sine Grænser. Ti man kan ikke stole paa den Endrægtighed, der
kun tilvejebringes ved Indjagen af Frygt; det er ikke ved Vaaben, men ved Kærlighed
og Højsind, at Sjæle kan besejres. (Tillæg, Kap. 16).

En nærmere Udvikling af Forholdet mellem den enkelte og Staten giver \textsc{Spinoza}
i den teologisk-politiske Afhandling Kap. 16.

\subsection*{e. Sætning 38---58.}
\begin{center}\textit{Værdiskala for Glæder og Sorger.}\end{center}

\textit{α.} Idet \textsc{Spinoza} gaar over til en Vurdering af de forskellige Følelser i Henhold
til den Maalestok, han i det foregaaende har vundet, begynder han med en
biologisk og sociologisk Betragtning. Godt er alt, hvad der gør Legemet skikket til
at virke paa saa alsidig Maade som muligt, og hvad der bevarer det konstante Forhold
mellem Hvile og Bevægelse indenfor Legemet. Godt er fremdeles alt, hvad der
fører til Endrægtighed i Samfundet og derved gør det muligt, at Menneskene kan
leve efter sand Erkendelse.

Derefter gennemgaas saa de forskellige Følelser. Hovedsætningen er, at Glæde
% --- p. 79 ---
\setcounter{footnote}{0}%
\opage{79}altid i og for sig er god, Sorg (Ulyst) i og for sig ond\footnote[1]{Om Berettigelsen af denne Sætning se min \textit{Etik} VII, 1, hvor tillige nogle historiske Notitser om den er givne.}. Ti, hedder det i Tillægget
(Kap. 31), jo større Glæde vi føler, des større Fuldkommenhed gaar vi over til, og
des mere bliver vi delagtige i den guddommelige Natur. Men det er Glæde som
Totalfølelse (hilaritas), det gælder, --- Glæde som omfattende Menneskets hele Væsen.
I denne Betydning er det nu lettere at tænke sig Glæden end at finde den i Erfaringen.
Ti for det meste er der en enkelt speciel Følelse, der ytrer sig i overvættes
Grad og faar Sindet til at fæstne sig saaledes ved en enkelt Genstand, at alt
andet slet ikke existerer for det (44 Schol.). Hvad den Glæde angaar, der kun omfatter
enkelte Dele af Menneskets Væsen, altsaa partiel Glæde, „Kildren“ (titillatio)
(„Lyst“ i Ordets snevreste Betydning), da kan den overdrives saaledes, at det gaar
ud over hele Individets Fuldkommenhed. ---
Som partiel Følelse kan Sorg eller
Smerte være god. Den tyder jo paa, at der er Liv endnu, --- at den beskadigede
Del ikke er raadden. Skønt Anger og Skamfølelse f. Ex. i sig selv er Sorg, er dog
den, der er stedt i saadan Sorg, fuldkomnere end den frække, der ikke længere
har nogen Trang til Ære. (58 Schol.). En Følelse, i hvilken Sorg er et Element,
kan ikke kaldes god i og for sig. Dette gælder om Frygt og Haab, der begge forudsætter
Ulyst (som ventes at komme eller som ventes fjernet), og det gælder om
Følelser, der forudsætter forudgaaende Haab og Frygt, saaledes Tryghedsfølelse og
Glæden over overstaaet Sorg (securitas et gaudium) (47 Schol.).

Ved Totalfølelse tænker \textsc{Spinoza}, som det ses, paa Følelser, der vel angaar
Menneskets hele Væsen, men er simple, ensartede og usammensatte, --- ikke paa
saadanne, i hvilke Elementer og Eftervirkninger af flere indbyrdes forskellige, maaske
stridende Følelser gør sig gældende. Følelser af denne sidste Art vilde han i
hvert Tilfælde tillægge en ringere Værdi\footnote[2]{Se om saadan Totalfølelses Betydning \textit{Den store Humor} Kap. 1. Smlgn. ovenfor p. 55 og 68 f.}.

\textit{β.} Had mod Mennesker kan aldrig være godt. Det samme gælder Misundelse,
Spot, Foragt, Hævn. Had udryddes ikke ved Had, men ved Højsind. De, der besejres
paa denne Maade, viger med Glæde, ikke af Afmagt, men netop fordi deres
aandelige Kraft er øget. Følelser som Anger, Medlidenhed og Ydmyghed er ikke i
sig selv gode, men, som Menneskene nu engang er, kan de nævnte Følelser være
til Gavn og lette Overgangen til Visdom. Nedstemtheden, der er tilbøjelig til at
kaste sig bort (abjectio), er heller ikke i sig selv god, men er dog lettere at helbrede
end Selvovervurdering og Stolthed (superbia), med hvilken den iøvrig meget
vel kan være beslægtet (57 Schol.), idet den især er Sorg over egen Afmagt formedelst
Sammenligning med andres Magt og Lykke, hvilket igen betinger, at den
let gaar over til Misundelse. --- Disse sidste Bemærkninger minder om \textsc{Nietzsche}'s
Karakteristik af de ulykkelige og deres „Slaveopstand“\footnote[3]{Smlgn. \textit{Moderne Filosofer} p. 140.}. Men der er en principiel
Modsætning i Maalestokken: \textsc{Nietzsche}'s Maalestok er Magten i og for sig, \textsc{Spinoza}'s
% --- p. 80 ---
\setcounter{footnote}{0}%
\opage{80}hentes fra de aandelige Goder, der kan blive tilgængelige for alle, og fra hvilke
Stolthed og Ydmyghed hver paa sin Maade kan udelukke.

\textit{γ.} Ved Slutningen af denne Værdiskala af Glæder og Sorger minder \textsc{Spinoza}
om, at hele den Vurdering, han her har anstillet, kun gælder ud fra det Synspunkt,
som menneskelig Stræben angiver. Naturens Love, tilføjer han, angaar Tilværelsens
almindelige Orden, af hvilken Mennesket kun er en Del. Det har ikke været hans
Opgave at berette om Menneskers Fejl og urimelige Færd, men at paavise Tingenes
Natur og Egenskaber. Det er stadig hans Mening at betragte menneskelige Følelser
og deres Egenskaber paa samme Maade som andre Naturfænomener. Og selv om
de menneskelige Følelser just ikke vidner om Menneskers Kløgt og Visdom, saa
vidner de dog om Naturens Magt og Visdom, ikke mindre end andre Fænomener,
vi glæder os ved at undersøge (57 Schol.). For \textsc{Spinoza} er der ingen Splid mellem
Kosmologi og Etik.

\subsection*{f. Sætning 59---66.}
\begin{center}\textit{Værdiskala for Drifter.}\end{center}

\textit{α.} Skønt \textsc{Spinoza} klart har set, at der er noget kunstigt i at skelne mellem
Følelse og Drift, da Driften bestemmes ved Forestillingen om Følelsens Aarsag (III,
12---15; 57 Dem.), siger han dog (IV, 58 Schol.), at han nu, efterat have talt om
Værdien af Glæder og Sorger, vil gaa over til at give en Vurdering af Drifterne.

De Handlinger, som udspringer af passive Følelser, kan ogsaa øves af sand
Erkendelse. At handle efter sand Erkendelse (ratio) er nemlig ikke andet end at
handle efter vor Naturs inderste Væsen. Naar Sorg, og undertiden ogsaa Glæde, er
af det onde, er det, fordi de er passive Følelser, Udtryk for hæmmet Virkekraft.
Falder Hæmningen bort, saa Selvhævdelse og Overgang til større Virkekraft bliver
mulig, udføres de Handlinger, som før skyldtes passiv Følelse, paa langt fuldkomnere
Maade. Som Exempel anføres Forskellen mellem blind Hævn og besindig Straf (59).

Det er af afgørende Betydning, om Driften kun udspringer af en enkelt Del af
vort Væsen, eller om den tager Hensyn til alle Sider i vor Natur. Fremdeles ---
om den udspringer af en Aarsag, der kun virker i det nærværende Øjeblik, eller
om den tager Hensyn til Fremtiden. Paa den sande Erkendelses Standpunkt falder
hele Forskellen mellem Fortid, Nutid og Fremtid bort. Forskellen mellem Erkendelsens
Trin beror væsentlig paa, at Tidsforskelle og Tidsforhold spiller en Rolle
paa de laveste Trin, men falder bort paa de højere Trin, hvor man ser Tingene
„under Evighedens Synspunkt“. (Smlgn. ovenfor p. 47---54). Erkendelsen af godt og
ondt hører ikke hjemme paa Erkendelsens højeste Trin, netop fordi den forudsætter
Tidsforskelligheders Indflydelse. Vi ser jo ikke lige klart paa Fortid, Nutid og Fremtid
med Hensyn til, hvad der er at foretrække. Derfor er den Erkendelse af godt og
ondt, vi har, selv om den er rigtig, dog kun abstrakt og generel og kan ikke i det
enkelte faa den Indflydelse, det i og for sig ligger i dens Væsen at udøve. Erkendelse
af det onde er nu altid en ufuldstændig (uadækvat) Erkendelse, da den, som
% --- p. 81 ---
\setcounter{footnote}{0}%
\opage{81}betinget ved Overgang til ringere Virkekraft, ikke finder sin Forklaring ved vort
eget Væsen, men kun forstaas ved Afhængigheden udadtil. Havde vi kun fuldkommen
Erkendelse, vilde vi intet Begreb om ondt danne. Den, der ledes af sit
eget Væsen, d. v. s. af selvstændig dannede klare Tanker, gør stedse, hvad han vil,
nemlig, hvad der for ham er det ypperste i Livet. Han kan i Sandhed kaldes et
frit Væsen (60---66).

\textit{β.} Ser vi nu tilbage paa den hele Udvikling om ufuldkommen og fuldkommen
Handlen, som \textsc{Spinoza} giver i fjerde Bog, viser det sig, at hvad der i etisk Henseende
er irrationelt, ifølge hans Opfattelse dels skyldes Forskelle i Følelsernes Styrkegrad
(IV, 1---18), dels Forskelle mellem indre og ydre Aarsager (19---28), dels Forskelle i
Henseende til Tid (59---66). Tilsidst kommer vi tilbage til den Kendsgerning, at der
overhovedet gives Forskelligheder i Tilværelsen, hvilket gør, at vor Tænken kommer
til at bestaa i Sammenligning (se Fortalen til fjerde Bog), ikke i Skuen af det indre
Baand mellem Tingene. Saalænge vi ikke ser alt i ét, handler vi mere eller mindre
i Blinde. Idealet (exemplar) er det frie Menneske, til hvis Karakteristik \textsc{Spinoza} nu
gaar over.

Skønt der altsaa hos \textsc{Spinoza} gør sig en skarp Modsætning gældende mellem
Mennesket, som Erfaringen viser os det, omtumlet i Blindhed og i Illusioner, og det
ideale Menneske, Forbilledet for vor Tænken og Handlen, føler han dog en dyb Uvillie
mod dem, som han snart kalder de overtroiske, snart Melankolikerne, snart Satirikerne,
som bedre forstaar at revse Lasterne end at vise Vejen til sand Frihed, --- som snarere
lærer Menneskene at sky det onde end at elske det gode, --- og som virker ved at
indgyde Frygt istedetfor at lede Menneskene fremad gennem sand Erkendelse. (63
Schol. og Tillæg Kap. 13 og 31, smlgn. allerede 35 Schol. og 45 Schol.).

\subsection*{g. Sætning 67---73.}
\begin{center}\textit{Det frie Menneske.}\end{center}

\textit{α.} Det frie Menneske betegner den højeste Udfoldelse af den menneskelige
Natur, et Liv efter dennes egne Love. Som det klart udtrykkes i Tractatus politicus
(II, 7. 11), er Frihed en Dyd eller en Fuldkommenhed; derfor kan Mennesket kun
kaldes frit, naar det følger den sande Erkendelse, existerer og handler ud af den
menneskelige Naturs indre Love, af Aarsager, der forstaas ud fra Mennesket selv.
Frihed i denne Betydning ophæver ikke Nødvendigheden, men forudsætter den.

Det frie Menneske tænker ikke paa Død, men paa Liv. Han higer efter, hvad
godt er, for dets egen Skyld, ikke af Frygt for Døden. Hans Visdom er en Gransken
af Livet (meditatio vitæ). Forskellen mellem godt og ondt existerer ikke for
ham; det er den rigtige Tanke i den gamle Fortælling om Forbuddet mod at spise
af Træet til at lære godt og ondt at kende. Hans Aandskraft viser sig lige saa meget
i at afværge Farer som i at overvinde dem. Da han handler af sand Erkendelse,
øver han ikke Svig. Kunde sand Erkendelse tillade Svig som almindelig Regel,
vilde hele Samfundet blive umuligt. Det frie Menneske lever i den Overbevisning,
% --- p. 82 ---
\setcounter{footnote}{0}%
\opage{82}at alt følger af den guddommelige Naturs Nødvendighed, og naar derfor noget staar
for ham som ondt, skrækkeligt eller uretfærdigt, vil han antage, at Grunden er
dunkel og brudstykkeagtig Erkendelse, og han vil stræbe efter at forstaa Tingene,
som de er i sig selv, uden at lade sig paavirke af Had og Haan, af Vrede og
Hovmod.

Denne Karakteristik er, som \textsc{Spinoza} selv bemærker, en videre Udførelse af,
hvad der allerede i tredie Bog (III, 59 Schol.) beskreves som Aandskraft (fortitudo)
i dens to Hovedformer som Livsmod (animositas) og som Højmod (generositas). At
Mennesket hæver sig over passive Følelser, som Medlidenhed og Sorg, fører derfor
ingenlunde til Selvoptagethed eller kold Ligegyldighed. Paa Højdepunktet af aandelig
Udvikling gør Modsætningen mellem eget og andres Vel sig ikke mere gældende,
fordi der er Kraft og Sind til at arbejde for begge, og --- kan vi føje til i \textsc{Spinoza}'s
Aand --- fordi der i ens egen Personligheds fulde Udvikling ydes det bedste Bidrag
til Menneskesamfundets sande Tarv. Statens eget Maal er Frihed. (Tract. theol. polt. % PRINTED AS IS: polt. (polit.)
Kap. 20).

\textit{β.} I \textsc{Spinoza}'s „frie Menneske“ fremtræder den mest gennemklarede Form for
den Tanke, der, som allerede ovenfor (p. 61) omtalt, bærer hele Renaissancens Etik.
Denne Tanke har udviklet sig under Indflydelse af Stoicismen. Men der er en betydningsfuld
Forskel mellem den antike Stoicisme\footnote[1]{Smlgn. om Stoicismen min Afhandling \textit{Hedenske Sandhedssøgere} (Mindre Arbejder. Første Række) og Afhandlingerne om \textit{Epiktet og Marcus Aurelius} (Mindre Arbejder. Tredie Række).} og dens Efterfølgere paa Renaissancens
Tid. Medens Selvhævdelsen hos Stoikerne var en Tilflugt, et Udtryk for
Resignation, fik den hos Renaissancens Tænkere (fra \textsc{Telesio} gennem \textsc{Bruno}, \textsc{Descartes}
og \textsc{Hobbes} til \textsc{Spinoza}) en aktiv og positiv Karakter, og den udsprang af
Tillid til Livet og Begejstring for dets store Love og næredes ved den Udvidelse af
Naturforstaaelse og Verdensanskuelse, som den nye Forsken bragte.

Den Sætning (68), at hvis Menneskene fødtes frie, vilde de ikke danne Begreber
om godt og ondt, er af stor Betydning. Af den følger nemlig, at Bevisbyrden paahviler
den, der anser det for nødvendigt at anstille moralsk Bedømmelse af menneskelige
Handlinger. Det etiske Problem ligger ikke blot i det Spørgsmaal, hvorvidt
og hvorledes etiske Domme kan begrundes, men først og fremmest i det Spørgsmaal,
om etiske Domme overhovedet er nødvendige. Det, at et Menneske fælder Domme
om godt og ondt, er for \textsc{Spinoza} et Tegn paa, at det ikke lever Livet ud af indre
Kraft (med Livsmod), og at det heller ikke (med Højmod) formaar at hjælpe andre
til at leve Livet ud af dem selv. Som det hedder i en allerede ovenfor citeret Udtalelse:
„Der er ingen Maade, paa hvilken et Menneske bedre kan vise sin Kløgt og sin
Begavelse, end ved at paavirke andre saaledes, at de kan leve under Tankens egen
Myndighed“. (Tillæg. Kap. 9).

I denne Tanke indeholdes den inderligste Forbindelse mellem individuel og
social Etik, som overhovedet er mulig\footnote[2]{Smlgn. herom min \textit{Etik}, III, 10 og XIII, 6.}. Ved at give efter for Lysten til at moralisere
kommer man let til at lægge Hindringer i Vejen for det, der netop skulde være
% --- p. 83 ---
\setcounter{footnote}{0}%
\opage{83}Maalet for al Moral: Udviklingen af selvstændige Personligheder. Det at vore Handlinger
bedømmes --- af os selv eller af andre --- er kun etisk berettiget, naar det er
Led i en Opdragelse til Frihed. I \textsc{Rousseau}'s pædagogiske Evangelium var det derfor
ogsaa en Hovedsætning, at Opdragelse først og fremmest maa foregaa indirekte,
eller, som han udtrykker det, være „negativ“. Det er selve Personlighedsprincipet,
som kræver dette.

I sin Begrundelse af, at det frie Menneske ikke handler med Svig (72), benytter
\textsc{Spinoza} en Tankegang, der senere spiller en stor Rolle hos \textsc{Kant} ved Fastsættelsen
af det almindelige moralske Princip. \textsc{Kant} hævder, at det ikke kunde opstilles som
almindeligt Princip, at man har Lov at tilegne sig betroet Gods eller at give et Løfte
med den Hensigt ikke at holde det\footnote[1]{Tankegangen er gammel og forekommer allerede hos \textsc{Platon} (Staten I, p. 362 C).}. \textsc{Kant} saa ikke, at der til Grund for denne
Tankegang ligger den Forudsætning, at Samfundet skal bestaa, og at det kun kan
bestaa ved gensidig Tillid\footnote[2]{Smlgn. \textit{Den nyere Filosofi's Historie\textsuperscript{2}}, p. 72---84.}. \textsc{Spinoza} derimod gaar ud fra Samfundet som et historisk
Faktum, hvis Betingelser det gælder at finde. Den Fornuft, det frie Menneske
ledes af hos \textsc{Spinoza}, indbefatter ogsaa Betingelserne for Samfundets Bestaaen og
erklærer det for absurd, at det kunde bestaa uden gensidig Tillid (72 Schol.). Trods
al \textsc{Spinoza}'s Spekulation mærkes det stedse i hans Etik, som i de andre Dele af
hans Filosofi, at Erfaring for ham er det første Trin af Erkendelse.

% --- p. 84 ---
\setcounter{footnote}{0}%
\opage{84}\section*{Femte Bog.}
\subsection*{a. Fortale og Axiomer.}

\textit{α.} Naar \textsc{Spinoza} begynder Fortalen til femte Bog med at sige, at han nu
„gaar over til Etikens anden Del“, betragter han vel fjerde Bog som den første Del,
da Etik i snævrere Forstand begynder med den. Men naar han siger, at denne
anden Del skal omhandle „Vejen, der fører til Frihed“, staar dens Indhold ikke i
samme Forhold til fjerde Bog som Frihed til Trældom. Allerede i fjerde Bog har
han vist Muligheden af en Udvikling fra Afhængighed af udefra vakte Følelser og
Drifter til aandelig Frihed, og Bogen naaede sit Højdepunkt i Skildringen af det frie
Menneske. Fjerde Bog har underligt nok Overskriften „Om den menneskelige Trældom“,
men denne Overskrift passer ikke til Indholdet. Ja, allerede i tredie Bogs
Slutning (Sætning 58---59) vises jo Muligheden af en Overgang fra passive til aktive
Følelser, altsaa fra Trældom til Frihed.

Fjerde Bog burde være sluttet med Sætning 66, og Sætning 67---73 (Skildringen
af det frie Menneske) burde have udgjort én Bog sammen med de tyve første Sætninger
i femte Bog, der handler om Midlerne mod Afhængighed af Følelserne (de
remediis affectuum), altsaa om Arbejdet paa at blive et frit Menneske. Den øvrige
Del af femte Bog udvikler, hvorledes den, der har naaet den sande Erkendelse, er
uafhængig af Tiden og staar som en evig Form for Gud eller Substansen. Allerede
i fjerde Bog (74 Schol. og Tillæg Kap. 4---5) var han kommen ind paa den intuitive
Erkendelse som det højeste Maal og som det, der bestemmer Værdien af alt andet.

\textsc{Joachim}\footnote[1]{\textit{The Ethics of Spinoza}, p. 268.} bestemmer Forholdet mellem fjerde og femte Bog saaledes, at fjerde % PRINTED AS IS: the mark at Joachim reads ²), the note below ¹) (the page's only note)
Bog beskriver Menneskets Erkendelse og Liv paa et Stadium, hvor fuldkomne Tanker
vel leder dets Færd, men ikke fylder hele dets Sjæl. Denne Opfattelse strider dog
afgjort mod Skildringen af det frie Menneske i Slutningen af fjerde Bog. Hvad der
kommer til i femte Bog er den Konsekvens, at det frie Menneske er hævet over
Tiden, ikke blot tænker under Evighedens Synspunkt, men selv er evig. I Slutningen
af 20 Schol. siger \textsc{Spinoza}, at han, efter at have afsluttet Betragtningen af
alt, hvad der angaar det nærværende Liv, vil gaa over til, hvad der angaar Sjælens
Bestaaen afset fra Legemet (sine relatione ad corpus). Det er, bortset fra de tyve
første Sætninger, mere en religiøs-mystisk end en etisk Betragtningsmaade, der er
karakteristisk for femte Bog.

% --- p. 85 ---
\setcounter{footnote}{0}%
\opage{85}\textit{β.} Fortalen er væsentlig en Kritik af \textsc{Descartes}' Psykologi, fremkaldt ved, at
\textsc{Descartes} ligesom Stoikerne uden videre tilskriver Mennesket Evne til at beherske
Følelserne. Som \textsc{Brochard} bemærker i sin Sammenligning mellem \textsc{Descartes}' og
\textsc{Spinoza}'s Psykologi\footnote[1]{\textit{Le traité des passions de Descartes et l'Éthique de Spinoza} (i hans Bog Études de philosophie ancienne et de philosophie moderne, p. 330).}, dvæler \textsc{Spinoza} kun ved de Punkter, i hvilke han er uenig
med \textsc{Descartes}, og ikke ved Overensstemmelserne; saaledes nævner han ikke den
Vægt, \textsc{Descartes} (ligesom allerede Stoikerne) lagde paa Forstaaelse af Følelsernes
Natur og i det hele paa Forestillingernes Indflydelse paa Følelseslivet.


\textsc{Spinoza}'s Kritik angaar fire Punkter:


1. \textsc{Descartes} giver ingen Forklaring af, hvorledes Sjæl og Legeme (Tanke og
et Udstrækningsfragment) kan være forbundne. Han antager en saadan Modsætning
mellem dem, at han maa beraabe sig paa Gud, hele Universets Aarsag, for at kunne
løse Gaaden. Men det er ingen videnskabelig Løsning, ingen Forklaring ved en
speciel Aarsag (per proximam causam)\footnote[2]{Proxima causa hos \textsc{Spinoza} svarer til vera causa hos \textsc{Kepler}.}.

2. Det bliver uforstaaeligt, hvorledes det kan afgøres, om Sjælen eller „Livsaanderne“
er stærkest. Hvem af dem formaar i det afgørende Øjeblik at give Koglekirtlen
det afgørende Stød? Hvorledes kan de to Kræfter, der her skal mødes,
sammenlignes? Villen og Bevægelse er ikke kommensurable.

3. Anatomien godtgør, at Koglekirtlen ikke har sin Plads midt i Hjernen, som
\textsc{Descartes} mente. Den kan derfor ikke bevæges saa let og paa saa mange Maader,
som han forudsatte. Desuden fører ikke alle Nerver ud til Hjernehulhederne\footnote[3]{Disse og lignende anatomiske Indvendinger var gjort af \textsc{Thomas Bartholin} (1651) og \textsc{Steno} (1665). Se herom \textit{Den nyere Filosofis Historie\textsuperscript{2}}, I, p. 517. --- \textsc{Steno}'s Kritik kan \textsc{Spinoza} have kendt fra sin personlige Omgang med ham i Aarene 1661---1663. \textsc{Steno}'s Afhandling om Hjernens Anatomi er oversat paa Dansk af Dr. \textsc{Maar}.}.

4. Der maa kræves en rent psykologisk Forklaring af Herredømmet over Følelserne.
Da al aandelig Kraft betinges ved Forstaaelse alene, maa man, hvor Talen er
om Midler mod Trældom under Følelser, lægge Forstaaelse til Grund. Men istedetfor
at give en rent psykologisk Forklaring af, hvorledes Forstaaelsen kan udvikle sig
og vinde Magt, beskriver \textsc{Descartes} en rent imaginær Kamp mellem Sjæl og Hjerne.
Den rette Fremgangsmaade vilde have været at beskrive de sjælelige Fænomener i
deres indbyrdes Sammenhæng og derpaa raadspørge Anatomien og Fysiologien.
Ifølge \textsc{Spinoza}'s Opfattelse har enhver Kamp mellem Motiver i Sjælen sit samtidige
Udtryk i Forholdet mellem forskellige Hjerneprocesser. Der kan da ikke være Tale
om, at Sjælen skulde kunne tage en Beslutning, medens samtidig Hjernen skulde
kunne indlede en Bevægelse i helt anden Retning, % PRINTED AS IS: the paragraph ends with a comma

Hvad den anden Indvending angaar, beror den paa de Grunde og Synsmaader,
der førte \textsc{Spinoza} til hans Attributlære. (Se ovenfor p. 38---45). Men han overser, at
han selv antager Proportionalitet mellem sjælelige og legemlige Fænomener, saa at
han kommer til at staa overfor Vanskeligheden ved at paavise en fælles Maalestok
for de to Grupper af Fænomener. Vel er der den Forskel, at hos \textsc{Descartes} skulde
% --- p. 86 ---
\setcounter{footnote}{0}%
\opage{86}den ene Art Fænomener paa et vist Punkt danne Fortsættelse af den anden (Koglekirtlens
Bevægelse afløse Beslutningen, eller Sansefornemmelse og Følelse afløse
Koglekirtlens Bevægelse) og danne en Række, indenfor hvilken Kontinuiteten paa et
enkelt Punkt var brudt, medens hos \textsc{Spinoza} to hver for sig kontinuerlige Rækker
forløber proportionalt med hinanden, saaledes at Menneskets Væsen netop lægger
sig for Dagen i denne Proportionalitet. Men Vanskeligheden bliver da den, at finde
et Led i den ene Række, der svarer til visse Led i den anden Række, en Vanskelighed,
der bliver des større, jo mere Vægt der lægges paa Kontinuiteten indenfor
hver Række\footnote[1]{Jeg har berørt dette Punkt i \textit{Filosofiske Problemer} (Universitetsprogram 1902), p. 23---25.}. For at raade Bod paa denne Vanskelighed maa der opstilles et
Begreb om psykisk Arbejde, der kan svare til Begrebet om materielt Arbejde\footnote[2]{Smlgn. herom \textit{Den menneskelige Tanke}, p. 4---6.}, og
dertil naaede \textsc{Spinoza} ikke, skønt han saa ofte taler om Virkekraft baade paa det
aandelige og det materielle Omraade. Det er hans kartesianske Naturfilosofi, der
hindrer ham deri. Først \textsc{Leibniz} begyndte at se klarere her\footnote[3]{Se \textit{Den nyere Filosofis Historie\textsuperscript{2}}, I, p. 357---363.}.

\textit{γ.} Der opstilles ingen nye Definitioner i femte Bog, kun to Axiomer. --- Det
første fastslaar, at naar to forskellige Virksomheder (actiones) opstaar hos samme
Person (subjectum), maa der ske Forandringer i dem begge, eller dog i den ene,
indtil Modsætningen er hævet. Det er aabenbart en Erfaringssætning, der her skal
opstilles. Men det er klart, at denne Sætning hviler paa den Forudsætning, at det
ikke er ligegyldigt for de to Virksomheder, at de foregaar hos samme Subjekt. Der
forudsættes en fundamental Tendens til at sammenarbejde de forskellige Elementer
og Processer, der kan forekomme indenfor samme Helhed. \textsc{Spinoza} vilde ved at
fremdrage denne Forudsætning være naaet til et Begreb om Arbejde, der vilde have
kunnet være frugtbart for hans Filosofi. Axiomet benyttes i Sætning 7, hvor der
tales om Modsætningen mellem saadanne Følelser, der er knyttede til Tanker om
Love, altsaa om „det, der gælder paa lige Vis i Delen som i Helheden“, og saadanne,
der er knyttede til en enkelt, isoleret Ting. Der drages den Slutning af Ax. 1, at
Følelser af sidste Art i Længden maa vige for Følelser af første Art. Hvad \textsc{Spinoza}
her slutter, er vel ogsaa rigtigt, skønt det kan tage meget lang Tid, inden en utaalmodig,
om et enkelt Emne koncentreret Følelse viger for Følelser, der er bestemte
ved Forhold, der atter og atter gør sig gældende. Følelseslivet udvikler sig nu engang
langsommere end Forestillingslivet, enkelte Opbrusninger undtagne. Men selv
om \textsc{Spinoza} har Ret i den specielle Anvendelse af Axiomet, har han neppe Ret i at
opstille Axiomets Indhold som eneste Mulighed. Det kan være, at der svinges mellem
de to Tendenser, eller at der vedbliver at bestaa en Spændingstilstand mellem dem,
eller at den ene faar Eneherredømmet, eller at der opstaar en Totalfølelse med en
efter Modsætningernes indbyrdes Forhold forskellig Klangfarve\footnote[4]{Smlgn. \textit{Den store Humor}, p. 39, 73, 88 ff.}. Det er dogmatisk
at antage en saa simpel mekanisk Udjævning, som \textsc{Spinoza} (og nogle moderne

% --- p. 87 ---
\setcounter{footnote}{0}%
\opage{87}Psykologer med ham) tænker sig. Der kan bestaa flere sjælelige Modsætninger, end
man i Regelen tænker sig\footnote[1]{Smlgn. \textsc{August} \textsc{Wimmer}: \textit{Psykogene Sindssygdomme} (Jubilæumsskrift for St. Hans Hospital, p. 107), hvor dette begrundes ad patologisk Iagttagelses Vej.}.

Det andet Axiom --- at Virkningens Kraft bestemmes ved Aarsagens Kraft ---
bringer intet nyt. Det begrundes da ogsaa ved en Henvisning til III, 7, der igen
begrundes ved I, 29. 36.

\subsection*{b. Sætning 1---20.}
\begin{center}\textit{Midler overfor Følelserne.}\end{center}

Det er ifølge \textsc{Spinoza} egentlig ikke Følelserne selv, der behøves Midler imod.
Hvad det gælder om, er at ophæve den passive Karakter, nogle Følelser har, ved
at omsætte dem til aktiv Form. Gennem en rent psykologisk Udvikling viser \textsc{Spinoza}
Muligheden heraf, som iøvrigt allerede var antydet i tredie og fjerde Bog.

Ifølge \textsc{Spinoza}'s Attributlære og den deri formulerede Identitetshypotese kan
og maa denne Mulighed undersøges baade fra den psykiske og fra den fysiske Side
og det er ligegyldigt, fra hvilken Side vi begynder. I anden, tredie og fjerde Bog
begyndte \textsc{Spinoza} i sin Bevisgang fra den fysiske Side, idet han dog stadig (se f. Ex.
III, 12 Dem.) fastholder den tvesidige Betragtning. I den første Del af femte Bog
(1---20) begynder han med den psykiske Side, aabenbart fordi man, hvad det her
omhandlede Spørgsmaal angaar, kender denne Side bedst. Senere (fra Sætning 21
af) gaar han ud fra den fysiske Side, da han her mener at have klare og simple
Forudsætninger at gaa ud fra. Det er ikke, som man undertiden har ment, selvmodsigende,
men det er fuldt konsekvent, naar \textsc{Spinoza} saaledes snart ser Sagen
fra den ene, snart fra den anden Side, og enhver videnskabelig Behandlingsmaade
af de Spørgsmaal, der her er Tale om, gaar i Virkeligheden frem paa denne Maade,
hvad enten man er sig det bevidst eller ikke.

De Veje, \textsc{Spinoza} anviser til Herredømme over eller Metamorfose af de passive
Følelser, er tre.

\textit{α.} Det gælder om at bringe Følelserne ind i nye Forestillingsforbindelser derved,
at Tanken om deres ydre Aarsag afløses af andre Tanker (2). Af særlig Betydning
er her Tankerne om Tingenes fælles Egenskaber og deres almene Love, ---
om det, der i lige Grad gør sig gældende i det enkelte og i det hele og stedse er
nærværende (7). Ved at knyttes til saadanne Tanker har Følelsen ikke længere en
enkelt og forsvindende Grund, men betinges ved sand Erkendelse (ratio) og kan
forstaas som Led i Tingenes store Sammenhæng. Dunkle Tanker afløses af klare
Tanker, og Følelsen faar nu en anden Karakter. Saaledes gaar f. Ex. den Hersketrang
eller Stolthed, der vil tvinge andre til at leve efter vort Tykke, over til det
Højmod, der stræber at udvikle andre til at leve efter klar Forstaaelse. Det er en
og samme Trang, som ligger til Grund i passive og i aktive Tilstande, men der er
foregaaet en Metamorfose ved Overgangen til aktiv Form formedelst klar Erkendelse
% --- p. 88 ---
\setcounter{footnote}{0}%
\opage{88}af Forholdene (4 Schol.). Sorgen ophører at være Sorg, naar dens Aarsag klart erkendes
(15 Dem.; 18 Schol.): den afløses nemlig af Erkendelsesglæden. En saaledes
omformet Følelse vil betyde en større Kraft i Sindet end en Følelse, der er knyttet
til en enkelt Genstand. Fast Tankesammenhæng giver ogsaa Sammenhæng i Følelseslivet.
Der kræves større Kraft til at hæmme Følelser, som er ordnede og forbundne
efter klar Erkendelse, end til at hæmme ubestemte og skiftende Følelser
(10 Schol.). Tanken om det evige og uforanderlige kan tilsidst i højeste Grad begrænse
de passive Følelsers Herredømme (20 Schol.).

Overgangen til klar Erkendelse er for \textsc{Spinoza} et specielt Tilfælde af Motivforskydning.
I og for sig kan Motivforskydning ogsaa føre til, at en passiv Følelse
afløses af en anden passiv Følelse. Saaledes kan Glæde ved en Ting eller Attraa
efter den føre til Misundelse overfor den, der besidder denne Ting (III. 32). Men
\textsc{Spinoza} har utvivlsomt Ret i, at en Forskydning er mulig fra Glæde ved en Ting
til Forstaaelse af denne Glæde og dens Aarsag, saaledes at derved begge, baade
Følelsen og dens Aarsag (Genstand), indordnes som Led i en hel Natursammenhæng,
hvorved Hildethed og Betagethed af det enkelte Emne ophører. Den klare
Forstaaelse er \textsc{Spinoza}'s store Universalmiddel, og han kan i denne femte Bog saa
meget des bedre lægge Hovedvægten paa dette Middel, som han her --- efter de Svingninger,
som ovenfor er paaviste i tredie og fjerde Bog --- afgjort vender tilbage til
sin intellektualistiske Opfattelse, ifølge hvilken „Aandens Kraft alene bestaar i Erkendelse“
(mentis potentia sola intelligentia definitur) (se Fortalens Slutning og 20
Schol.). Kun gennem klar Forstaaelse bliver Sjælen helt og fuldt sig selv.

\textit{β.} Men Udviklingen tager Tid. Sand Erkendelse naas kun Skridt for Skridt.
Det gælder at benytte Mellemrummene mellem stærke Sindsbevægelser til Besindelse.
\textsc{Spinoza} taler her af egen Erfaring. I den Gæringsperiode, som han beskriver
i Begyndelsen af Afhandlingen om Forstandens Forbedring, følte han sig friere, naar
han blot kunde hengive sig til alvorlig Eftertanke (modo possem serio deliberare!).
Og de Mellemrum, i hvilke dette skete, blev efterhaanden hyppigere og længere.
Derfor formaner han nu (V, 10 og 20 Schol.) til at benytte Mellemrummene til at
indpræge sig visse Livstanker (vitæ dogmata), som der kan blive Brug for, naar
passive Følelser igen truer med at overvælde os. Saadanne Livstanker er f. Ex., at
Had ikke overvindes ved Had, men ved Højsind, og at Menneskenes Færd, ligesom
alt i Naturen, er underkastet nødvendige Love. Ved at indpræge sig saadanne
Tanker kan man vinde frem, selv om det ikke sker uden Svingninger (sine animi
fluctuatione). Derimod er det en uheldig Vej, man gaar, naar man f. Ex. trøster
sig over ikke at være rig ved at betænke, hvor ofte Rigdom misbruges, og hvor
lastefulde de rige ofte er. Derved volder man kun ny Sorg hos sig selv. Nej, det
ledende Motiv skal være Trang til aandelig Frihed (amor libertatis). Man ledes da
til at forstaa værdifulde Karakteregenskabers Oprindelse, og Sindet vil fyldes af
Erkendelsesglæde.

\textit{γ.} Sit Højdepunkt naar denne Erkendelsesglæde, naar Forstaaelsen af vore
Følelser og deres Aarsager gaar tilbage til det sidste Grundlag for vor Erkendelse,
% --- p. 89 ---
\setcounter{footnote}{0}%
\opage{89}Tanken om Gud (Substansen, Naturen). I dette Grundlag har Erkendelsesglæden
sin sidste Aarsag, og da al Glæde bliver til Kærlighed, naar den forbindes med Forestillingen
om sin Aarsag (III, 53 og V, 15 Dem.), vil der opstaa Kærlighed til Gud,
--- en Kærlighed, der ikke kræver Genkærlighed og ikke giver Anledning til Skinsyge
eller Misundelse, men som i højeste Grad kan optage Sindet, saalænge Livet
varer (14---20).

\subsection*{c. Sætning 21---40.}
\begin{center}\textit{Aandens Evighed.}\end{center}

\textit{α.} I dette Afsnit tager \textsc{Spinoza} igen sit Udgangspunkt i den fysiske (fysiologiske)
Side af Menneskets Væsen. Han skelner mellem Legemet i dets specielle eller
aktuelle Existens, med bestemt Existenstid og Tidsvarighed, og Legemets Natur som
indeholdt i den evige Naturorden, der er udtrykt i det konstante Forhold mellem
Hvile og Bevægelse i Universet, som jo for ham er et eneste stort Individ, i hvilket
der ingen Forandring sker, saa store Forandringer der end foregaar i de enkelte
Dele. (Se ovenfor p. 28---30; 43). I et af sine Breve siger han, at vore Legemer ikke
skabtes, men existerede forud for vor Fødsel, skønt i andre Former (qvamvis alio
modo formata) (Ep. 4, smlgn. Fortalen til anden Del af den korte Traktat og L. M.s
Fortale til \textsc{Spinoza}'s Fremstilling af \textsc{Descartes}' Filosofi). Forsaavidt det enkelte
Legemes Natur er indeholdt i den materielle Naturorden og gennem denne i Udstrækningens
Attribut, er det en evig Modus, Genstand for Erkendelse „under Evighedens
Synspunkt“. Paa tilsvarende Maade er ogsaa Sjælen indeholdt i den aandelige
Naturorden og gennem den i Tænkningens Attribut som evig Modus, Genstand
for Erkendelse „under Evighedens Synspunkt“. Kun saalænge den bestemte Tidsexistens
varer, har Sjælen Sansning og Erindring. Der gives i Gud (Substansen, Naturen)
en Tanke, der svarer til Legemets Existens under alle Former (før og efter
Fødselen og efter Døden)\footnote[1]{„Vort Legeme frembød før Fødselen et andet Forhold mellem Hvile og Bevægelse [end efter] og vil ligeledes efter Døden bestaa under et andet Forhold mellem Hvile og Bevægelse [end før Døden]. Men ikke desto mindre var der da [ɔ: før Fødselen], og vil der siden [efter Døden] ligesaa vel som nu [i Levetiden] være en til Legemet svarende Tanke i Substansen som tænkende, kun ikke den samme, da Legemet vil være stedt i et andet Forhold mellem Bevægelse og Hvile“. \textit{Kort Traktat}. II. Fortale (Note 10).}. Ligesaalidt som Legemet bliver derfor Sjælen til intet
ved Døden (21---23).

Begrundelsen sker ved Hjælp af II, 8 Cor., hvor der allerede (se ovenfor p. 41 f)
skelnedes mellem en dobbelt Existensform for alle enkelte Ting (res singulares), ---
den ene, forsaavidt Tingene er indeholdte i Guds Attributer, den anden, forsaavidt
de tillige har Existens i Tiden. Denne Sætning tager \textsc{Spinoza} nu op, ikke blot i
Beviset for 21 og 23, men bestemtere i 29 Schol. (hvor der vises tilbage til II, 45,
der igen begrundes paa II, 8 Cor.). Hin førstnævnte Existensform er evig, ligesom
Forholdet mellem Grund og Følge er det. Der har jo aldrig været en Tid, hvor
Grunden (Præmisserne) var, men Følgen (Konklusionen) ikke. Naar to Chorder
% --- p. 90 ---
\setcounter{footnote}{0}%
\opage{90}skærer hinanden indenfor en Cirkel, er de Rektangler, der dannes af de afskaarne
Stykker, ligestore. Det er en evig Sandhed. I hvert enkelt Tilfælde er kun visse
bestemte Chorder dragne; men de andre existerer alligevel som indeholdte i Cirkelens
Begreb (II, 8 Cor.). Det er atter her \textsc{Spinoza}'s Identificering af Grund og Aarsag,
der gør sig gældende. Som tidsbestemt er Mennesket Led i en Aarsagsrække, i
hvilken ethvert Led paa én Gang er Virkning af det foregaaende og Aarsag til det
følgende; men i sit evige Væsen er det en utilbleven, uforanderlig og uforgængelig
Modus for Substansen. Modsætningen mellem disse to Existensformer vilde staa
som meget grel, dersom det ikke for \textsc{Spinoza} forholdt sig saaledes, at Tidsforskellen
er aldeles uvæsentlig for Aarsagsforholdet, medens det afgørende er det tidløse Forhold
mellem Grund og Følge.

\textsc{Spinoza} minder her om \textsc{Platon}. Den nyplatoniske Tankegang har jo ogsaa
paavirket ham gennem jødisk og skolastisk Teologi. Ligesom for \textsc{Platon} Ideernes
Verden udtrykker Tilværelsens Indhold som evigt, medens de Ting, Iagttagelsen viser
os, er endelige, ufuldkomne og tidsbestemte Former for dette evige Indhold, saaledes
er for \textsc{Spinoza} det evige Forhold mellem Grund og Følge den sande Virkelighed,
medens de i Erfaringen givne Fænomener er tidsbestemte Udformninger (affectiones)
af dette evige Indhold. Forskellen fra \textsc{Platon} bestaar, som allerede ovenfor fremhævet
(se p. 48 f) deri, at \textsc{Spinoza} (som allerede \textsc{Baco} og \textsc{Bruno}) sætter Begrebet
Lov istedetfor Begrebet Ide.

\textit{β.} Ved Spørgsmaalet om, hvorledes Forholdet er mellem den Del af Menneskets
Natur, der hører Tiden til, og den Del, der hører Evigheden til, støtter \textsc{Spinoza} sig
(24---31) til sin Erkendelsesteori. Jo mere vi forstaar Enkelttingene, des mere forstaar
vi Gud, som er sidste Forudsætning for deres Væsen og deres Bestaaen. Vi
bygger da ikke paa løse og tilfældige Iagttagelser, men paa den lovmæssige Sammenhæng
mellem Enkelttingene, altsaa paa Grundlaget for den videnskabelige Forstaaelse,
i hvilken Erkendelse af det individuelle og Indsigten i de almene Love
gaar i ét. (Se ovenfor p. 51---54). Derved tilfredsstilles ikke blot den menneskelige
Aands inderste Trang, men den Erkendelse, ved hvilken Mennesket bliver delagtig
i Guds Væsen, er kun mulig formedelst det evige i selve Menneskets Væsen. Jo
mere Mennesket gaar op i fuldkommen Erkendelse, des mere kan han vide, at han
er i Gud og forstaas ud af Guds eget Væsen, som jo er det evige Forhold mellem
Grund og Følge (30). Til en Erkendelse af nævnte Art svarer indenfor Udstrækningens
Attribut ikke Legemets tidsbestemte Existens, men Legemets Væsen som evig
Modus, set under Evighedens Synspunkt (29 og 31).

Idealet af Erkendelse falder for \textsc{Spinoza} sammen med Idealet af Existens, altsaa
med Substansen. Men den højeste Erkendelsesart kan være udviklet i højst
forskellige Grader (26), og derfor kan ogsaa Forholdet mellem den evige og den forgængelige
Del af Sjælen være højst forskelligt (38).

Naar Intellektualismen naar sit Højdepunkt ved Tydningen af den højeste
menneskelige Erkendelse som Guds Iboen i Mennesket, saa medvirker dog herved en
mystisk Tendens, der aldrig forsvandt hos \textsc{Spinoza}, selv efterat han havde brudt
% --- p. 91 ---
\setcounter{footnote}{0}%
\opage{91}med sin teologiske Fortid, men træder stærkere frem i femte Bog end i de tidligere
Bøger. Det er en umiddelbar Enhedsfølelse, der trænger sig frem. Hans religiøse
Trang vil ikke anerkende nogen Forskel mellem Tanken selv og Tankens højeste
afsluttende Emne. De skal ganske gaa i ét. Menneskets Tænken af Gud er Guds
Tænken af Mennesket. --- Sammenligner vi femte Bog med den korte Traktat, har
Mystiken faaet et mere intellektualistisk Præg i Løbet af \textsc{Spinoza}'s Udvikling. Men
hans Intellektualisme viser sig ikke i passive Kontemplationer, men i Tankevirksomhed.
Det er ved at blive fuldstændig selvvirksomt i sin Erkendelse, ved at leve
og virke ud af sig selv, at Mennesket bliver ét med Substansen, bliver ligesom en
af de Straaler, der fremgaar af den som Lyskilde. Gud er ikke Genstand for
Tanken, saa at denne skulde staa i et Afhængighedsforhold overfor Gud. I selve
Tankevirksomheden rører Gud sig. (Smlgn. ovenfor p. 7---9). \textsc{Spinoza}'s Mystik
er her en Konsekvens af hans oprindelige Forudsætning om Forholdet mellem Erkendelsen
og dens Indhold.

\textsc{Spinoza} undersøger ikke, om der ikke skulde gives andre Former for psykisk
Energi, og dermed for aandelig Kultur, end den rent intellektuelle. Naar \textsc{Goethe}
fandt Grundlaget for sin Udødelighedstro i den Tanke, at vor Aand er af en stedse
virkende Natur, og at naar han indtil sin Død var rastløs virksom, var Naturen
forpligtet til at fremskaffe en ny Tilværelsesform for ham, naar den nærværende ikke
længere kunde holde, --- saa tænkte han ikke paa blot Tankeenergi, men ogsaa paa
andre Former, i hvilke Menneskeaanden kan virke. Tænkeren og Digteren er dog
enige i, at kun de Mennesker er udødelige, i hvilke en stor Energi har udviklet sig:
„Wir sind nicht auf gleiche Weise unsterblich, und um sich künftig als grosse
Entelechie zu manifestiren, muss man auch eine sein“. (Gespräche mit \textsc{Eckermann}.
1. Septbr. 1829). \textsc{Spinoza} siger det samme paa sit Sprog, naar Udødelighed betinges
ved Opnaaelsen af den højeste Erkendelsesform og dennes Herredømme i Sjælen,
og naar han i disse Henseender anerkender forskellige Grader (38---39). Ud fra den
højeste Viden som Udtryk for Evighed kunde han ikke komme til Udødelighed for
alle\footnote[1]{Smlgn. \textit{Den nyere Filosofis Historie\textsuperscript{2}}, I, p. 552 (Note 67), hvor der er nævnt en Række Forfattere
fra Oldtiden og den nyere Tid, hos hvem den Tanke kommer frem, at kun de er udødelige, som
har naaet den højeste aandelige Udvikling. Til de der nævnte Navne maa føjes \textsc{Renouvier}, og maaske \textsc{Sibbern}. --- Ogsaa hos Teologer forekommer denne Antagelse. Saaledes fortolker \textsc{Otto Møller} „den evige Fordømmelse“ som Tilintetgørelse (Brevveksling med \textsc{Rørdam}, II, p. 321).}.

\textsc{Spinoza} mener alligevel, at alle Mennesker er sig Udødelighed bevidst, kun at
de forveksle den med Varighed eller Tidslængde og tro, at der efter Døden kan være
Tale om Erindring og om billedlige Forestillinger (34 Schol.). Strengt taget kan
man hverken tale om en Begyndelse eller om et Ophør eller om en Fortsættelse af
Sjælens sande Væren (31 Schol. og 34 Schol.). Den Udødelighed, \textsc{Spinoza} taler om,
ligger ikke i Fortsættelsen af Tidsexistensen, men er i hvert Øjeblik og har altid
været. Den kommer til vor Bevidsthed, naar vor Tankevirksomhed rører sig paa
afgjort aktiv Maade. „Vi mærker og erfarer“, siger \textsc{Spinoza}, „at vi er evige, ti vi
% --- p. 92 ---
\setcounter{footnote}{0}%
\opage{92}mærker ikke mindre det, vi har en forstaaende Opfattelse af (res, qvas mens intelligendo
concipit), end det, vi erindrer. Sjælens Øjne er selve de tankemæssige Begrundelser“.
Den udvortes Forskel mellem Tanke og Genstand falder bort i den
sande Erkendelse. Evigheden rører sig i selve Tanken, ikke blot i dens Resultater.

Der er to Spørgsmaal, som trænger sig frem overfor hele denne Tankegang
hos \textsc{Spinoza}. Det ene er, om et Menneskes Væsen, naar det helt skal forstaas indenfor
den Sammenhæng, Erkendelsen paaviser, ikke ogsaa maa omfatte dets Forhold
til (Betingethed ved og Indgriben i) de bestemte Tids- og Stedsomstændigheder,
under hvilke det existerer. Dette Forhold præger og betinger dets inderste Natur
og kan ikke være ligegyldigt for den fulde Forstaaelse --- ogsaa fra Evighedens
Standpunkt, forsaavidt det er os muligt at indtage dette, hvilket jo ifølge \textsc{Spinoza}
selv meget sjælden er muligt. Dette første Spørgsmaal angaar Forstaaelsen af Tidsexistensen.
Det andet Spørgsmaal angaar Værdien af den. Der kræves jo ifølge
\textsc{Spinoza} en lang teoretisk og praktisk Udviklingsgang (se ovenfor p. 89---90) for at
naa det Stade, hvor Evighedsexistensen først bliver selvfølgelig. Men saa bliver jo
Evigheden afhængig af Tiden, ikke blot omvendt. \textsc{Spinoza} selv vilde vistnok meget
nødig have indrømmet, at hele det aandelige Arbejde, han saa indgaaende har beskrevet,
skulde være en aldeles ligegyldig Episode. Naar Tidsformens reale Gyldighed
for vor Erkendelse ikke anerkendes, kan det aandelige Arbejdes Realitet
heller ikke fastholdes. Ja, selve Begrebet Værdi forudsætter Tidens Realitet, idet
det forudsætter Trang og Stræben.

Denne Forudsætning vilde \textsc{Spinoza} nu ikke anerkende. Han forlanger jo, at
Tingenes Fuldkommenhed skal vurderes efter deres „Natur og Magt“, og ikke efter
vor Trang. Men Værdibegrebet er derfor ikke forsvundet fra \textsc{Spinoza}'s Filosofi; det
har, som ovenfor (p. 21 f. 33) vist, kun lagt sig til Ro i Grundtanken om Tilværelsens
inderste Lov. Og det er i Forhold til den kosmologiske Værdi, at al anden Værdi
tilsidst bliver forsvindende.

\textit{γ.} I denne Værdi bliver Mennesket delagtig gennem Erkendelsesglæden. Da
Gud som sidste Grund for Erkendelsens Mulighed er Aarsag til Erkendelsesglæden,
opstaar der en intellektuel Kærlighed til Gud (amor intellectualis dei). Den er ikke
rettet mod nogen Genstand i Tiden, men mod et evigt Væsen, som kun Tanken
fatter (32 Cor.). Der kan egentlig ikke tales om nogen Genstand for denne Kærlighed.
Ti ligesom Gud er og virker i selve den menneskelige Erkendelse paa dennes
højeste Trin, saaledes er den intellektuelle Kærlighed til Gud en Del af Guds Kærlighed
til sig selv som den, der virker i Menneskets Væsen, naar dette har udfoldet
sig til evig Tankemodus. Den intellektuelle Følelse er derfor en i højeste Grad
aktiv Følelse (36). I Forhold til den blegner enhver anden Følelse.

Af \textsc{Spinoza}'s egen Psykologi følger en Vanskelighed med Hensyn til Begrebet
„intellektuel Kærlighed“, naar denne skal være en Følelse, der opstaar og bestaar
efter psykologiske Love. Den skal være evig og uforanderlig som Substansen, i
hvis Væsen Mennesket bliver delagtig ved den højeste Erkendelsesgrad. Men ifølge
\textsc{Spinoza}'s Psykologi (se ovenfor p. 63 f) forudsætter enhver Følelse en Tilstands
% --- p. 93 ---
\setcounter{footnote}{0}%
\opage{93}ændring og bestaar i en Overgang --- noget, han selv i denne Sammenhæng minder
om (33 Schol.), i samme Øjeblik som han erklærer, at den intellektuelle Kærlighed
ikke er Glæde (lætitia), der bestaar i Overgangen til Fuldkommenhed, men Salighed
(beatitudo), der bestaar i Besiddelsen af Fuldkommenhed.

Ogsaa paa et andet Punkt nødes \textsc{Spinoza} paa Højdepunktet af sin Mystik til
at ændre sine tidligere opstillede psykologiske Sætninger. I Axiomet i Begyndelsen
af fjerde Bog var det hævdet, at der ikke i Tilværelsen (in rerum natura) gaves
noget Enkeltvæsen, som ikke kunde komme til at staa overfor noget, der var stærkere.
Men nu lærer han (37), at der ikke gives noget i Naturen (in natura), som
er stærkere end den intellektuelle Kærlighed og kunde ophæve den. Han har selv
følt Modsigelsen, men hævder (37 Schol.), at Axiomet i fjerde Bog kun angaar Enkeltvæsener
som existerende paa bestemt Tid og Sted, ikke evige Modi, der er ét
med Substansen. Men --- en psykologisk Udvikling til klar Erkendelse og til Erkendelsesglæde
foregaar dog bestandig under bestemte Sted- og Tidsforhold. Og
\textsc{Spinoza} lærer selv, at den intellektuelle Kærlighed først efterhaanden kan gennemtrænge
hele Sjælen, og at der behøves et stort Arbejde for at naa til, at „den største
Del af Sjælen“ beherskes af den (38---39). Den intellektuelle Kærlighed optræder
altsaa paa dén psykologiske Arena, og hvori ligger Sikkerheden for, at den stedse
vil sejre? Man kan nu engang ikke slippe for Psykologien. Den er allestedsnærværende
paa Aandslivets Omraade, helt op i de højeste intellektuelle, etiske og
religiøse Regioner.

Ogsaa senere Tænkere har søgt at frigøre sig fra Psykologien, hvad disse ideale
Regioner angaar. \textsc{Kant} afviste ethvert Følelsesmotiv paa det etiske Omraade, og
naar han dog maa indrømme, at den moralske Lov kun kan blive en Magt i vort
Indre gennem den Agtelse og Ærefrygt, den indgyder, saa viger han ikke tilbage
for den Konsekvens at erklære, at Agtelse og Ærefrygt slet ikke er Følelser, da de
ikke har noget med Lyst og Ulyst at gøre, ja, overhovedet slet ikke kan forklares
ad Erfaringens Vej\footnote[1]{Se \textit{Den nyere Filosofi's Historie\textsuperscript{2}}, II, p. 83 f. --- \textsc{Kant}'s „Agtelse“ hører til, hvad \textsc{Spinoza} kalder „aktive Følelser“.}. En lignende Forskel som den, \textsc{Spinoza} gør mellem Salighed
og Glæde gør \textsc{Carlyle} mellem Velsignelse (blessedness) og Lykkefølelse (happiness)\footnote[2]{\textit{Sartor Resartus}, Book II, Chap. 9: Man can do without happiness, and instead hereof find blessedness. --- Smlgn. mine Bemærkninger om Terminologien \textit{Etik} III, 11 og VII, 1.}.

Det vilde være urigtigt her at finde en blot Ordstrid, naar fremragende Tænkere
vægrer sig ved at bringe deres højeste Idealer ind under psykologisk Belysning.
Naar Kontinuiteten synes dem brudt, er det, fordi det ikke er lykkedes Psykologien
at paavise lovbestemte og motiverede Overgange fra de laveste til de højeste Former
for Aandsliv. \textsc{Spinoza} har selv (i tredie og fjerde Bog) givet værdifulde Bidrag i
denne Retning, men trods alt, hvad der siden er ydet, er der dog her stedse Opgaver
for beskrivende og analyserende Psykologi. I Distinktioner som dem, \textsc{Spinoza},
\textsc{Kant} og \textsc{Carlyle} hævder, indeholdes der Opgaver for Psykologien, som endnu
kræver at behandles.

% --- p. 94 ---
\setcounter{footnote}{0}%
\opage{94}\subsection*{d. Sætning 41---42.}
\begin{center}\textit{Etiken er uafhængig af Spekulation og Dogme.}\end{center}

\textsc{Spinoza}'s sidste Ord er dog ikke Spekulation og Mystik. Sætning 41, den
næstsidste i hele Værket, lyder saaledes: „Selv om vi ikke vidste, at vor Aand er
evig, vilde vi dog betragte Pligtfølelse, Religion og alt, hvad der i fjerde Bog henregnes
til Livsmod og Højmod som det højeste“. Og han gør i sit Bevis for denne
Sætning opmærksom paa, at han i sin Begrundelse af disse Karakteregenskabers
Værdi ikke har gjort Brug af sin Lære om Udødeligheden. De fleste Mennesker,
tilføjer han, tænker ganske vist anderledes; de vilde ikke agte hine Karakteregenskaber,
hvis der intet var at haabe eller frygte efter Døden. I et af sine Breve
(Ep. 43) hævder han bestemt, at hvad enten de etiske Sætninger (documenta moralia)
er Guds Lovbud eller ikke, er de lige guddommelige og heldbringende. Og
allerede i den korte Traktat (II, 26, 4) spotter han over, at mange, der anses for
store Teologer, siger, at hvis der intet evigt Liv var, vilde de blot søge deres eget
Bedste. Det er, som om en Fisk vilde sige, at naar der efter Livet i Vandet ikke
følger noget andet Liv, vilde den forlade sit eget Element og gaa op paa Land!

Ethica's sidste Sætning (42) udsiger, at Salighed (beatitudo) ikke er Belønning
for værdifulde Karakteregenskaber, men ét med dem. Vi behersker ikke vore
passive Følelser for at opnaa Salighed; men omvendt --- vi kan beherske de passive
Følelser, fordi vi ejer Salighed. Ti det højeste Tankeliv og den Glæde, der er
forbunden med det, formaar at sikre mod at blive passive Følelsers Slave. Det
viser sig atter i denne Slutningsbetragtning, hvor inderlig \textsc{Spinoza} er overbevist om
Tankens Magt til at harmonisere Livet.

% --- p. 95 ---
\setcounter{footnote}{0}%
\opage{95}\section*{Slutningskarakteristik.}


\textit{a.} Den inderste Nerve i \textsc{Spinoza}'s Tænkning er hans Overbevisning om Tilværelsens
Rationalitet. Det Urfænomen\footnote[1]{Om dette fra \textsc{Goethe} stammende Begreb og dets Betydning for ethvert Forsøg paa Verdensanskuelse, se \textit{Den menneskelige Tanke}, p. 311 f.}, paa hvilket hans System er bygget, var
det, at Tingene eller Begivenhederne i Verden kunde opfattes som Led i logiske og
matematiske Rækker, i hvilke ethvert Led forholder sig til de andre som Grund til
Følge eller som Følge til Grund. Ad to Veje har, ifølge den i nærværende Afhandling
opstillede Formodning, Overbevisningen om dette Urfænomens Realitet fæstnet
sig hos ham. I det Gudsbegreb, til hvilket han fra sin Barndom af blev ledet,
fandt han, naar det gennemtænktes med Konsekvens, Principet om en rationel
Grund til alt udtrykt. Og det samme Princip fandt han som Forudsætning for den
nye Naturvidenskabs Stræben efter matematisk Formulering af fysiske Love. Begrebet
Gud og Begrebet Natur mødtes for ham i Begrebet Substans, defineret som
det, hvis Begreb ikke igen forudsætter noget andet Begreb. Rækken af Grunde og
Følger (Aarsager og Virkninger) maa, dersom Tilværelsen er absolut rationel,
hvile paa et Princip, som igen ikke forudsætter noget, men „har sin Grund i sig
selv“ eller er „Aarsag til sig selv“.

Det var desuden hans Overbevisning, at alt Aandsliv ikke blot forudsætter
Tanke, men paa sit Højdepunkt er et Tankeliv. Alle aandelige Fornødenheder tilfredsstilles
da, naar fuldkommen Erkendelse er naaet. Al Kultur er for ham tilsidst
intellektuel. Al Værdi er tilsidst indeholdt i Erkendelsesglæden.

Al Videnskab og al Filosofi arbejder og maa arbejde i den Retning, \textsc{Spinoza}
anviser. Det gælder stedse at finde den indre Sammenhæng mellem de forskellige
Erfaringsfænomener (Modi) og mellem Tilværelsens forskellige Sider eller Egenskaber
(Attributer). Forsaavidt havde \textsc{Hegel} Ret, naar han sagde: „Entweder hast du den
Spinozismus, oder du hast keine Philosophie“.

\textit{b.} Der udkom herhjemme for omtrent 60 Aar siden (1857) en grundig og
skarpsindig Monografi om \textsc{Spinoza} af \textsc{Hans} \textsc{Bröchner}. I mangt og meget stemmer
min Opfattelse af den store Tænker med \textsc{Bröchner}'s. Det var da ogsaa \textsc{Bröchner},
der først ret lærte mig at værdsætte \textsc{Spinoza} og hans Tænkning, selv om der gik
% --- p. 96 ---
\setcounter{footnote}{0}%
\opage{96}mange Aar, inden jeg kom til den Opfattelse, jeg har gjort gældende i nærværende
Skrift. Og Mindet om min Lærer og Ven, der i Karakter og Livsførelse havde saa
meget fælles med \textsc{Spinoza}, har da ogsaa fulgt mig under min fornyede Fordybelse
i Emnet.

Da \textsc{Bröchner}'s Skrift udkom, havde v. \textsc{Vloten} endnu ikke gjort det betydningsfulde
Fund, der bragte „den korte Traktat“ og den oprindelige Text af Brevene for
Dagen og derved ikke blot berigtigede forskellige Personalia, men ogsaa kastede Lys
over \textsc{Spinoza}'s Udviklingsgang. Ikke blot ved, at jeg, ligesom de mange Forfattere,
der i de sidste 50 Aar har beskæftiget sig med \textsc{Spinoza}, har kunnet benytte disse
Dokumenter, men ogsaa paa adskillige enkelte Punkter afviger min Opfattelse fra
\textsc{Bröchner}'s, og iøvrig fra de Spinozaforskere, jeg kender. Jeg har særlig søgt at
opfatte hans Ideer i Sammenhæng med Grundtanker, der forelaa i Datidens Naturvidenskab,
og er derved kommen til at betone den realistiske Karakter af hans Filosofi,
som ganske blev miskendt i den romantiske Periode, saa meget man end sværmede
for \textsc{Spinoza}. Det er denne realistiske Karakter af hans Tænken, der har vakt
saa megen Interesse for ham i den sidste Del af det nittende Aarhundrede. \textsc{Spinoza}
havde den stille Overbevisning, at for at fastholde Aandslivets Realitet i Verden er
det ingenlunde nødvendigt paa Forhaand at sætte Grænser for Erfaringsvidenskabens
Arbejde, hverken paa det naturvidenskabelige eller paa det historiske Omraade.
\textsc{Bröchner} opfatter \textsc{Spinoza} „som Repræsentant for den nyere Idealismes Begyndelsespunkt“,
som Indleder af en Retning, der skulde faa sin Afslutning ved \textsc{Hegel}. Efter
min Opfattelse blev den romantiske Idealisme netop kun mulig ved et Brud med
den Retning, \textsc{Spinoza} var slaaet ind paa. \textsc{Spinoza} er ikke Idealist, hvis dertil hører
principiel Fornægtelse af Erfaringsvidenskabens Synsmaader og Metoder.

Men den væsentligste Forskel mellem min og \textsc{Bröchner}'s Opfattelse af \textsc{Spinoza}
vil fremgaa af nogle Bemærkninger, med hvilke jeg vil slutte denne Afhandling.

\textit{c.} Kan vi virkelig, som \textsc{Spinoza} forudsætter, afslutte vor Videnskab med
Tanken om noget, der forstaas ved sig selv og bestaar ved sig selv? --- Tænkning
bestaar i Sammenfatten og Sammenligning, hin som Forudsætning for denne. Men
enhver Sammenfatten medfører en Begrænsning, der stedse kan ophæves ved ny
Erfaring, og enhver Sammenligning forudsætter et Synspunkt, der selv kan blive
Genstand for Sammenligning med andre Synspunkter. Vi kan ende med et Spørgsmaal,
men ikke med en en Gang for alle afsluttende Tanke. Hvad der for \textsc{Spinoza} % PRINTED AS IS: en doubled
var Urfænomenet, det, der havde sin Grund i sig selv, var Lovmæssigheden i Verden.
Men denne Lovmæssighed er en stor Hypotese, som Videnskaben under sit stadige
Arbejde søger at finde Bekræftelse paa gennem stedse ny Anvendelse. Og (som
\textsc{Spinoza} selv har indrømmet) de enkelte Fænomener og Begivenheder kan ikke udledes
af selve Lovmæssighedens Begreb, selv om de kan forstaas (og kun kan forstaas)
ved Anvendelse af dette Begreb.

I sin førkritiske Periode var \textsc{Kant} inde paa en Tankegang, der minder om
\textsc{Spinoza}'s. Han sluttede fra den lovmæssige Sammenhæng i Naturen til et absolut
% --- p. 97 ---
\setcounter{footnote}{0}%
\opage{97}Væsen, der laa til Grund for alt. Det var i denne Periode for ham „den eneste
mulige Bevisgrund“ for Gudsbegrebets objektive Gyldighed. Men Overgangen til
den kritiske Filosofi blev gjort, da \textsc{Kant} ikke længere spurgte om Verdens objektive
Grund, men om Muligheden af et objektivt Verdensbegreb, og da han fandt denne
Mulighed i selve Grundformen for menneskelig Tænken, hvis Gyldighed beror paa,
at uden streng logisk Forbindelse af Iagttagelser er Erfaring i streng Forstand ikke
mulig\footnote[1]{Smlgn. Om \textit{Kontinuiteten i Kant's filosofiske Udviklingsgang} (Vid. Selsk. Skr. 1893), p. 14---15.}. Den aktive Karakter af Tænkningen, som \textsc{Spinoza} saa stærkt havde fremhævet,
blev nu lagt til Grund for hele Betragtningen. Syntesebegrebet afløste Substansbegrebet,
og medens \textsc{Spinoza}'s Filosofi først og fremmest var, eller vilde være,
Kosmologi (Metafysik), en Lære om Tilværelsens inderste Natur, trængte nu Erkendelsesproblemet
sig frem foran alle andre filosofiske Problemer, og en videnskabelig
Kosmologi kom til at staa som et Grænseproblem. Til \textsc{Spinoza}'s Substans („det,
som er i sig selv“) svarer i \textsc{Kant}'s Filosofi „Tingen i sig selv“, om hvilken intet
kan vides (ikke engang, at den „forstaas ved sig selv“).

En nærmere Drøftelse af Erkendelsesproblemet førte igen til at sondre mellem
forskellige Grupper af Kategorier, svarende til \textsc{Spinoza}'s Skelnen mellem Substans,
Attributer og uendelige Modi af forskellig Grad, men uden den umiddelbare metafysiske
Belydning, han tillagde dem. Der skelnes mellem Grund og Aarsag, altsaa % PRINTED AS IS: Belydning (Betydning)
mellem formale og reale Kategorier, og Værdibegreberne fremtræder i deres Forskel
fra begge disse Grupper. Deraf følger igen, at al aandelig Kultur ikke, som \textsc{Spinoza}
i Grunden forudsætter, kan være ét med intellektuel Kultur.

Men den ideale Opfattelse af Tankearbejdet, som \textsc{Spinoza} har hævdet, og den
inderlige Erkendelsesglæde, som for kam kastede Glans over Livet, kan ikke mangle % PRINTED AS IS: kam (ham)
hos nogen sand Forsker. Enhver Iagttagelse og enhver Tanke, der giver et Indblik
i den store Sammenhæng, som al Forstaaelse af Tilværelsen forudsætter, lader os
føle Betydningen af \textsc{Goethe}'s, denne trofaste Spinozadyrkers Ord: „Ort für Ort
sind wir im Inneren“, skønt vi ikke, som \textsc{Spinoza} mente, en Gang for alle kan afsløre dette Indre.

% --- p. [99], INDHOLD (PDF 487), read at the image 2026-10-04.  p. [98] is blank.
\clearpage
\begin{center}{\large INDHOLD}\end{center}
\medskip
\hfill{\scriptsize Side}\par
\tocline{Indledning}{1---13}
\tocline{Første Bog}{14---35}
\tocline{Anden Bog}{36---57}
\tocline{Tredie Bog}{58---69}
\tocline{Fjerde Bog}{70---83}
\tocline{Femte Bog}{84---94}
\tocline{Slutningskarakteristik}{95---97}

\end{document}
